\begin{savequote}
\includegraphics[width=4cm]{./images/gabriel-christine-aass/2010-06-14_010.jpg}
\end{savequote}


\chapter
[\CO{[EHI]} Erste Hilfe und erweiterte Erste Hilfe]
{\CC{[EHI]}\\Erste Hilfe und erweiterte Erste Hilfe}
\label{chp:EHI}


\ChapterIntro
{%
Unter Erster Hilfe versteht man Maßnahmen, um das Leben und
die Gesundheit von Patienten zu retten und bis zum
Eintreffen professioneller Hilfe zu erhalten, sowie um
bedrohende Gefahren abzuwenden. Erste-Hilfe-Maßnahmen sind
in der Regel simpel und mit wenig Aufwand durchzuführen. Sie
sind ein essentieller Bestandteil der Rettungskette und
haben, wenn sie frühzeitig und richtig durchgeführt werden,
einen enormen Einfluß auf das spätere Behandlungsergebnis. 
}
{%
\begin{description}
  \item[Maintainer:] Roman Koch
  \item[Co-Maintainer:] Christof Koller
  \item[Reviewer:] Standard-Reviewprozess
  \item[Version:] {\DocVersion} ({\DocVersionInfo})
  \item[SHA1:] {\DocGitCommit}
\end{description}

{\TxTeilkapitelAASS}
}

\section{Einleitung}
\label{topic:behandlungspflicht-ehi}


\subsection{Allgemeine Einleitung}


Erste Hilfe ist die Rettung von Menschen aus lebensbedrohlichen
Situationen unter Einsatz von einfachen Mitteln und Techniken. Sie
kann von jedem schnell und effektiv
durchgeführt werden und ist -- trotz moderner Medizin -- noch
immer Beginn und Grundpfeiler jeder Behandlung. Daher ist
Erste Hilfe auch für den Rettungssanitäter die Basis jeder 
Patientenversorgung und folglich Ausgangspunkt seiner
Ausbildung.

% {
% Die Erste Hilfe umfasst die Behandlung von (uU. lebensbedrohlichen)
% Notfällen und Unfällen. Unfälle sind alle durch äußere
% Gewalt verursachten Verletzungen der körperlichen Integrität.
% Notfälle hingegen umfassen Erkrankung und Vergiftungen. Beide
% Zustände können in ihrer Folge zur Beeinträchtigung der Psyche
% führen, wodurch der Ersthelfer nicht nur Maßnahmen zur
% Aufrechterhaltung der physischen, sondern auch der psychischen
% Unversehrtheit treffen muss.}

\Layout{0}{Grundsätze}{
Bei jeder Behandlung hat sich der Ersthelfer an folgende
Grundsätze zu halten:

\subparagraph{1. Grundsatz:} Jeder Mensch ist zur
Durchführung von Erster Hilfe verpflichtet! Sie darf nur bei
Unzumutbarkeit abgelehnt werden!

\Z{Unzumutbar ist die Erste Hilfe für den Rettungssanitäter
nur, wenn er durch Setzen   der Maßnahme sein eigenes Leben
und/oder seine Gesundheit gefährden würde.}

\subparagraph{2. Grundsatz:} Die zumutbare
Erste-Hilfe-Maßnahme ist zumindest zu versuchen und darf
nicht von vornherein als unmöglich abgelehnt werden!

\Z{
  Gemeint sind hierbei Fälle in denen für den Ersthelfer
  zwar keine
Gefahr droht, er sich aber an seine physischen Grenzen
gesetzt fühlt. Ein denkbares Beispiel wäre die   Rettung
eines stark übergewichtigen Menschen aus einem
Gefahrenbereich durch   eine körperlich schwache Person.}

\subparagraph{3. Grundsatz:} Die Gesamtbehandlung des
Patienten unterliegt einer Interessensabwägung!

\Z{
Es muss immer das nächstgrößere Übel beseitigt werden. Es kann mitunter auf
Grund der Dringlichkeit in der Situation nicht immer Rücksicht auf das Eigentum
des Patienten genommen werden.
%   Da die Erste Hilfe nicht nur dem Gebot der Schnelle unterliegt,
% sondern auch mit   einfachsten Mitteln auskommen muss, ist es teilweise
% notwendig eine geringfügige   Beeinträchtigung des Patienten oder
% seiner Besitztümer für den Schutz von höherwertigen Interessen
% in Kauf zu nehmen. Als Maßstab fungiert die Beachtung einer  
% vernünftigen Mittel-Ziel-Relation. 
Als Beispiel dient das Aufschneiden der Kleidung des Patienten zur besseren Untersuchung und
Behandlung.
}

\subparagraph{4. Grundsatz:} Bei jeder Einzelbehandlung ist der Patient so
schnell und so schonend wie möglich zu behandeln!

\Z{
  Dieser Grundsatz erfasst alle, am Patienten durchführbaren
Maßnahmen. 
% Neben diesen allgemeinen Grundsätzen herrschen bei der Rettung des
% Patienten aus einem Gefahrenbereich noch weitere Grundsätze die an
% passender Stelle näher behandelt werden.
}
}{
\begin{enumerate}
  \item Verpflichtung zur Hilfeleistung
  \item Hilfeleistung muss zumindest versucht werden
  \item Interessensabwägung
  \item Hilfeleistung so schnell und schonend wie möglich
\end{enumerate}
}{}{}{}

% In ihrer Durchführung folgt die Erste Hilfe stets einem gewissen
% Ablauf, der ähnlich einem Baum mit einer Anfangsmaßnahme als
% Stamm beginnt und sich durch Ansammeln weiterer Informationen über
% den Zustand des Patienten in verschiedene Behandlungsmaßnahmen als
% \ Äste verzweigt. In der folgenden Ausarbeitung wird versucht die
% Erste Hilfe anhand dieses Ablaufschemas praxisnah wiederzuspiegeln und
% den Ersthelfer vom Eintreffen am Not- oder Unfallort bis hin zur
% Übergabe an weiterbehandelnde Gesundheitsberufe zu begleiten.
% 
% Als der Verständlichkeit dienenden Überblick ist das
% Inhaltsverzeichnis zu empfehlen, da es auf einen Blick den
% Behandlungsalgorithmus widerspiegelt.



\subsection{Rettungskette}
\index{Rettungskette}
\label{topic:rettungskette}


Jede Bewältigung eines Notfalls sollte immer nach dem
selben Schema ablaufen, der Rettungskette.
% Als eine andere Form des überblickhaften Aufbaus der Gesamtbehandlung
% eines Patienten wird seit langem die Rettungskette verwendet. Sie baut
% sich aus nachstehenden fünf Gliedern auf. 

% TODO EHI: Bild Rettungskette einfügen


\Layout{0}{Rettungskette}
{
%%%%%%%%%%%%%%%%%%%%%%%%%%%%%%%%%%%%%%%%%%%%%%%%%%%%%%%%%%%
% TODO Rettungskette: Grafik Neuerstellung
\TODO{GABS}{yyyy-mm-dd}{Rettungskette: Grafik Neuerstellung}{}
%%%%%%%%%%%%%%%%%%%%%%%%%%%%%%%%%%%%%%%%%%%%%%%%%%%%%%%%%%%
Die Rettungskette besteht aus fünf Gliedern. Am Beginn
stehen die \EE{Lebensrettenden Sofortmaßnahmen}, welche noch
in \E{Allgemeine Sofortmaßnahmen} und \E{am Patienten
gesetzten Lebensrettenden Sofortmaßnahmen} unterteilt werden
können. Zu den \E{Allgemeinen Sofortmaßnahmen} zählt das
Absichern der Unfallstelle (Eigensicherung!), das Retten von
Personen aus Gefahrenzonen, sowie das Absetzen des Notrufs
durch eine 2. Person. Zu den am Patienten zu setzenden
\E{Lebensrettenden Sofortmaßnahmen} zählt die stabile
Seitenlage bei Bewusstlosigkeit, die
Herz-Lungen-Wiederbelebung bei Atem-Kreislaufstillstand,
Blutstillung bei starken Blutungen, sowie Schockbekämpfung
beim Schock. Nach den Lebensrettenden Sofortmaßnahmen wird
der \EE{Notruf} abgesetzt, falls dies noch nicht durch eine
2. Person geschehen ist.

Anschließend wird \EE{weitere Erste Hilfe} geleistet, wie
zum Beispiel die Versorgung von einfacheren Verletzungen und
die psychologische Betreuung des Patienten. Der eintreffende
Rettungsdienst übernimmt den Patienten für die
\EE{Sanitätshilfe} und führt den \EE{Transport} durch. Im
Spital erfolgt schließlich die \EE{ärztliche Betreuung} des
Patienten.

% \Hand{Je nachdem wie
% gut die Rettungskette funktionert, landet der Patient im
% Erdgeschoss oder im Keller des Spitals.}
}{
\begin{itemize}
  \item Lebensrettende Sofortmaßnahmen
  \item Notruf
  \item Weitere Erste Hilfe
  \item Sanitätshilfe - Transport
  \item Ärztliche Betreuung
\end{itemize}
}
{}{}{}





% {
% Von der Kette gehören die ersten drei Glieder zu den ersthelferischen
% Maßnahmen. Das erste Glied mit dem Thema lebensrettende
% Sofortmaßnahmen besteht nach herrschender Meinung aus folgenden
% Maßnahmen:}
% 
% \begin{enumerate}
%     
%     
%     
%     \item i) Gefahrenbekämpfung:
%     
%     {\scriptsize
%     Um Gefahren wirksam bekämpfen zu können, muss der Ersthelfer sie
%     zuerst erkennen,  daraufhin den Gefahrenbereich absichern und den
%     Patienten aus dem Gefahrenbereich retten.}
%     
%     \item ii) Kontrolle der Vitalfunktionen:
%     
%     {\scriptsize
%     Nach erfolgreicher Gefahrenbekämpfung ist der Patient auf seine
%     Lebensfunktionen hin zu  untersuchen. Dabei wird mit der
%     Bewusstseinskontrolle begonnen und der Atemkontrolle  aufgehört.}
%     
%     \item  iii) stabile Seitenlagerung:
%     
%     {\scriptsize
%     Sollte der Patient bewusstlos, aber bei Atmung und Kreislauf sein, muss
%     er in die stabile  Seitenlagerung umgelegt werden.}
%     
%     \item  iv) Wiederbelebung:
%     
%     {\scriptsize
%     Konnten weder Bewusstsein noch Atmung und Kreislauf beim Patienten
%     festgestellt werden,  ist mit einer Wiederbelebung zu beginnen.}
%     
%     
%     \item  v) Schockbekämpfung:
%     
%     {
%     {\scriptsize{ Befindet sich der Patient in einem Zustand
%     lebensbedrohlicher Kreislaufstörungen (Schock),}}\footnote{Für
%     einen Ersthelfer sollte an dieser Stelle zur besseren Verständnis
%     zwischen dem \E{Schock} und dem \E{Schreck} unterschieden werden. Der
%     Begriff Schock bezeichnet vorrangig, und damit abweichend als im Volksmund
%     geläufig, eine Störung der Kreislauffunktionen, also der Physis.
%     Der Begriff Schreck hingegen umschreibt eine Störung der geistigen
%     Unversehrtheit, also der Psyche. Der somit im Volksmund in Situationen
%     der Angst und Unruhe verwendete Begriff Schock, ist in Wahrheit der
%     Schreck. }{\textit{  ist sofort mit deren Bekämpfung zu
%     beginnen. Eine Bekämpfungsart ist die Stillung von stark  blutenden
%     Wunden. Die Blutstillung ist streng zu trennen von der einfachen
%     Wundversorgung,  welche zum Glied der weiteren Erste Hilfe gehört.}}}
% \end{enumerate}
% 
% {
% Nähere Beschreibungen der einzelnen Maßnahmen siehe unten.}
% 
% {
% Das zweite Glied erfasst den Notruf und beschäftigt sich mit den
% Fragen: {\quotedblbase}Wo rufen ich an?{\textquotedblleft} und
% {\quotedblbase}Was muss ich beim Notruf sagen?{\textquotedblleft}.}
% 
% {
% Das letzte ersthelferische Glied ist die Weitere Erste Hilfe. Diese
% besteht aus folgenden Maßnahmen:}
% 
% \begin{enumerate}
%     
% \item  i) Wundversorgung:
% 
% {\scriptsize
%  Sie beschäftigt sich mit der Versorgung von nicht stark blutenden
% Wunden und ist somit von  der Blutstillung als lebensrettende
% Sofortmaßnahme zu unterscheiden.}
% 
% \item  ii) Erkrankungsbekämpfung:
% 
% {
% {\scriptsize{ Sollte der Patient erkrankt sein, muss der Ersthelfer
% versuchen die Erkrankung durch  vorrangige Linderung der Symptome zu
% bekämpfen.}}{ }}
% 
% \item  iii) Vergiftungsbekämpfung:
% 
% {\scriptsize{ Ein an Vergiftungssymptomen leidender Patient ist so zu
% behandeln, dass versucht wird seine  Symptome zu lindern und das Gift
% }}{\textit{vom}}{\textit{ Körper, nicht aber
% }}{\textit{aus}}{\textit{ dem Körper, zu beseitigt.}}
% \end{enumerate}


Der Ersthelfer steht am Anfang der Rettungskette. Von seinem Handeln hängt 
zunächst alles ab, im Extremfall auch das Leben des Patienten.

% Die Rettungskette
% erlaubt es auf einfache Weise wichtige Faktoren der gesamten Patientenbehandlung darzustellen. Dazu zählt die Tatsache,
% dass von den fünf Gliedern der Kette ganze drei Glieder als deutliche
% Mehrheit dem Ersthelfer obliegen. Untermauert wird dies durch den auch
% hier gültigen Satz: 

\Fazit{"`Eine Kette ist immer nur so stark 
wie ihr schwächstes Glied."'}

Dieser Satz soll  die Wichtigkeit der Erste-Hilfe-Maßnahmen
verdeutlichen. Schließlich werden von den fünf Gliedern der
Rettungskette die ersten drei vom Ersthelfer durchgeführt.

%  Dies verdeutlicht
% die Wichtigkeit ersthelferischer Maßnahmen. Des Weiteren zeigt sie,
% dass bei jeder Behandlung die lebensrettenden Sofortmaßnahmen der
% Notrufabgabe, aber vor allem der weiteren Ersten Hilfe vorgehen.
% \footnote{Dieser Umstand verdeutlicht den oben definierten 3. Grundsatz einer
% Interessenabwägung.
% 
% 
% Leider ist die Rettungskette nicht frei von Makel, da sie für viele
% Situationen die Erste Hilfe zu generalisiert behandelt und es in dem
% tatsächlichen Ablauf der Behandlung zu Differenzen kommt. Dies
% betrifft einerseits die Beziehung zwischen dem ersten Glied der
% lebensrettenden Sofortmaßnahmen und dem Notruf als unabhängiges
% zweites Glied. In der realen Behandlung gibt es nur wenige Situationen
% in denen eine Sofortmaßnahme bis zum Schluss durchgeführt wird,
% um erst nach deren Beendigung die Notrufabgabe einzuleiten.
% Tatsächlich muss der Notruf schon während der Sofortmaßnahme
% abgegeben werden und ist daher auch eher als ein Bestandteil der
% lebensrettenden Sofortmaßnahme zusehen als ein eigenständiges
% Glied. Andererseits ist die grobe Trennung zwischen lebensrettenden
% Sofortmaßnahmen und Weiterer Ersten Hilfe nicht zielgerecht, da
% viele unter die Weitere Erste Hilfe fallende Situationen sich als
% lebensbedrohliche Gefahren offenbaren können und es daher zu ihrer
% Bekämpfung stets lebenserhaltende Sofortmaßnahme benötigt. }
% 
% \Z{
% Die nachstehende Ausarbeitung folgt zwar in ihrem Rahmen der
% Unterteilung in lebensrettende Sofortmaßnahmen und Weitere Erste
% Hilfe, setzte sich dafür aber bei den Themen der Weiteren Ersten
% Hilfe näher mit den uU. notwendigen lebensrettenden
% Sofortmaßnahmen auseinander.}



\section{Lebensrettende Sofortmaßnahmen}

\subsection{Verhalten am Notfallort}
\label{topic:verhalten-notfallort}


% Vor dem Setzen irgend einer Maßnahme ist es für den Erfolg der
% Behandlung (und uU. für die Gesundheit des Ersthelfers) wichtig, dass
% sich der Ersthelfer:

\Layout{0}{Verhalten am Notfallort
}{
Vor dem Setzen von Maßnahmen sollte der Ersthelfer tief
durchatmen und sich danach einen \EE{Überblick} über die
Situation verschaffen: Was ist passiert? Um welche Art von
Notfall handelt es sich? Welche Gefahren können drohen?

% \dictum[Dickie, in: Samuel Shem: House of God.
% \cite{HouseOfGodDe}]{Alles in Ordnung, auch wenn Sie in Panik geraten sind und sich beschissen fühlen.
% Ich kenne das. Es ist furchtbar, und es wird nicht das letzte Mal sein.
% Aber vergessen Sie nicht, was Sie gesehen haben. Regel Nr. 3: Bei Herzstillstand
% zuerst den eigenen Puls fühlen.}

Es gilt \EE{Ruhe zu bewahren} und \EE{Ruhe zu
vermitteln}. Dabei stellt man sich ruhig und gezielt auf die
bevorstehenden Maßnahmen ein und reiht diese nach
\EE{Dringlichkeit}. Erst nach diesen anfänglichen
Überlegungen wird mit einer Behandlung begonnen! 
}{
\begin{itemize}
  \item Überblick verschaffen
  \item Ruhe bewahren, Ruhe vermitteln
  \item Maßnahmen nach Dringlichkeit sortieren
\end{itemize}

}{}{}{}

\subsection{Gefahrenbereich}
\index{Gefahrenbereich}
\label{topic:gefahrenzone}

\Layout{0}{Gefahr
}{
Durch das Verschaffen des Überblicks kann der Ersthelfer
\EE{Gefahren erkennen}. Diese können für den Patienten, den
Ersthelfer und dazukommende Dritte bestehen. Der Ersthelfer
muss \EE{Absperrmaßnahmen durchführen} und damit den
Notfallort absichern.

Muss der Patient aus dem Gefahrenbereich gerettet werden,
ist unbedingt auf Selbstschutz zu achten. Anschließend
müssen ggf. \EE{Spezialkräfte} angefordert und ihnen
sämtliche verfügbaren Informationen über die Gefahr bekannt
gegeben werden.
}{
\begin{description}
  \item [G]efahr erkennen
  \item [A]bsperrungsmaßnahmen durchführen - Absichern
  \item [S]pezialkräfte anfordern
\end{description}
}{}{}{}

\subsubsection{Absichern des Gefahrenbereichs}

% Sobald die Gefahr erkannt wurde ist es primär notwendig den
% Gefahrenbereich für den Patienten, für sich selber und für
% dazukommende Dritte abzusichern. 
Die Methoden der Absicherung sind vielfältig und müssen den jeweiligen
Gegebenheiten angepasst werden. Als Grundsatz jeder
Absicherung gilt:

\Fazit{Selbstschutz geht vor Fremdschutz!}

% Dieser Leitsatz bedeutet, dass vor der Sicherung des
% Patientenlebens der Ersthelfer sein eigenes Leben und seine Gesundheit bewahren und sich dementsprechend
% absichern und verhalten muss.
% Diese Grundregel leitet sich aus dem 1. Grundsatz der Ersten Hilfe ab
% und berührt nicht nur den Laien, sondern jede Person im
% Gesundheitsberuf. 
In der Praxis ist es oft möglich, durch
eine einfache Absicherungsmaßnahme das eigene Leben, sowie
das des Patienten, effektiv zu schützen.

\Layout{0}{Beispiel: Vorgehen bei einem Verkehrsunfall}
{
% Für die Arten der Absicherung muss nach folgenden Situationen
% unterschieden werden:
%
% \subparagraph{Absicherung im Gefahrenbereich Verkehr}~
Wenn sich der Ersthelfer mit einem Fahrzeug einer
Unfallstelle nähert, hält er nach Möglichkeit sein Fahrzeug
in einem angemessenen \EE{Sicherheitsabstand} zum
Gefahrenbereich an. Besondere \EE{Gefahrenquellen} wie
Kurven, Gefälle, Hangkuppen, sowie die Windrichtung bei
Rauchentwicklung müssen beim Absichern der Unfallstelle
berücksichtigt werden. Zum Absichern selbst werden
Hilfsmittel, wie \EE{Warnleuchte} oder \EE{Warndreieck},
sowie Personen, welche Handzeichen geben, eingesetzt.
}{
\begin{itemize}
  \item Sicherheitsabstand
  \item Gefahrenquellen beachten (Kurven, Gefälle,
  Windrichtung, Gefahrenstoffe,\ldots)
  \item Hilfsmittel
  \begin{itemize}
    \item Warnleuchte, Warndreieck, Handzeichen
  \end{itemize}
  \item Selbstschutz
  \begin{itemize}
    \item Warnweste
    \item Mit geöffnetem Warndreieck
    Verkehr entgegen gehen
  \end{itemize}
\end{itemize}
}{}{}{}



\Fazit{Niemals selbst gefährden!}

Daher hat der Ersthelfer vor dem Aussteigen aus dem Auto
eine \EE{Warnweste} anzuziehen und mit bereits
geöffnetem \EE{Warndreieck} (Pannendreieck) dem 
Verkehr entgegenzugehen. 


\Layout{0}{Pannendreieck}
{
Das Pannendreieck ist
dem jeweiligen Anhalteweg entsprechend zur Unfallstelle
aufzustellen:
\begin{itemize}
  \item \E{Ortsgebiet}: unmittelbar bis zu ca. 50\Meter*
  \item \E{Freilandstraße}: ca. 150\Meter*
  \item \E{Autobahn}: ca. \VonBis{250}{400}\Meter*
\end{itemize}
}{
\begin{itemize}
  \item \E{Ortsgebiet}: unmittelbar bis zu ca. 50\Meter*
  \item \E{Freilandstraße}: ca. 150\Meter*
  \item \E{Autobahn}: ca. \VonBis{250}{400}\Meter*
\end{itemize}
}{}{}{}

\Layout{0}{Eisenbahnkreuzung}
{\index{Eisenbahnkreuzung}
Bei einem Unfall an einer Eisenbahnkreuzung ist zu bedenken, dass der
herannahende Zug einen \E{viel längeren Bremsweg} als ein
Auto hat. (je nach Geschwindigkeit und
Gewicht des Zuges, ca. \VonBis{600}{800}\Meter*). Das
\EE{Haltesignal} für den Lokführer ist eine \EE{kreisende Armbewegung}. Dadurch wird er gewarnt und leitet 
eine Notbremsung ein.
}{
\begin{itemize}
  \item Viel längerer Bremsweg
  \item Haltesignal: kreisende Armbewegung
\end{itemize}
}{}{}{}



% {
% Für die Absicherung sind mehrere Möglichkeiten denkbar, die sich
% daran unterscheiden, wie (d.\,h. mit welchem Fahrzeug) der Ersthelfer zu
% der Unfallstelle hinkommt. Bei allen Varianten besteht die Pflicht zur
% Einhaltung eines Sicherheitsabstandes von ca. 50 m zum
% Gefahrenbereich.}

% {\scriptsize
% {Eine Möglichkeit ist, dass der Ersthelfer mit einem
% }{Kraftfahrzeug}{ die Unfallstelle antrifft. In diesem
% Fall muss der Ersthelfer zuerst sich selbst absichern, indem er die
% Warnblinkanlage einschaltet und mit angezogener Sicherheitsweste das
% Fahrzeug verlässt. Sicherheitswesten müssen in Österreich
% verpflichtend im Kraftfahrzeug mitgeführt und im Fahrerraum
% griffbereit aufbewahrt werden. Ein Überhängen der Weste über den
% Fahrer- oder Beifahrersitz ist grds. erlaubt, außer das Fahrzeug
% besitzt Seitenairbags in den Sitzen. In einem solchen Fall würden die
% Westen bei einem Seitenaufprall des Fahrzeuges die Wirkung der
% Seitenairbags verringern oder ganz beseitigen. Verpflichtend
% mitzuführen ist nur eine Weste, sinnvoll sind aber so viele wie
% Personen mit Fahrzeug mitfahren können, da jeder Begleiter denselben
% Anspruch auf Sicherheit hat wie der Fahrer.}}

% {
% Nach dem Verlassen des Fahrzeuges ist das ebenso verpflichtend
% mitzuführende Pannendreieck entsprechend des Bremsweges aufzustellen.
% In Zahlen ausgedrückt bedeutet dies:}
% 
% \begin{compactitem}
% \item Autobahn:  250-400\,m
% \item Freilandstraße: 150-200\,m
% \item  Ortsgebiet:  0-50\,m\
% % \item footnote{
% Die im Ortsgebiet anwendbare Distanz von null Metern entspricht
% natürlich nicht dem Bremsweg. Sie ist aber in manchen Situationen
% notwendig, um den tatsächlich gefährdeten bzw. gefährdenten
% Verkehr zu warnen. Ein Beispiel dazu wäre ein Erliegen mit dem
% Fahrzeug 1-2m hinter einer Kreuzung. Würde man in einem solchen Fall
% das Pannendreieck 50m weit entfernt \ aufstellen, wäre das
% Warnzeichen in einem Bereich wo der herannahende Verkehr keinen Bezug
% zu einem Hindernis erkennen kann, wodurch er maximal irritiert aber
% nicht gewarnt ist. Der Ersthelfer muss direkt hinter seinem
% Kraftfahrzeug das Pannendreieck aufstellen, damit der aus der Kreuzung
% herausfahrende Verkehr die Absicherung verstehen und sich daran halten
% kann.}
% \end{compactitem}
 

 





% \Z{
% Für Krafträder gelten andere Bestimmungen.
% Ein Motorrad- oder Mopedfahrer ist gesetzlich nicht zur Mitnahme einer
% Warnweste oder eines Pannendreieck verpflichtet, obwohl es natürlich
% auch für diese Verkehrsteilnehmer passende Sicherungsmittel gibt. Der
% Kraftradfahrer ist nur verpflichtet zur Mitnahme einer Verbandsbox, die
% den Anforderungen des {\S}\,102 KFG
% entspricht.\footnote{Nähere Informationen zu den Arten von
% Verbandsboxen siehe unten} Der Ersthelfer muss sein Kraftrad
% quer zur Fahrtrichtung als sichtbare Absperrung aufstellen und die zum
% herannahenden Verkehr gerichtete Blinkerseite einschalten. An Stelle
% des Pannendreiecks empfiehlt es sich den Helm in denselben Abständen
% aufzustellen, um somit dem herannahenden Verkehr ein Zeichen von
% Abnormität im Verkehr mitzuteilen. Denn selbst für den
% passioniertesten Fahrer ist ein alleine auf der Fahrbahn liegender Helm
% nicht alltäglich.
% Fahrradfahrer, Fußgänger oder sonstige
% Verkehrsteilnehmer verfügen im Normalfall über keine
% vergleichbaren Absicherungsmittel. Eine Möglichkeit ist durch Wink-
% oder Handzeichen aus einem sicheren Bereich heraus dem nachkommenden
% Verkehr das Stehenbleiben anzudeuten, der wiederum mit seinen
% Absicherungsmitteln helfend eingreifen kann. 
% }





% \subparagraph{Absicherung im Gefahrenbereich Haushalt}
% 
% {
% Auch im Haushalt sind Gefahrenbereiche als Folge eines Unfalls denkbar.
% Als Beispiel zu nennen wäre ein noch eingeschaltetes \EE{Elektrowerkzeug}
% durch das sich der Patient verletzt hat. Das Elektrogerät ist vor
% Behandlung des Patienten zu sichern. Da es neben diesem Beispiel noch
% unzählige weitere Gefahrenbereiche geben kann, die auf
% unterschiedliche Weise abzusichern sind, kann hier keine Grundregel zur
% Absicherung aufgestellt werden. Der Ersthelfer muss sich den
% erwähnten Grundsätzen gerecht verhalten!}

% \subsubsection{Tipps und Tricks für den Ersthelfer bei
% Gefahrenbereichen}
% 
% Als Quelle von Gefahren kommen im Alltag unzählige Situationen in
% Betracht, wodurch es an dieser Stelle unmöglich ist alle Gefahren zu
% behandeln. Die folgende Auflistung erfüllt daher keinen Anspruch auf
% Vollständigkeit, sondern versucht nur die relevantesten Gefahren zu
% skizzieren.
% 
% % \subparagraph{
% % Gefahren im Haushalt und am Arbeitsplatz:}
% % 
% % \Z{
% % Die Mehrheit von Gefahren verwirklichen sich in den eigenen vier
% % Wänden oder dem Arbeitsplatz. Hierbei werden gerade alltägliche
% % Situationen zur Quelle von Gefahren, weil der Umgang mit ihnen als
% % harmlos und alltäglich empfunden wird und sich dadurch eine sorglose
% % Handhabe einstellt, die Gefahren erzeugt oder erhöht. Zu den
% % möglichen Gefahrenquellen gehören:
% % 
% % \begin{enumerate}
% %     \item  falsch verwendete Elektrogeräte
% %     \item  unvorsichtiges Ausüben der Wohnungsreinigung
% %     \item  unvorsichtiges Verwenden von Werkzeug
% % \end{enumerate}
% % 
% % 
% % Als am häufigsten betroffene Personengruppe im Haushalt sind Kinder zu
% % nennen, da sie ihre Umgebung aus einem anderen Blickwinkel wahrnehmen
% % und gleichzeitig einen natürlich starken Entdeckerdrang haben.
% % Gefahrenquellen für Kinder sind zum Beispiel:
% % 
% % \begin{enumerate}
% %     \item  giftige Zimmerpflanzen (Bsp: Weihnachtsstern aus nicht
% % kontrollierten Anbau)
% %     \item  ungesicherte Steckdosen
% %     \item  in Getränkeflaschen abgefüllte Haushaltsreiniger
% % \end{enumerate}
% % 
% % }
% % 
% % \begin{itemize}
% %     \item Eine der größten Gefahren für den Ersthelfer auf Strassen oder
% %     Autobahnen ist, das Überfahrenwerden von nachkommenden Fahrzeugen.
% %     \item Beim Übersteigen einer Leitschiene muss vorsichtig geprüft werden, ob
% %     auf der anderen Seite auch ein fester Untergrund vorhanden ist. Sonst liegt
% %     man ev. 30\,m tiefer unter einer Brücke.
% %     \item Bei jedem Verkehrsunfall herscht Explosionsgefahr
% %     auf Grund des explosiven Gemisch im Tank bzw. der Benzinleitung. Daher in diesem Bereich
% %     nicht rauchen!
% %     \item Eine Besonderheit stellen Gefahrenguttransporter dar (\See{topic:gefahrengutunfall}). 
% % Sie sind -- wie es ihr Name schon sagt -- mit gefährlichen Gütern versehen und tragen
% % somit eine höhere Gefahrenpotentialität in sich. Zu erkennen sind
% % Gefahrenguttransporter durch eine orange, rechteckige Warntafel mit
% % schwarzen Rand. 
% % 
% % 
% % 
% % \begin{tabular}{m{3.225cm}m{3.225cm}m{3.225cm}m{3.225cm}}
% % \multicolumn{2}{m{6.65cm}}{\centering Allgemeine Kennzeichnung (Sammeltransporte)} 
% % &
% % \multicolumn{2}{m{6.65cm}}{\centering Spezielle
% % Kennzeichnung} 
% % \\\midrule
% % \includegraphics[width=\linewidth]{./images/khd/warntafel-allgemein.png}
% %  &
% %  &
% % \includegraphics[width=\linewidth]{./images/khd/warntafel-un-60-1710.png}
% %  &
% % Gefahrnummer\par (\T{Kemler-Nummer})
% % 
% % Stoffnummer\par (\T{UN-Nummer})
% % \\\bottomrule
% % \end{tabular} 
% % 
% % In der Ausführung gibt es neben der unbeschrifteten Tafel eine mit
% % zwei Zifferkombinationen ausgestattete Form der Warntafel. Die
% % Beschriftung zeigt in seinem oberen Feld die Gefahrennummer
% % (Kemmlernummer) und in seinem unteren Bereich die Stoffnummer
% % (UN-Nummer) an. Die Nummern sind beim Notruf anzugeben.
% % 	\item Bei einem beginnenden Fahrzeugbrand kann durch stoßweises Sprühen mit
% % 	einem Feuerlöscher direkt auf den Brandherd	bekämpft werden.
% % 	\item Bei brennende Kleidungsstücke sollten die Flammen sofort erstickt werden
% % 	und wenn möglich die verbrannte Kleidung vom Patient entfernen.
% % 	\item Achtung beim Zimmerbrand. Nicht ungeschützt die Türklinke angreifen
% % 	(heiss!) und die Türe nur langsam öffnen. Den Raum nur zum Retten von Personen
% % 	betreten. Nicht atmen! Sauerstoffmangel! Kein Fenster öffnen. Dadurch kommt
% % 	nur wieder Sauerstoff in den Raum und Flammen benötigen diesen zum Leben.
% % 	\item In einem gasgefüllten Raum kein Licht aufdrehen, nicht telefonieren,
% % 	klingeln und rauchen. Hier darf bzw. soll das Fenster geöffnet werden. Wenn
% % 	möglich die Gaszufuhr abdrehen. Auch hier gilt, den Raum nur zum Retten von
% % 	Patienten betreten.
% % 	\item 
% % \end{itemize}
% % 
% % 
% % 
% 
% 
% 
% 
% 
% 
% \Layout{0}{Gefahrenquellen}
% {
% Eine der größten Gefahren für den Ersthelfer auf Straßen oder
% Autobahnen ist das \E{Überfahrenwerden} durch
% nachkommende Fahrzeuge. Beim Übersteigen einer Leitschiene muss vorsichtig geprüft werden, ob
% auf der anderen Seite auch ein fester Untergrund vorhanden ist (sonst liegt
% man ev. 30\,m tiefer unter einer Brücke).
% Bei jedem Verkehrsunfall herrscht \Ca{Explosionsgefahr}
% auf Grund des explosiven Luft-Gas-Gemisches im Tank bzw. in der
% Benzinleitung. Bei Verkehrsunfällen besteht
% absolutes
% \EE{Rauchverbot!}
%     
% 
% 
% Ein beginnender \EE{Fahrzeugbrand} kann unter Umständen
% durch stoßweises Sprühen direkt auf den Brandherd mit einem Feuerlöscher bekämpft
% werden. Sollten \EE{Kleidungsstücke} brennen, werden die Flammen
% sofort mit einer Decke o.ä. erstickt. Wenn möglich
% wird die verbrannte Kleidung vom Patienten
% entfernt.\footnote{Nur wenn leicht möglich. Mit der Haut
% verbundene, "`eingebrannte"' Kleidungstücke müssen belassen werden.}
% 
% Achtung beim \EE{Zimmerbrand}! Die
% Türschnalle nicht ungeschützt angreifen (heiß!), die Türe
% nur langsam öffnen und den Raum nur zum Retten von Personen
% betreten. Aufgrund der Rauchgase und des Sauerstoffmangels
% darf \E{nicht geatmet werden!} Bereits ein bis zwei
% Atemzüge können eine Bewusstlosigkeit zur Folge haben. Es
% darf kein Fenster geöffnet werden, sonst käme wieder
% Sauerstoff in den Raum und das Feuer würde erneut angefacht
% werden.
% 
% In mit \EE{Gas gefüllten Räumen}\footnote{Austreten von Gas aus der
% Gasleitung, Stickgase durch defekte Heizungen, Gärkeller, \ldots} besteht
% akute Explosions- und Erstickungsgefahr. Jeder Funke muss vermieden werden! Es
% darf daher kein Licht aufgedreht, nicht telefoniert, geklingelt oder geraucht
% werden. Das Fenster ist wenn möglich zu öffnen und die Gaszufuhr abzudrehen.
% Auch hier gilt es, den Raum nur zum Retten von Patienten zu betreten und nicht zu atmen!
% 
% Bei einem \EE{Stromunfall} wird zwischen
% \E{Niederspannung} und \E{Hochspannung} unterschieden. Bei einem
% Niederspannungsunfall kann der Ersthelfer versuchen das Gerät auszuschalten,
% den Stecker zu ziehen oder die Sicherung herauszudrehen. Zu guter Letzt bleibt
% noch die Möglichkeit, den Verunglückten von der Stromquelle loszureißen. Bei
% einem Hochspannungsunfall muss über die
% Polizei das Energieversorgungsunternehmen informiert und die Abschaltung
% veranlasst werden.
% 
% Viele Gefahren laueren aber auch bei \EE{Freizeit- und Sportunfällen}. Zum
% Beispiel beim Eislaufen auf Natureis (Seen und Teiche). Kommt es zu einem
% \E{Einbruch druch Eis} hat sich der Ersthelfer flach am
% Bauch liegend mit Brett, Stange oder Leiter dem Patienten zu
% nähern. Es dürfen keine Rettungsversuche ohne Absicherung
% durchgeführt werden! Auf  \E{Skipisten} gelten \E{zwei
% gekreuzt aufgestellte Ski} als Warnzeichen
% und Absicherung.
% 
% Wenn ein \EE{Verschütteter} gerettet werden soll, muss der Ersthelfer darauf
% achten, dass kein Schüttmaterial nachrutscht. Es ist mindestens der Kopf und
% der Brustkorb des Patienten freizuschaufeln, um
% die Atmung zu ermöglichen.
% }{
% \begin{itemize}
%   \item Straßenverkehr:
%   \begin{itemize}
%     \item Überfahrenwerden
%     \item Explosions- und Brandgefahr
%   \end{itemize}
%   \item Gefahrgut: Warntafel beachten und beim Notruf mitteilen
%   \item Beginnender Fahrzeugbrand: Feuerlöscher
%   \item Brennende Kleidung: Flammen ersticken; vom Patienten entfernen
%   \item Zimmerbrand: kein Fenster öffnen; nur zum Retten betreten
%   \item Gas: Funken vermeiden: Kein Licht etc., Explosions- und
%   Erstickungsgefahr!
%   \item Stromunfall:
%   \begin{itemize}
%     \item Niederspannung: Gerät abschalten; Stecker herrausziehen;
%     Sicherungen herrausdrehen
%     \item Hochspannung: Abschalten durch Fachleute veranlassen
%   \end{itemize}
%   \item Freizeit- und Sportunfälle:
%   \begin{itemize}
%     \item Einbruch auf Eis: Flach mit Brett oder Leiter nähern; Nie ungesichert!
%     \item Ski-Unfälle: Absicherung mit zwei gekreuzten Skiern
%   \end{itemize}
%   \item Verschütteter: Kopf und Brustkorb freimachen (Atmung)
% \end{itemize}
% }{}{}{}

\subsubsection{Gefahrengutransport}

Eine Besonderheit stellen \EE{Gefahrenguttransporte} dar
(\See{topic:gefahrengutunfall}). Sie sind -- wie der Name
schon sagt -- mit gefährlichen Gütern versehen und haben
somit ein erhöhtes Gefährdungspotential. Zu erkennen sind
Gefahrenguttransporter durch eine orange, rechteckige
Warntafel mit schwarzem Rand
(\prettyref{tab:gefahrentafel-ehi}).

\begin{table}[H]
\Captionof{table}{In der Ausführung gibt es neben der unbeschrifteten Tafel
eine mit zwei Ziffernkombinationen ausgestattete Form der Warntafel. Die
Beschriftung zeigt in seinem oberen Feld die Gefahrennummer
(Kemler-Nummer) und in seinem unteren Bereich die Stoffnummer
(UN-Nummer) an. \E{Die Nummern sind beim Notruf
anzugeben}.\label{tab:gefahrentafel-ehi}}

\begin{tabular}{m{3.225cm}m{3.225cm}m{3.225cm}m{3.225cm}}
\multicolumn{2}{m{6.65cm}}{\centering Allgemeine Kennzeichnung (Sammeltransporte)} 
&
\multicolumn{2}{m{6.65cm}}{\centering Spezielle
Kennzeichnung} 
\\\midrule
\includegraphics[width=\linewidth]{./images/khd/warntafel-allgemein.png}
 &
 &
\includegraphics[width=\linewidth]{./images/khd/warntafel-un-60-1710.png}
 &
Gefahrnummer\par (\T{Kemler-Nummer}\index{Nummer!Kemler-})

Stoffnummer\par (\T{UN-Nummer}\index{Nummer!UN-})
\\\bottomrule
\end{tabular} 

\end{table}


% {
% Die häufigste nicht natürliche Todesursache für junge Menschen
% liegt im Straßenverkehr. Am meisten gefährdet sind jene
% Verkehrsteilnehmer, die ungeschützt am Verkehr teilnehmen. Vorrangig
% zu nennen sind Fußgänger, Radfahrer oder Motorradfahrer. Am
% besten geschützt sind jene Personen, die über eine ausreichend
% starke Schutzzone bzw. Knautschzone in Verbindung mit hoher Masse ihres
% Fahrzeuges verfügen. Vorrangig zu nennen sind LKW-Fahrer oder Züge.
% Nennbare Gefahrenquellen sind:}
% 
% \begin{enumerate}
%     \item  Auffahrunfall
%     \item  Fahrzeugbrand
%     \item  Eisenbahnkreuzung
% \end{enumerate}






% \subparagraph{Weitere Gefahren}
% 
% \Z{
% Die sonst häufig vorkommenden Gefahren liegen im Freizeit- und
% Sportbereich. Beispiele wären:
% 
% \begin{itemize}
%     \item  Sturz auf der Skipiste
%     \item  Absturz beim Klettern
%     \item  Überbelastung bei Ausdauersport
% \end{itemize}
% 
% }






\subsubsection{Retten aus dem Gefahrenbereich}
\label{topic:wegziehen}
\label{topic:rautekgriff}

Durch Anwenden der oben genannten Absicherungsmaßnahmen kann
in vielen Situationen kein vollständiger Schutzbereich errichtet werden.
Insbesondere im Straßenverkehr drohen, solange sich der Patient auf
der Fahrbahn befindet, weiterhin Gefahren, da Absicherungsmittel leicht
vom nachkommenden Verkehr übersehen werden können. Daher
ist der Patient von der Fahrbahn in einen
vollständig sicheren Raum neben der Fahrbahn zu retten. 

\Fazit{Der Patient ist so rasch und schonend wie nur möglich aus der
Gefahrenzone zu bringen.}

% {
% Der Ersthelfer hat bei der Rettung aus dem Gefahrenbereich zusätzliche
% Grundsätze der Ersten Hilfe einzuhalten, die allesamt denselben
% Prinzipen wie die allgemeinen Grundsätze folgen:}
% 
% \subparagraph{1. Grundsatz:} Bei der Rettung darf immer nur Zug auf den
% Patientenkörper ausgeübt    werden, niemals Druck!
% 
% {
% \scriptsize Durch Zug werden die Muskeln des
% Körpers angespannt, wodurch eine Stabilisierung   des Körpers
% eintritt. Des Weiteren können Nervenbahnen mit Zug sehr gut umgehen, 
%  während sie bei Druck schnell beeinträchtigt werden.}
% 
% \subparagraph{2. Grundsatz:} Der Patient ist bei Ausübung von Zug immer nur an
% den Fixpunkten seines   Körpers anzufassen!
% 
% {\scriptsize
%   Als Fixpunkte des Körpers werden Gelenke bezeichnet, da an Ihnen
% eine erhöhte   Stabilität vorzufinden ist. Bei Ausübung von Zug
% an den Gelenken kann der Ersthelfer   bei weniger Kraftaufwand eine
% schneller, für ihn leichtere und den Patienten schonen-  dere Rettung
% durchführen.}
% 
% {
% In Einhaltung dieser Grundsätze haben sich folgende Rettungsmethoden
% entwickelt, die entweder beim sitzenden oder liegenden Patienten
% anzuwenden sind.}




\paragraph{Sitzender Patient:}
%%%%%%%%%%%%%%%%%%%%%%%%%%%%%%%%%%%%%%%%%%%%%%%%%%%%%%%%%%%
% TODO : Neuerstellung
\TODO{GABS}{yyyy-mm-dd}{Neuerstellung von Text notwendig}{}
%%%%%%%%%%%%%%%%%%%%%%%%%%%%%%%%%%%%%%%%%%%%%%%%%%%%%%%%%%%

% \label{topic:rautekgriff}

\TakeNotesLarge

% TODO EHI: Bilder Rautek aus Auto einfügen

% \subparagraph{4.3.1.1) Rautekgriff beim sitzenden Patienten:}
% 
% {
% {Durch das große Engagement eines Wieners namens Franz
% Rautek, konnten in den 1930er Jahren Richtlinien für die
% Personenrettung entwickelt werden, die noch heute allgemeine
% Gültigkeit haben.}\footnote{\ {Franz Rautek (1902-1989): Ein
% aus Wien stammender Jiu-Jitsu-Lehrer, der sich interessehalber mit der
% Rettung von Personen beschäftigte, wobei er im Zuge seiner Recherchen
% ein dermaßen umfangreiches Fachwissen ansammeln konnte, dass er
% schon zu Lebzeiten international als führender Experte gehandhabt
% wurde und noch heute international geltende Richtlinien der
% Personenrettung entwickeln konnte. Er war ua. auch für den
% Arbeiter-Samariter-Bund Österreichs unterrichtend tätig.}}{
% In Anlehnung an diesen Pionier der Ersten Hilfe, wurde der folgende von
% ihm entwickelte Rettungsgriff als
% {\quotedblbase}Rautekgriff{\textquotedblleft} benannt.}}
% 
% 
% { Erklärung des
% Rautekgriffs anhand von Abbildungen auf der linken Seite und
% erläuternden Bemerkungen auf der rechten Seite.}\par
% 
% ~
% 
% { Ablauf (Der Vorgang beschreibt die
% Rettung des Fahrers eines links gesteuerten Fahrzeuges):}
% 
% \begin{enumerate}
%     \item  Tür des Fahrzeugs öffnen
%     \item  Patient ansprechen und an der
%     Schulter berühren
%     \item  Halbmondgriff: um den Brustkorb
%     des Patienten in Form eines Halbmondes umwandern und dabei drei
%     Teilschritte durchführen:
%     \begin{enumerate}
%         \item  Schlüssel abziehen und auf das
%         Armaturenbrett legen
%         
%         {\scriptsize Bestimmte Fahrzeuge besitzen
%         keine klassische Schlüsselzündung, sondern einen Initialdruckknopf
%         \ oder eine \ Kartenzündung. In diesem Fall ist durch Drücken des
%         Initialdruckknopfs oder Herausziehen der Fahrzeugkarte der laufende
%         Motor abzustellen.}
%         \item  Handbremse anziehen
%         
%         {\scriptsize Manche Fahrzeugtypen oder
%         --marken besitzen nicht die klassische Feststellbremse in Form eines
%         Hebels in der Mittelkonsole. An Stelle dieses Hebels kann es ein
%         viertes Pedal links neben der Kupplung oder einen herauszuziehenden
%         Hebel rechts unter dem Lenkrad geben. }
%         
%         \item  Gurt öffnen
%         
%         {\scriptsize Sollte der Gurt nicht zu
%         öffnen sein, kann sich der Ersthelfer behelfen durch Rückschieben
%         des Sitzes oder Aufschneiden des Gurtes mittels Schere oder Gurtmesser.
%         Eine Schere findet sich standardmäßig in einer
%         ÖNORM-Autoapotheke, ein Gurtmesser hingegen müsste extra besorgt
%         werden. Bevor der Gurt geöffnet wird ist der Patient, am besten an
%         der linken Schulter, anzufassen, um den Patienten in seiner Position zu
%         stabilisieren. Auch in den nächsten Schritten ist immer darauf zu
%         achten, dass eine Hand den Patienten stabilisiert.}
%     \end{enumerate}
%     
%     \item  Patient soweit nach vorne lehnen,
%     dass er mit dem Kopf am Lenkradrand ankommt. Dabei wird mit der linken
%     Hand die Stirn stabilisiert, während die rechte Hand zwischen den
%     Schulterblättern die Vorwärtsbewegung ausübt.
%     \item  linke Hand wandert auf die linke
%     Schulter, rechte Hand fasst die rechte Hüfte an
%     \item linke Hand wandert zur linken
%     Hüfte, mit beiden Händen den Rücken des Patienten zum Ersthelfer
%     drehen
%     \item Knie und Oberschenkel zur
%     Wirbelsäule des Patienten stellen
%     \item mit beiden Armen unter die Achseln
%     des Patienten greifen, einen der beiden Arme in Nähe der Gelenke
%     anfassen. Dabei ist darauf zu achten, dass die eigenen Daumen nicht
%     zwischen Arm und Brustkorb des Patienten eingeengt werden.
%     \item  Patient in einer kurzen
%     Hebebewegung auf den Oberschenkel des Ersthelfers heben und auf dem
%     Oberschenkel stabilisierend vom Fahrzeug wegziehen
%     
%     
%     
% \end{enumerate}

% {\bfseries
% 4.3.2) liegender Patient:}

\paragraph{Liegender Patient}
%%%%%%%%%%%%%%%%%%%%%%%%%%%%%%%%%%%%%%%%%%%%%%%%%%%%%%%%%%%
% TODO : Neuerstellung
\TODO{GABS}{yyyy-mm-dd}{Neuerstellung von Text notwendig}{}
%%%%%%%%%%%%%%%%%%%%%%%%%%%%%%%%%%%%%%%%%%%%%%%%%%%%%%%%%%%

\TakeNotesLarge

% TODO EHI: Bilder Wegziehen einfügen

% 4.3.2.1) Scherengriff:
% 
% {
% Bei der Rettung eines liegenden Patienten besteht im Unterschied zur
% Rettung aus sitzender Position die erhöhte Gefahr einer Verletzung am
% Kopf des Patienten. Würde man den Patienten der Einfachheit halber an
% den Knöcheln oder den Händen anfassend wegziehen, würde der Kopf
% im ersten Fall am Boden nachschleifen oder im zweiten Fall mit stark
% überdehnter Halswirbelsäule zu Boden hängen und uU. darauf
% aufschlagen. Um dieser Gefahr entgegenzutreten, wurde eine
% Rettungstechnik entwickelt, bei der die Arme des Patienten wie die
% Klingen einer Schere überkreuzt werden, um dadurch den Kopf
% stabilisierend zwischen beiden Oberarmen einzuengen und so an den
% Handgelenken anfassend den Patienten aus dem Gefahrenbereich
% wegzuziehen. }
% 
%  { Erklärung des
% Scherengriffs anhand von Abbildungen auf der linken Seite und
% erläuternden Bemerkungen auf der rechten Seite.}\par
% 
% { Ablauf:}
% 
% \begin{enumerate}
%     \item den linken Arm des Patienten mit
% dem Oberarm parallel zum Kopf legen
% 
% {\scriptsize Prinzipiell kann auch mit dem
% rechten Arm begonnen werden, der weitere Verlauf muss aber
% spiegelverkehrt ablaufen.}
%     \item das linke Knie auf die linke
% Außenseite des Oberarms vom Patienten stellen, danach den Kopf des
% Patienten auf den Oberarm hinauflegen und mit der linken Hand die
% rechte Gesichtshälfte des Patienten festhalten
%     \item mit der rechten Hand den rechten
% Arm des Patienten mit dem Oberarm parallel zum Kopf legen, dabei darauf
% achten, dass die rechte Hand des Patienten auf den eigenen linken
% Oberschenkel zu liegen kommt, danach wird an dem rechten Handgelenk des
% Patienten anfassend der Arm soweit nach links gezogen, bis der Kopf
% zwischen den beiden Oberarmen eingeengt ist.
%     \item linke Hand greift auf des rechte
% Handgelenk des Patienten und übernimmt den zuvor aufgebauten Zug;
% rechte Hand greift auf das linke Handgelenk und zieht nach rechts
%     \item die Hände des Patienten wie bei
% einer Schere zusammenführen, dabei darauf achten, dass der Kopf
% zwischen den Oberarmen eingeengt bleibt
%     \item Patient aus dem Gefahrenbereich
% herausziehen
% 
% \end{enumerate}

% {
% Der Scherengriff hat neben dem Vorteil einer Stabilisierung des Kopfes
% samt Halswirbelsäule den weiteren Vorteil einer Stabilisierung der
% gesamten restlichen Wirbelsäule, da hierbei der Körper möglichst
% flach am Boden weggezogen werden kann ohne Gefahr zu laufen den
% Körper zu sehr zu beugen. Des Weiteren kann der Ersthelfer beim
% Wegziehen unterstützend sein Körpergewicht in die Zugrichtung
% hineinlegen, wodurch er sogar schwere Personen befördern kann. Daher
% ist der Scherengriff dem nachstehenden Rautekgriff grds. vorzuziehen.}
% 
% \Z{\scriptsize
% In der Praxis hat sich gezeigt, dass es insbesondere eine Personengruppe
% gibt, bei der ein Scherengriff nicht umsetzbar ist. Grund für das
% Scheitern ist eine Fehlstellung des Schultergelenks, wodurch diese
% Personen nicht in der Lage sind mit den Oberarmen den Kopf zu
% berühren. Aufgrund dieser Körperfehlstellung kann der Scherengriff
% nicht umgesetzt werden und muss durch den Rautekgriff oder einer seiner
% Alternativen ersetzt werden.}


% 
% 
% \paragraph{
% 4.3.2.2) Rautekgriff beim liegenden Patienten und Alternativen:}
% 
% {
% Das Prinzip dieses Rautekgriffs ist ident zu dem beim sitzenden
% Patienten. In der Ausführung unterscheidet sie in wenigen Punkten.}
% 
% { Erklärung des
% Rautekgirffs anhand von Abbildungen auf der linken Seite und
% erläuternden Bemerkungen auf der rechten Seite.}\par
% 
% { Ablauf:}
% 
% \begin{enumerate}
%     \item  Ersthelfer kniet sich hinter den
%     Kopf des Patienten, legt dessen Kopf auf die eigenen Oberschenkel und
%     greift unter den Achseln einen der beiden Arme des Patienten in
%     Gelenksnähe mit herausstehenden Daumen an
%     
%     \item  Ersthelfer lehnt sich mit dem
%     Patienten nach vorne, steht auf und lehnt den Patienten auf einen
%     seiner Oberschenkel
%     
%     \item  Patient aus den Gefahrenbereich
%     herausziehen
% \end{enumerate}
% 
% { Alternativ:}
% 
% \begin{enumerate}
%     \item siehe Schritt 1) oben
%     \item  Ersthelfer steht nicht auf, sondern
%     zieht den Patienten auf den Knien rutschend aus dem
%     Gefahrenbereich
% \end{enumerate}
% 
% 
% {
% Beim Arbeiter-Samariter-Bund Österreichs wird noch eine weitere
% alternative Rettungstechnik gelehrt: Sie funktioniert wie folgt:}
% 
%  { Erklärung des
% ASBÖ-alternativen-Rettungsgriffs anhand von Abbildungen auf der
% linken Seite und erläuternden Bemerkungen auf der rechten Seite.}\par
% 
% { Ablauf:}
% 
% \begin{enumerate}
%     \item Ersthelfer kniet sich hinter den
% Kopf des Patienten, legt dessen Kopf auf die eigenen Oberschenkel und
% greift in die Achseln des Patienten, wobei er darauf achtet mit seinen
% Unterarmen eine Auflage für den Kopf des Patienten zu formen 
% 
%     \item Ersthelfer zieht den Patienten auf
% den Knien rutschend aus dem Gefahrenbereich
% \end{enumerate}


% 
% \paragraph{4.3.2.3) Wegziehen an den Knöcheln:}
% 
% {
% In manchen Fällen ist der Zugang zum Kopf des Patienten durch
% umherstehende, unverrückbare Gegenstände versperrt. Folglich ist
% auch eine Rettung vom Kopf aus des Patienten unmöglich. Entsprechend
% dem Grundsatz der Interessensabwägung muss der Patient trotzdem von
% seiner Position wegbewegt werden, um sein Leben zu retten. In solchen
% absoluten Ausnahmesituationen ist es zulässig den Patienten an den
% freiliegenden Knöcheln anzufassen und soweit wegzuziehen, bis eine
% Rettung vom Kopf aus möglich wird.}



\subsection{Helmabnahme}
\label{topic:sturzhelmabnahme}

\TODO{GABS}{2010-08-01}{Mit Aussendung der ASBÖ-Akademie
(\protect\bibentry{ASBOEAkademieUpdate2010Sturzhelm},
\cite{ASBOEAkademieUpdate2010Sturzhelm}) bezüglich Konformität kontrollieren!}{}

\Layout{0}{\TxBeschreibung}
{Ein Patient mit aufgesetztem Sturzhelm stellt eine eine
besondere Herausforderung dar. Es gilt:
% KORR ZWING 2010-03-xx, ok KOCH
\Fazit{Ein Sturzhelm ist grundsätzlich immer abzunehmen!}

% Bei einem Unfall mit einem motorisierten einspurigen Kraftfahrzeug kann es
% vorkommen, dass ein Patient mit aufgesetztem Sturzhelm vorgefunden wird. 
% 

% In der
% Regel wird davon ausgegangen, dass der Verunglückte bei Bewußtsein ist und
% den Helm schon selbst abgenommen hat. 

% Nach dem Retten aus der Gefahrenzone muss die Abnahme des
% Sturzhelms rasch und schonend erfolgen. Das Visier ist zu
% öffnen und der Verletzte anzusprechen. Danach wird der Helm
% abgenommen.



Die Annäherung an den Patienten erfolgt in dessen
Blickrichtung, um ihn nicht zu überraschen. Bevor der
Patient angesprochen wird, muss der Helm fixiert
werden. Der Patient wird daraufhin aufgefordert sich nicht
zu bewegen. 
% \SpeechMargin{"`Bitte bleiben Sie
% ruhig liegen und bewegen Sie sich nicht!"'}

Am besten wird der Sturzhelm von zwei Helfern abgenommen
(\II{Zweihelfermethode}), zur Not funktioniert es jedoch
auch mit nur einem Helfer (\II{Einhelfermethode}). Der Kopf
muss in Neutralposition stabilisiert werden. Auf
die \EE{Halswirbelsäule} darf \EE{kein Zug} ausgeübt
werden und sie darf \EE{nicht gestaucht}
werden!~\cite{ITLSde6,AsboeLehrmeinung2011-07}

Im professionellen Umfeld muss bei Verdacht auf eine
Verletzung der (Hals-)Wirbelsäule (\prettyref{m:einschaetzung-indikation-ws-immobilisation})
anschließend eine HWS-Schiene angelegt werden
(\prettyref{topic:hws-schiene-anlegen}). 
}{
\begin{itemize}
  \item Zweihelfermethode (1. Wahl)
  \item Einhelfermethode
  \item Ziel:
  \begin{enumerate}
    \item Kopf in Neutralposition
    stabilisieren
    \item HWS nicht stauchen
    \item Keinen Zug ausüben
  \end{enumerate}
\end{itemize}
}
{}
{}
{}


% \subsubsection{Zweihelfermethode:}

\Layout{57}{}
{}{}
{
\begin{tabularx}{\linewidth}{XXX}
\includegraphics[width=\linewidth]{./images/koch-aass/00800-gamma/HelmZweihelfer1b.JPG}
&
\includegraphics[width=\linewidth]{./images/koch-aass/00800-gamma/HelmZweihelfer2b.JPG}
&
\includegraphics[width=\linewidth]{./images/koch-aass/00800-gamma/HelmZweihelfer3b.JPG}
\\
% Schritt 1 
Bewegungsverbot: \Speech{"`Bitte bleiben Sie ruhig liegen
und bewegen Sie sich nicht!"'} 
Kopf und Helm fixieren, Visier öffnen
\Speech{"`Mein Name ist {\ldots} Wir
werden jetzt Ihren Helm abnehmen."'} 
&
% Schritt 2 
Ansprechen, Kinnriemen öffnen
&
% Schritt 3 
Kopf stabilisieren, Helm abziehen
\\\addlinespace
\includegraphics[width=\linewidth]{./images/koch-aass/00800-gamma/HelmZweihelfer4b.JPG}
&
\includegraphics[width=\linewidth]{./images/koch-aass/00800-gamma/HelmZweihelfer5b.JPG}
&
\includegraphics[width=\linewidth]{./images/koch-aass/00800-gamma/HelmZweihelfer6b.JPG}
\\
% Schritt 4 
Helm abziehen \ACh{GABS}{2010-02-19}{Helm kippen bis
Nasenspitze sichtbar?}
&
% Schritt 5 
Stabilisierung übernehmen
&
% Schritt 6
Kopf halten: Ohne Zug, nicht stauchen!.

\Z{(Ggfs. HWS-Schiene anlegen,
\prettyref{topic:hws-schiene-anlegen})

\Speech{"`Bitte bleiben Sie noch ruhig liegen, wir werden
Ihnen jetzt eine Halskrause anlegen, um Ihren Hals zu
stabilisieren."'}}
\\\addlinespace\bottomrule
\end{tabularx}
}{Bilderserie: Helmabnahme zu zweit
(Zweihelfermethode)\index{Bilderserie!Helmabnahme
(Zweihelfermethode)}\index{Helmabnahme!Bilderserie (Zweihelfermethode)}}{}

% \subsubsection{Einhelfermethode:}

\Layout{57}{}
{}{}
{
\begin{tabularx}{\linewidth}{XXX}
\includegraphics[width=\linewidth]{./images/koch-aass/00800-gamma/HelmEinhelfer1b.JPG}
&
\includegraphics[width=\linewidth]{./images/koch-aass/00800-gamma/HelmEinhelfer2b.JPG}
&
\includegraphics[width=\linewidth]{./images/koch-aass/00800-gamma/HelmEinhelfer3b.JPG}
\\
% Schritt 1 
Bewegungsverbot: \Speech{"`Bitte bleiben Sie ruhig liegen und bewegen Sie
sich nicht!"'}
Kopf und Helm fixieren, Visier öffnen
&
% Schritt 2 
Ansprechen, Kinnriemen öffnen

\Speech{"`Mein Name ist {\ldots} Ich
werde jetzt Ihren Helm abnehmen."'} &
% Schritt 3 
Helm kippen \ldots
\\\addlinespace
\includegraphics[width=\linewidth]{./images/koch-aass/00800-gamma/HelmEinhelfer4b.JPG}
&
\includegraphics[width=\linewidth]{./images/koch-aass/00800-gamma/HelmEinhelfer5b.JPG}
&
\includegraphics[width=\linewidth]{./images/koch-aass/00800-gamma/HelmEinhelfer6b.JPG}
\\
{\ldots} bis
Nasenspitze sichtbar ist.
\ACh{Instruktor 1}{0.8.0}{Weg lassen. }\ACh{GABS}{2010-02-19}{?} 
&
% Schritt 5 
Kopf stabilisieren, Helm abziehen
&
% Schritt 6
Kopf ablegen

\Z{(Ggfs. HWS-Schiene anlegen)

\Speech{"`Bitte bleiben Sie noch ruhig liegen, wir werden
Ihnen jetzt eine Halskrause anlegen, um Ihren Hals zu
stabilisieren."'}}
\\\addlinespace\bottomrule
\end{tabularx}
}{Bilderserie: Helmabnahme alleine
(Einhelfermethode)\index{Bilderserie!Helmabnahme
(Einhelfermethode)}\index{Helmabnahme!Bilderserie (Einhelfermethode)}}{}

 


% {
% Ein Kraftradfahrer mit aufgesetztem Helm ist im Zuge der
% ersthelferischen Maßnahmen an manchen Stellen besonders zu
% behandeln. Auffälligster Unterschied besteht in Zusammenhang mit der
% Helmabnahme, da hierbei zwischen dem Ruf um Hilfe und der
% Verdachtskontrolle der Atemwege ein eigener Behandlungsschritt nur
% für das Abnehmen des Helms besteht. Im Zuge der vorhergehenden
% Behandlungen ist der Kraftradfahrer grds. gleich wie die anderen
% Patienten zu behandeln. Auch seine Unfallstelle muss abgesichert werden
% und er selbst muss gerettet werden. Bei der Bewusstseinskontrolle muss
% hingegen vor dem Ansprechen das Visier geöffnet werden. Der Kopf darf
% bei aufgesetztem Helm nicht überstreckt werden. }
% 
% {
% Der Helm ist ein Schutzobjekt, das seine Wirkung nur während des
% Aufpralls verwirklicht. Nach Abschluss des Unfalls beeinträchtigt er
% den Patienten, da er einen stetigen Druck auf den Kopf ausübt, den
% Mundraum unvollständig versperrt und eine genauer Untersuchung und
% Behandlung unmöglich macht. \ Dadurch hat sich ein weiterer Grundsatz
% entwickelt.}
% 
% \EE{Der Helm ist immer abzunehmen, außer der Patient verweigert dies!}
% 
% {
% Bei der Umsetzung sind zwei Möglichkeiten zulässig. Bei der einen
% agiert ein Ersthelfer alleine (Einhelfermethode), bei der anderen
% arbeiten zwei Ersthelfer zusammen (Zweihelfermethode). Vorteil der
% Zweihelfermethode ist eine höhere Stabilisierung des Kopfes und der
% Halswirbelsäule.}
% 
% 
% { Erklärung der
% Einhelfermethode anhand von Abbildungen auf der linken Seite und
% erläuternden Bemerkungen auf der rechten Seite.}\par
% 
% { Ablauf:}
% 
% \begin{enumerate}
%     \item 
% 
% { 1) Ersthelfer geht von vorne auf den
% Patienten zu, spricht ihn an und berührt ihn bei geöffneten Visier}
% 
% { 2) Ersthelfer positioniert sich
% hinter dem Kopf des Patienten, fasst mit einer Hand den Helm an und
% öffnet mit der anderen den Kinnriemen}
% 
% {\itshape Ab nun darf der Helm bis zur
% beendeten Abnahme nicht mehr losgelassen werden. Der Kinnriemen ist
% meistens durch einen roten Druckknopf zu öffnen. Als Alternativen zum
% Druckknopf gibt es am Markt ua. auch Bajonettverschlüsse.}
% 
% { 3) die Finger des Ersthelfers greifen
% unter die Laschen des Kinnriemens und ziehen den Helm leicht
% auseinander, die Daumen wandern in das Visier}
% 
% { 4) die eigentliche Abnahme besteht
% aus drei Schritten:}
% 
% { 4.1) Helm von oben zum Ersthelfer hin
% drehen, bis die Nasenspitze des Patienten zu sehen ist}
% 
% { {4.2) Helm von unten zum Ersthelfer hin
% drehen }{und}{ ziehen, bis kurz vor der äußersten
% Wölbung des Hinterhauptes }}
% 
% {\itshape Bemerkbar ist dies durch den
% rapide leichter werdenden Helm und dem Gefühl der Kopf würde bei
% weiterem Ziehen herausfallen.}
% 
% { \ 4.3) mit den Händen
% 90{\textdegree} versetzt umgreifen, sodass die eine Hand das
% Hinterhaupt festhält, während die andere Hand die Spitz des Helm
% stabilisiert; mit beiden Händen gleichmäßig den Helm gerade zum
% Ersthelfer hin abziehen}
% 
%  5) Helm stabil auf die Seite und Kopf
% sanft auf den Boden legen
% 
% \end{enumerate}
% 
% { Erklärung der
% Zweihelfermethode anhand von Abbildungen auf der linken Seite und
% erläuternden Bemerkungen auf der rechten Seite.}\par
% 
% { Ablauf:}
% 
% \begin{enumerate}
%     \item 
% 
% { 1) Helfer 1 positioniert sich hinter
% dem Kopf des Patienten und hält den Helm beidseitig fest, Helfer 2
% geht von vorne auf den Patienten zu, spricht ihn an und berührt ihn
% bei geöffnetem Visier.}
% 
% { 2) Helfer 2 öffnet den Kinnriemen
% und stabilisiert die Halswirbelsäule durch den Doppelten-C-Griff.}
% 
% {\scriptsize Beim Doppelte-C-Griff werden
% beide Hände zu einem gebärdensprachlichen C geformt, wobei die eine
% Hand auf den Boden aufgestützt das Genick stabilisiert, während die
% andere Hand den vorderen Halsbereich
% umschließt. Beide Hände üben einen Zug des Kopfes vom Körper
% weg aus.}
% 
% { 3) Helfer 1 greift unter die Laschen
% des Kinnriemens und übt die ersten beiden Abnahmeschritte der
% Einhelfermethode aus.}
% 
% { 4) Helfer 2 fährt mit seiner
% unteren Hand das Genick entlang bis er das Hinterhaupt umfassen kann,
% dadurch braucht der Helfer 1 nicht den dritten Abnahmeschritt der
% Einhelfermethode durchzuführen.}
% 
% { 5) Helfer 1 fasst den freigelegten
% Kopf links und rechts an und übt einen Zug aus, Helfer 2 kann den
% Doppelten-C-Griff lösen und sich mit der Untersuchung des Patienten
% beschäftigen.}
% 
% \end{enumerate}






\subsection{Kontrolle der Vitalfunktionen}
\label{topic:lebensfunktionen-kontrolle}

\Layout{0}{Kontrolle der Lebensfunktion
}{
Nachdem der Patient aus dem Gefahrenbereich gerettet wurde, ist eine
\E{Kontrolle seiner Lebensfunktionen} notwendig, um eine genauere
Einschätzung des Gesundheitszustandes zu erlangen. 
Bei der Kontrolle
beschränkt man sich auf zwei einfach festzustellende
Vitalwerte -- das \EE{Bewusstsein} und die \EE{Atmung}. 

Eine Kontrolle des
Kreislaufs, durch Tasten des Pulses an der Halsschlagader (oder irgendeiner anderen
Stelle), wird normalerweise von einem Ersthelfer \E{nicht}
durchgeführt. Die Pulskontrolle ist 
% in Anbetracht der Beziehung
% zwischen Kreislauf und Atmung 
für einen Ersthelfer auch gar nicht
notwendig, denn Kreislauf und Atmung sind derart eng miteinander
verknüpft, dass ohne dem Einen das Andere nicht funktioniert. 
}{
\begin{itemize}
  \item Bewusstseinskontrolle
  \item Atemkontrolle
\end{itemize}
}{}{}{}



% Um den Kreislauf und die Atmung sowie ihre Beziehung zueinander
% besser zu verstehen, empfiehlt es sich das passende Kapitel im Lehrbuch
% weiter hinter genau durchzulesen und anschließend an dieser Stelle
% weiter zu lesen.


% Bei der Bewusstseins- und
% Atemkontrolle (BAK)
% \footnote{BAK: die Abkürzung BAK bezeichnete in ihrer
% ursprünglichen Form die Wörter: Bewusstsein,
% Atmung und Kreislauf. Sie stammt
% aus einer Zeit in der standardmäßig neben der Bewusstseins- und
% Atemkontrolle auch die Kreislaufkontrolle vom Ersthelfer
% durchzuführen war. Die Abkürzung hat sich noch heute im Fachjargon
% als Bezeichnung für die allgemeine Vitalwertekontrolle gehalten und
% wurde zwischenzeitlich auf das aktuelle Kontrollsystem insoweit
% adaptiert, als sie heute für Bewusstseins- und
% Atemkontrolle steht.} wird
% ein streng einzuhaltender Maßnahmenkatalog vorgeschrieben, der aber
% in jedem seiner Hauptschritte aus drei Teilschritten aufgebaut ist.
% Erster Hauptschritt ist die Bewusstseinskontrolle.

\Layout{0}{Bewusstseinskontrolle
}{
Die Kontrolle baut sich aus folgenden drei Teilschritten auf. Der erste ist das
laute und deutliche \EE{Ansprechen} des Patienten. Idealerweise mit einfachen
Anweisungen wie:

\begin{quote}
   \emph{"`Können Sie mich hören? Machen Sie bitte die Augen
   auf!"'}
\end{quote}

\noindent Gleichzeitig kann man den Patienten an seinen Schultern \EE{berühren}
und leicht rütteln.

\Cave{Patienten, die in Folge eines Unfalls verunglückt
sind, dürfen nicht, oder nur sehr schwach gerüttelt
werden!}

Sollte der Patient darauf nicht reagieren, muss im
Bereich des Kopfes ein \EE{Schmerzreiz} gesetzt werden. Jede Reaktion wird als
Lebenszeichen gewertet. Erfolgt keine Reaktion muss der
Ersthelfer \EE{um Hilfe rufen}, oder besser -- wenn
möglich -- eine zweite Person mit dem Notruf beauftragen.

}{
\begin{itemize}
  \item Ansprechen
  \item Berühren
  \item Schmerzreiz setzen
  \item Ruf um Hilfe
\end{itemize}
}{}{}{}


% 
% \subsubsection{5.1) Bewusstseinskontrolle:}
% 
% % {
% Das Bewusstsein ist die Wahrnehmung der umgebenden Welt und des eigenen
% Selbst. Durch das Bewusstsein kann der Körper eine Vielzahl von
% Funktionen steueren und aufrecht halten. Bei Störungen oder
% gänzlichem Ausfall des Bewusstseins verliert der Patient
% lebenswichtige Kontrollfunktionen über seinen Körper. Dazu gehört
% die Fähigkeit Schmerzen und andersartige belastende Einwirkungen auf
% den Körper wahrzunehmen und sich dagegen zur Wehr zu setzen. Genauso
% entgleitet dem Patienten die Fähigkeit seine Körperspannung (Tonus)
% aufrecht zu halten. Ohne Tonus verfällt die Muskulatur in einen
% Zustande der Schlaffheit. Im Fall der Zunge bedeutet dies, dass die
% Zunge schlaff in den hinteren Rachenraum fällt und die Atemwege
% soweit blockiert, dass eine Aufrechterhaltung der Atmung und folglich
% des Kreislaufs nicht mehr möglich ist. Der Patient würde
% buchstäblich an seiner eigenen Zunge ersticken. Daher ist der Zustand
% von Bewusstlosigkeit auch nicht mit dem Zustand des Schlafens zu
% verwechseln, da beim letzteren der Körper durchgehend einen Tonus
% besitzt.}

% {
% Die Bekämpfung dieser Gefahr wird im nächsten Hauptschritt der
% Atemkontrolle zum ersten Mal wahrgenommen. Zuerst ist es aber für den
% Ersthelfer wichtiger herauszufinden, ob denn überhaupt eine
% Bewusstlosigkeit vorliegt. Die Kontrolle baut sich aus folgenden drei
% Teilschritten auf:}
% 
% \begin{enumerate}
% 
%     \item  Ansprechen:
% 
% {\scriptsize
%  Der Patient ist laut und deutlich mit Worten wie: {\quotedblbase}Hallo!
% Können Sie mich hören!{\textquotedblleft}  anzusprechen. Ein
% Klatschen neben dem Ohr ist alternativ möglich.}
% 
%     \item  Berühren:
% 
% {\scriptsize Der Patient ist an der Schulter anzufassen.}
% 
%     \item  Rütteln:
% 
% {
% {\textit{ Der Patient ist an beiden Schultern zu
% rütteln.}}\footnote{\ {Hierbei besteht ein Unterschied zu
% Angehörigen von Gesundheitsberufen, die an Stelle des Rüttelns
% einen Schmerzreiz in der Wange oder hinter dem Ohr zu setzen haben.
% Warum der Schmerzreiz für den Ersthelfer ASBÖ intern verboten
% wurde, ist weder begründbar noch haltbar, da der richtig angewendete
% Schmerzreiz am meisten zielgerichtet den Unterschied zwischen
% bloßer Bewusstseinsreduzierung und totaler Bewusstlosigkeit
% offenbart und gerade dem unerfahrenen Ersthelfer ein sicheres
% Instrument der Beurteilung ermöglicht.}}{\textit{ }}}
% 
% \Cave{ Cave: Patienten die in Folge eines Unfalls verunglückt sind dürfen
% nicht oder nur sehr  schwach gerüttelt werden!}
% 
% \end{enumerate}

% 
% {
% Sollte der Patient auf die Kontrolle reagieren und folglich bei
% Bewusstsein sein, ist durch Erfragen und/oder Beobachten
% herauszufinden, ob Verletzungen, Erkrankungen oder Vergiftungen
% vorliegen. Ihre Behandlung wird unter dem später folgenden Kapitel
% {\quotedblbase}Weitere Erste Hilfe{\textquotedblleft} näher
% behandelt.}
% 
% {
% Sollte der Patient keine Reaktion abgeben, muss der Ersthelfer davon
% ausgehen, dass eine Bewusstlosigkeit vorliegt. Um die oben erwähnte
% Gefahr einer blockierenden Zunge zu vermeiden, ist die für den
% Ersthelfer einfachste Methode einer Bekämpfung die Überstreckung
% des Kopfes um ca. 30{\textdegree}. Durch diese Maßnahme wird die
% Zunge einem Zug ausgesetzt, der sie in die Länge streckt und damit
% die Atemwege öffnet. }
% 
% {
% Ausgenommen vom sofortigen Überstrecken ist der Kraftradfahrer mit
% aufgesetztem Helm. Bei diesem gelten besondere Regelungen, die unten
% näher behandelt werden.}
% 
% [Warning: Draw object ignored]
% 
% 
% 
% 
% {\bfseries
% 5.1.1) Ruf um Hilfe:}
% 
% {
% Nach Feststellen der Bewusstlosigkeit weiß der Ersthelfer, dass sich
% der Patient in einem lebensbedrohlichen Zustand befindet, für deren
% Bekämpfung er unbedingt weitere, insbesondere professionelle,
% Hilfskräfte benötigt. Da er aber mit der Vitalwertekontrolle noch
% nicht fertig ist und den Patienten auch nicht ungeschützt liegen
% lassen darf, muss er sich eines anderen Ersthelfers zur Notrufabgabe
% bedienen. Um einen solchen Anzusprechen bestehen zwei Möglichkeiten:}
% 
% \begin{enumerate}
%     \item  direkte Ansprache einer Person:
% 
% {\scriptsize
%  Anzuwenden ist diese Form bei sich in unmittelbarer Nähe aufhaltenden
% weiteren Personen.  Diese sind direkt anzusprechen mit den Auftrag
% einen Notruf bei einem der unten stehenden  Notrufnummern abzugeben. }
% 
%     \item  Ruf um Feuer:
% 
% {\scriptsize
%  Sollten sich zwar keine dritten Personen in unmittelbarer Nähe
% aufhalten, aber in hörbarer  Distanz, muss der Ersthelfer durch
% lautes Rufen von {\quotedblbase}Feuer{\textquotedblleft} auf sich
% aufmerksam machen.  Ein Rufen mit dem Wort
% {\quotedblbase}Hilfe{\textquotedblleft} hat leider in der Praxis
% gezeigt, dass sich kaum oder gar keine  Personen bereit erklären
% Hilfe zu leisten. }
% \end{enumerate}
% 
% {
% Sind weder Personen in unmittelbarer Nähe noch hörbarer Distanz
% vorhanden, muss der Ersthelfer mit der Untersuchung fortfahren, um
% weitere Informationen zum Gesundheitszustand sammeln zu können.}
% 
% 
% 
% 
% 
% 
% 
% {\bfseries
% 5.1.3) Verdachtskontrolle der Atemwege:}
% 
% {
% Ist dem Ersthelfer im Zuge der Untersuchung aufgefallen, dass ein Objekt
% offensichtlich (!) die Atemwege verlegt, und er somit den begründeten
% Verdacht einer Atemwegsverlegung hat, darf er den Mundraum des
% Patienten öffnen und blockierende Fremdkörper entfernen. Besteht
% hingegen kein begründeter Verdacht, darf auch nicht kontrolliert
% werden.}
% 
% {
% Anm.: Zahnprothesen (sogen. Dritte) sind zwar Fremdkörper, aber
% blockieren nur bei Abrutschen von ihrer vorgesehenen Position die
% Atemwege. Sie sind daher nur in solchen Situationen zu entfernen. }

\Layout{0}{Atemkontrolle}{%
Damit die Atmung kontrolliert werden kann, muss der Ersthelfer die
\EE{Atemwege freimachen}. Im Verdachtsfall wird die
Mundhöhle auf Fremdkörper kontrolliert. Anschließend wird der Kopf
nackenwärts \EE{überstreckt}, damit die zurückfallende
Zunge die Atemwege nicht blockiert. Nun muss die
Atemkontrolle durch \EE{Hören}, \EE{Sehen} und
\EE{Fühlen} durchgeführt werden. Dafür wird das Ohr über
den Mund des Patienten gelegt und 10~Sekunden lang
überprüft, ob eine normale Atmung vorhanden ist.

%%%%%%%%%%%%%%%%%%%%%%%%%%%%%%%%%%%%%%%%%%%%%%%%%%%%%%%%%%%
% TODO : Neuerstellung Grafik Querschnitt
\TODO{GABS}{yyyy-mm-dd}{Neuerstellung von Grafik notwendig}{}
%%%%%%%%%%%%%%%%%%%%%%%%%%%%%%%%%%%%%%%%%%%%%%%%%%%%%%%%%%%

\Fazit{Es
gilt: Wenn eine Atmung vorhanden ist, dann ist auch ein
Kreislauf vorhanden.}
}{
\begin{itemize}
  \item Atemwege freimachen
  \item Atemkontrolle: 10 Sek
  \begin{itemize}
    \item Hören
    \item Sehen
    \item Fühlen
  \end{itemize}
  \item Wenn Atmung vorhanden, ist auch ein Kreislauf vorhanden
\end{itemize}
}{}{}{}





% \subsubsection{5.2) Atemkontrolle:}
% 
% {
% Als nächster Schritt ist die Atmung zu kontrollieren. Sie muss immer
% bei überstreckten Kopf durchgeführt werden, da ansonsten die den
% Rachen verlegende Zungen jegliches Wahrnehmen einer Atmung verhindern
% würde. Auch bei diesem Hauptschritt haben wir drei Teilschritte zu
% erledigen, wobei sich der Ersthelfer mit seinem Ohr zum Mund des
% Patienten begibt und den Blick auf den Brustkorb/Bauch des Patienten
% richtet:}
% 
% \begin{enumerate}
%     \item  Sehen:
% 
% {
% { }{\textit{Der Ersthelfer beobachtet den Bauch und
% Brustkorb auf Bewegungen der Atmung.}}}
%     \item  Hören:
% 
% {
% { }{\textit{Der Ersthelfer hört die aus dem Mund des
% Patienten entweichenden Atemzüge.}}}
%     \item  Fühlen:
% 
% {
% { }{\textit{Der Ersthelfer fühlt die auf seine Wange
% treffenden Atemzüge.}}}
% \end{enumerate}
% {
% Diese Kontrolle muss für mind. 10 Sek. durchgeführt werden.
% Innerhalb von 10 Sek. atmet ein Erwachsener durchschnittlich 2-3 mal. }
% 
% {
% Nimmt der Ersthelfer eindeutig keine Atemzüge wahr oder ist er sich
% unsicher ob eine ausreichende Atmung vorliegt, muss er davon ausgehen,
% dass der Patient keine Atmung und keinen Kreislauf besitzt. Der Patient
% ist folglich wiederzubeleben. Bewegungen des Patienten die eine Atmung
% lediglich vortäuschen, sind möglich und treten in Form der sogen.
% Schnappatmung auf. Die Schnappatmung ist erkennbar durch langsame, von
% größeren Pausen unterbrochene Atemzüge, die sich entsprechend
% ihres Namens wie das Schnappen nach einem Bissen äußern.}
% 
% {
% Nimmt der Ersthelfer hingegen eindeutig Atemzüge wahr, kann er davon
% ausgehen, dass Atmung und Kreislauf vorhanden sind. Er muss den Patient
% in eine Position bringen, in der dauerhaft die Zunge vom Rachenraum
% befreit wird. Eine solche Position ist die stabile Seitenlagerung.}

\subsection{Notfalldiagnose Bewusstlosigkeit}
\index{Bewusstlosigkeit!Erste-Hilfe}
\label{topic:notfalldiagnose-bewusstlosigkeit}

\Layout{0}{\TxBeschreibung
}{%
Nach der Kontrolle der Vitalfunktionen kann der Ersthelfer
zu der Erkenntnis kommen, dass der Patient zwar \EE{kein
Bewusstsein} hat, aber eine \EE{ausreichende Atmung} und
daher einen \EE{vorhandenen Kreislauf}. In einer solchen
Situation liegt die Notfalldiagnose \T{Bewusstlosigkeit}
vor. Der Patient besitzt keine Schutzreflexe mehr und ist
daher im wahrsten Sinne des Wortes schutz- und hilflos. 
}{
\begin{itemize}
    \item Kein Bewusstsein
    \item Ausreichende Atmung
    \item Vorhandener Kreislauf
\end{itemize}
}{}{}{}

\MassnahmenNG{Erste-Hilfe-Maßnahmen}{Bewusstlosigkeit}{}{
\begin{itemize}
    \item Stabile Seitenlage
\end{itemize}
}

\Layout{57}{}
{}{}
{
\begin{tabularx}{\linewidth}{XXX}
\includegraphics[width=\linewidth]{./images/koch-aass/00800-gamma/Seitenlage1b.JPG}
&
\includegraphics[width=\linewidth]{./images/koch-aass/00800-gamma/Seitenlage2b.JPG}
&
\includegraphics[width=\linewidth]{./images/koch-aass/00800-gamma/Seitenlage3b.JPG}
\\
Schritt 1 &
Schritt 2 &
Schritt 3 \\\addlinespace
\includegraphics[width=\linewidth]{./images/koch-aass/00800-gamma/Seitenlage4b.JPG}
&
\includegraphics[width=\linewidth]{./images/koch-aass/00800-gamma/Seitenlage5b.JPG}
&
\\
Schritt 4 &
Schritt 5 &
\\\addlinespace\bottomrule
\end{tabularx}
}{Bilderserie: Stabile
Seitenlage\index{Bilderserie!Stabile Seitenlage}\index{Stabile
Seitenlage!Bilderserie}}{}



% \Layout{0}{Lebensrettende Sofortmassnahme
% }{
% In einer solchen Situation ist der Patient in eine Position zu bringen in der
% die Zunge dauerhaft die Atemwege nicht verlegt, sonst würde er daran
% ersticken. Dies ist erreichbar durch die sogenannte \T{stabile
% Seitenlage}.
% }{
% \begin{itemize}
%     \item Stabile Seitenlage.
% \end{itemize}
% }



% { Erklärung der stabilen
% Seitenlagerung anhand von Abbildungen auf der linken Seite und
% erläuternden Bemerkungen auf der rechten Seite.}\par
% 

\Layout{2}{Stabile Seitenlage}{\index{Stabile Seitenlage}
Jeder Ersthelfer ist in der Lage mit den nachfolgend beschriebenen, einfachen
Handgriffen eine \T[Stabile Seitenlage]{stabile Seitenlage} herzustellen.
\begin{enumerate}
  \item  Den zum Ersthelfer näherliegenden Arm
  des Patienten im rechten Winkel auflegen.
  \item  Unter das gegenüberliegende
  Knie des Patienten greifen und anheben. Das
  Handgelenk des gegenüber liegenden Arms des Patienten
  auf das Kniegelenk des Patienten geben.
  \item  Hand- und Kniegelenk des Patienten
  zum Ersthelfer drehen.
  \item  Kopf des Patienten
  überstrecken
\end{enumerate}
Durch die stabile Seitenlage bewegt sich die Zunge auf Grund der Schwerkraft
vom Rachenraum weg und die Atemwege werden dadurch frei. Bei Erbrechen fließt
der Großteil des Erbrochenen in der stabilen Seitenlage allein aus dem Mund-
und Rachenraum. Die restlichen Bestandteile sind jedoch vom Ersthelfer
mittels eines um zwei Finger gewickelten Tuches auszuräumen.
}{}{}{}{}


\Layout{0}{Weiterführende Maßnahmen
}{
Nachdem der Patient in die stabile Seitenlage gebracht wurde, muss der
Ersthelfer weiterführende Maßnahmen setzen, um Verschlechterungen
des Gesundheitszustandes zu vermeiden. Sehr wichtig ist der \EE{Wärmeerhalt}. Der
Ersthelfer kann  sich hierfür unterschiedlicher Mittel bedienen, wie z.\,B. die,
in einer Autoapotheke gemäß \Abk{ÖNORM} befindliche Rettungsdecke. Weiters muss eine
\EE{laufende Kontrolle der Atmung} durchgeführt werden. Der Ersthelfer hat
sich alle 30 Sekunden durch Überprüfung der Atemtätigkeit
davon zu überzeugen, dass der Patientenzustand unverändert geblieben ist. Durch Öffnen eines Fensters in
geschlossenen Räumen, wird dem Patienten die benötigte \EE{Frischluft}
zugeführt. Auch der \EE{psychischen Betreuung} des Patienten
durch den Ersthelfer wird eine große Bedeutung beigemessen.
}{
\begin{itemize}
  \item Wärmeerhalt: zudecken
  \item Laufende Kontrolle der Atmung alle 30\Sec*
  \item Frischluftzufuhr: Fenster öffnen
  \item Psychische Betreuung
\end{itemize}
}{}{}{}




% \subsubsection{7.1) anschließende Behandlung:}
% 
% {
% Nach Positionierung des Patienten in die stabile Seitenlage muss der
% Ersthelfer weiterführende Maßnahmen setzen, um Verschlechterungen
% des Gesundheitszustandes zu vermeiden. Die nun folgenden drei
% Maßnahmen sind {\quotedblbase}ersthelferische
% Allzweckwaffen{\textquotedblleft}, da sie nach jeder Grundbehandlung am
% Patienten angewendet werden sollen, um für den Patienten eine
% durchgehende Versorgung gewährleisten. Ein dabei positiv auftretender
% Nebeneffekt ist, dass auch der Ersthelfer in Bewegung bleibt, seine
% Nervosität wegarbeiten kann und sich nicht als nutzlos vorkommt. Die
% Maßnahmen lauten:}
% 
% \begin{enumerate}
%     \item  Wärmeerhalt sichern:
%     
%     {\scriptsize
%     Bei 1{\textdegree}C gesunkener Körperkerntemperatur sinkt die
%     Überlebenschance des Patienten um  ganze 30\%. Eine Erhaltung der
%     Körperwärme ist daher absolut notwendig. Der Ersthelfer kann  sich
%     hierfür unterschiedlicher Mittel bedienen, eines ist die in einer
%     Autoapotheke gemäß  ÖNORM befindliche Rettungsdecke.}
%     
%     \item  psychische Betreuung:
%     
%     {\scriptsize
%     Ein sich im Stadium der Bewusstlosigkeit befindender Patient ist nicht
%     tot! Weil  ein Ersthelfer  nicht zuverlässig beurteilen kann, ob der
%     Patient in einem totalen Stadium der  Bewusstlosigkeit und somit
%     Wahrnehmungsunfähigkeit liegt, muss er den Patienten  psychologisch
%     betreuen. Diese Pflicht begründet sich nicht nur in einer moralischen
%     Komponente des zwischenmenschlichen Umgangs, sondern kann auch aus dem
%     allgegenwärtigen Anspruch auf Achtung der Menschenwürde abgeleitet
%     werden. Die  Wichtigkeit psychologischer Zuwendungen wurde ohnedies von
%     unzähligen medizinischen und  psychologischen Studien bewiesen.
%     Aufbauen kann sich die psychologische Betreuung aus  einem guten
%     Zuspruch und einer Hand die auf die Schulter des Patienten gelegt
%     wird.}
%     
%     \item  laufende Kontrolle:
%     
%     {\scriptsize
%     Der Ersthelfer soll sich alle 30 Sek. durch Überprüfung der
%     Atemtätigkeit davon überzeugen,  dass der Patientenzustand
%     unverändert geblieben ist.}
%     
% \end{enumerate}


\subsection{Notfalldiagnose Atem- und Kreislaufstillstand}
\label{topic:reanimation-eh}
\label{topic:notfalldiagnose-atem-kreislaufstillstand}

\Layout{0}{\TxBeschreibung}
{Nach der Kontrolle der Vitalfunktionen, kann der
Ersthelfer zum Ergebnis kommen, dass der Patient \EE{kein Bewusstsein} und 
\EE{keine normale Atmung} hat, d.\,h. \EE{kein Kreislauf vorhanden} ist.
In einer solchen Situation liegt die Notfalldiagnose
\T{Atem-/Kreislaufstillstand} vor.\bigskip
}{
\begin{itemize}
  \item Kein Bewusstsein
  \item Keine normale Atmung
  \item → Kein Kreislauf
\end{itemize}
}{}{}{}

\Layout{0}{Achtung}
{Entgegen der Vorgehensweise der Rettungskette wird bei einem Atem- und
Kreislaufstillstand grundsätzlich bereits \E{nach der Kontrolle der
Vitalfunktionen} der \E{Notruf abgesetzt} und erst danach mit der \E{lebensrettenden
Sofortmaßnahme} begonnen. \bigskip\bigskip
}{
\begin{itemize}
  \item Nach der Kontrolle der Vitalfunktionen Notruf absetzen.
  \item Anschließend lebensrettende Sofortmaßnahmen einleiten.
\end{itemize}
}{}{}{}

\MassnahmenNG{Erste-Hilfe-Maßnahmen}{Atem- und Kreislaufstillstand}{}{
\begin{itemize}
  \item Herz-Lungen-Wiederbelebung (Reanimation)
  \begin{itemize}
    \item Herzdruckmassage
    \item Beatmung
  \end{itemize}
\end{itemize}
}

% \Layout{0}{Lebensrettende Sofortmaßnahme
% }{
% Nach dem Absetzen des Notrufes muss sofort mit der
% \EE{Herz-Lungen-Wiederbelebung (Reanimation)} begonnen werden. Diese dient der
% Aufrechterhaltung einer von sich aus ausbleibenden Atmung und fehlenden Kreislauftätigkeit. 
% Zusätzlich zielt sie auf ein Wiederingangsetzen der eigenständigen
% Herz- sowie Lungentätigkeit ab. Die Reanimation besteht aus zwei
% Teilschritten, der \EE{Herzdruckmassage} zur Nachahmung einer
% Herztätigkeit und der \EE{Beatmung} zur Zuführung von Sauerstoff in die
% Lungen.
% }{
% \begin{itemize}
%     \item Herz-Lungen-Wiederbelebung (Reanimation)
%     \begin{itemize}
%         \item Herzdruckmassage
%         \item Beatmung
%     \end{itemize}
% \end{itemize} 
%  }

\subsubsection{Herzdruckmassage}
\index{Herzdruckmassage}

Die Herzdruckmassage bewirkt mittels Kompression des Herzens
zwischen Wirbelsäule und Brustbein ein künstliches Fließgeschehen
(Zirkulation) des Blutes.

\Layout{0}{Technik}
{
\EE{Position und Druckpunkt} Der Patient muss auf einem
harten Untergrund liegen. Ggfs. ist der Patient aus einem Bett
o.\,ä. auf den Boden zu bringen. Der Helfer kniet
dann seitlich neben dem Körper des Patienten und sucht
den richtigen \E{Druckpunkt} für die Herzdruckmassage
auf. Dieser liegt in der \EE{Mitte des Brustkorbes}.

\EE{Arm- und Körperhaltung} Die richtige
Arm- und Körperhaltung ist wichtig, um mit wenig Kraftaufwand eine
hervorragende Leistung zu erzielen. Die \E{Finger
sollten verschränkt sein}, nur der \E{Handballen liegt am Brustbein auf}.
Die \E{Arme bleiben gestreckt}, und die \E{Schultern liegen genau über
den Druckpunkt}. Nun wird mit dem Gewicht des Oberkörpers Druck auf das Brustbein ausgeübt.

\EE{Druckerzeugung} Die
Druckerzeugung hat \E{kräftig, aber nicht ruckartig} zu
erfolgen. Die \E{Eindrücktiefe} beträgt \E{\nicefrac{1}{3} des
Brustkorbes}.

\EE{Druckentlastung} Der Brustkorb muss sich wieder ganz
ausdehnen können. Es darf daher \E{kein Restdruck} bestehen. Die Druck-
und Entlastungsphasen dauern beide gleich lange.

\EE{Arbeitsgeschwindigkeit} Die Arbeitsgeschwindigkeit der
Herzdruckmassage beträgt \E{100 Massagen pro Minute}.

}{
\begin{itemize}
  \item Druckpunkt: \E{Mitte} des Brustkorbes
  \item Arm- und Körperhaltung
  \begin{itemize}
    \item Finger verschränkt,
    \item Handballen am Brustbein,
    \item Arme bleiben gestreckt,
    \item Schultern genau über den Druckpunkt.
  \end{itemize}
  \item Druckerzeugung
  \begin{itemize}
    \item kräftig aber nicht ruckartig
    \item Eindrucktiefe: \nicefrac{1}{3} des Brustkorbes
  \end{itemize}
  \item Druckentlastung
  \item 100\PerMin*
\end{itemize}
}{}{}{}

\Layout{0}{Mögliche Fehlerquellen}
{Das Schlimmste, das der Ersthelfer falsch machen kann, ist
 nichts zu tun. Wenn reanimiert wird, kann es zu dem
 einen oder anderen Fehler kommen. Der Patient muss auf
 einer festen Unterlage gelagert werden, ein Bett ist
 \EE{keine feste Unterlage}! Es darf nicht vergessen werden, den Patienten auf den Boden oder eine andere harte Unterlage zu
 legen.

Grundsätzlich
kann es auch zu \EE{Verletzungen durch den Ersthelfer} kommen. Dazu zählt der
Bruch des Brustbeines durch einen zu festen Druck, Lungen- oder Herzquetschung
durch Abweichung des Druckpunktes, Rippenbrüche durch Abweichen des
Druckpunktes oder nicht senkrecht ausgeführten Druck, Organquetschung bzw.
-einreißung durch einen zu tiefen Druckpunkt.  

Ein zu ruckartiger Druck führt zu einer zu kurzen Kompression des Herzens, was
wiederum eine verkürzte Auswurfphase bedeutet. Dies alles führt zu einer
\EE{ungenügenden Zirkulation} des Blutkreislaufs. 

\Cave{Wird einmal falsch Druck ausgeübt, oder eine Rippe gebrochen, nicht 
 aufhören, sondern den
richtigen Druckpunkt erneut aufsuchen und weiter machen. 

Gebrochene Rippen und
gequetschte Organe können wieder repariert werden. Das Herz, das nicht mehr
schlägt, benötigt unsere Hilfe!}
}{
\begin{itemize}
  \item keine feste Unterlage
  \item Verletzungen durch den Ersthelfer
  \item ungenügende Zirkulation
\end{itemize}
}{}{}{}

\subsubsection{Beatmung}
\index{Beatmung!Erste-Hilfe}

\Layout{0}{Arten}
{
Die bekannteste Methode zur Beatmung durch den Ersthelfer ist die
\T{Mund--zu--Mund--Beatmung}. Dabei wird bei überstrecktem Kopf mit der stirnseitigen Hand die Nase des Patienten zugedrückt. Anschließend mit dem eigenen Mund den des Patienten umschließen und
dabei die Lippen luftdicht an den Patienten drücken. Nun
wird die eigene Ausatemluft in den Patient eingeblasen. Vor der nächsten Atemspende
muss der Mund wieder abgesetzt werden, um selbst frische Luft einatmen zu
können. 

Die zweite Möglichkeit ist die \T{Mund--zu--Nase--Beatmung}. Hier wird bei überstrecktem Kopf mit der kinnseitigen Hand der Mund des Patienten
zugehalten. Anschließend mit dem eigenem Mund die Nase des Patienten umschließen und die Lippen luftdicht an den Patienten drücken. Nun wird die eigene
Ausatemluft in den Patienten eingeblasen. Vor der nächsten Atemspende
muss der Mund wieder abgesetzt werden, um selbst frische Luft einatmen zu
können. 

Bei Säuglingen und Kleinkindern wäre noch die
\T{Mund--zu--Mund/Nasen-Beatmung} möglich. In diesem
besonderen Fall wird der Mund und die Nase des Patient mit dem eigenem Mund umschlossen.
}{
\begin{itemize}
  \item Mund--zu--Mund-Beatmung
  \item Mund--zu--Nasen-Beatmung
  \item Mund--zu--Mund/Nasen-Beatmung
\end{itemize}
}{}{}{}

\Layout{0}{Wichtige Parameter für die Beatmung}
{
Die \EE{Atemfrequenz} des Erwachsenen beträgt
\VonBis{12}{16}\PerMin*, d.\,h. wir atmen
\VonBis{12}{16}\PerMin* ein und aus. Das
\EE{\underline{max.} (!) Atemzugsvolumen} in ml entspricht
dem 10-fachen des Körpergewichtes, d.\,h. ein 90\Kg*
schwerer Mann kann 900\Ml* Luft ein- und wieder ausatmen.
Ein normaler Atemzug beträgt jedoch nur etwa 500\Ml* (vgl.
\See{topic:atemvolumina}).

Die ausreichende \ce{O2}-Konzentration in der Umgebungs-/Einatmungsluft ist die
Voraussetzung für eine funktionierende Atmung. Ist sie
nicht gegeben, ist auch die beste Atemarbeit nutzlos. Der normale \ce{O2}-Gehalt der Umgebungsluft beträgt
21\PC*. Der \ce{O2}-Gehalt in der \EE{Ausatemluft} beträgt
noch 16\PC*, d.\,h. der Mensch verbraucht  nur 5\PC* des
eingeatmeten Sauerstoffs. 
}{
\begin{itemize}
  \item Atemfrequenz \VonBis{12}{16}\PerMin*
  \item Max. Atemzugsvolumen [ml]: 10-fache des Körpergewichtes
  \item Einatemluft 21\PC* \ce{O2}-Gehalt
  \item Ausatemluft 16\PC* \ce{O2}-Gehalt
\end{itemize}
}{}{}{}

\Layout{0}{Bei der Beatmung ist zu beachten:}
{
Ein \EE{zu hohes Atemzugsvolumen} führt zu einem
höheren Druck im Nasen-Rachenraum, sobald die Lunge
vollständig mit Luft befüllt ist. Die Luft findet nur den
Ausweg über die Speiseröhre in den Magen. Dies kann während einer Herz-Lungen-Wiederbelebung meist unbemerkt geschehen. Bei einem \EE{zu niedrigem Atemzugsvolumen} ist die
\ce{O2}-Versorgung nicht sichergestellt, da die Luft nicht bis in die
Lunge kommt. Bei Bartträgern kann es zu einer
\EE{ungenügenden Abdichtung} kommen. Eine \EE{zu hohe
Atemfrequenz} hindert den Patienten am Ausatmen. Wenn es beim Freimachen
der Atemwege zu einer \EE{mangelnden Überstreckung} des Kopfes
kommt, kann es zu einer unzureichenden Belüftung der Lunge kommen. Bei jeder Atemspende ist der \EE{Brustkorb} des Patienten zu \EE{beobachten}. Dieser
muss sich deutlich heben und senken. Immer auch auf den \EE{Selbstschutz}
achten. Wenn möglich nur mit Beatmungstuch oder Beatmungsmaske beatmen.
}{
\begin{itemize}
  \item zu hohes Atemzugsvolumen ${\rightarrow}$ Luft in den Magen
  \item zu niedriges Atemzugsvolumen ${\rightarrow}$ \ce{O2}-Versorgung nicht
  sichergestellt
  \item ungenügende Abdichtung bei Bartträgern
  \item mangelnde Überstreckung ${\rightarrow}$ unzureichenden Belüftung der
  Lunge
  \item Brustkorb beobachten ${\rightarrow}$ deutlich heben und senken
  \item Selbstschutz ${\rightarrow}$ mit Beatmungstuch oder Beatmungsmaske
  beatmen
\end{itemize}
}{}{}{}

\subsubsection{Algorithmus Herz-Lungen-Wiederbelebung}

\Layout{0}{Vorgehensweise}{
Nach der Kontrolle des Bewusstseins und der Feststellung, dass der Patient nicht
ansprechbar ist, wird um Hilfe gerufen. Anschließend durch
Überstrecken des Kopfes die Atemwege öffnen und die Atmung kontrollieren. Wenn keine normale
Atmung vorhanden ist, Notruf absetzen. Danach mit der Herz-Lungen
Wiederbelebung beginnen. Diese setzt sich aus 30 Herzdruckmassagen und 2
Beatmungen zusammen. 
}{
\prettyref{tab:bilderserie-hlw} \ldots{} Bilderserie. 

\prettyref{tab:algorithmus-bls-aed} \ldots{} Alg. AED.
}{}{}{}

\Layout{57}{}
{}{}
{
\begin{tabularx}{\linewidth}{XXX}
\includegraphics[width=\linewidth]{./images/koch-aass/00800-gamma/Rea1b.JPG}
&
\includegraphics[width=\linewidth]{./images/koch-aass/00800-gamma/Rea2b.JPG}
&
\includegraphics[width=\linewidth]{./images/koch-aass/00800-gamma/Rea3b.JPG}
\\
% Schritt 1 
Ansprechen, berühren, Schmerzreiz
&
% Schritt 2 
Hilferuf
&
% Schritt 3 
Atemwege freimachen
\\\addlinespace
\includegraphics[width=\linewidth]{./images/koch-aass/00800-gamma/Rea4b.JPG}
&
\centering\Huge{\ttfamily 144}
&
\includegraphics[width=\linewidth]{./images/koch-aass/00800-gamma/Rea5b.JPG}
\\
% Schritt 4 
Atemkontrolle
& 
Notruf
&
% Schritt 5
30 Herzdruckmassagen
\\\addlinespace
\includegraphics[width=\linewidth]{./images/koch-aass/00800-gamma/Rea6b.JPG}
&
\includegraphics[width=\linewidth]{./images/koch-aass/00800-gamma/Rea7b.JPG}
&
\\
% Schritt 6
2 Beatmungen
&
% Schritt 7
Wenn vorhanden: Defibrillation.\ACh{Instruktor 1}{0.8.0}{Weg lassen.
}\ACh{GABS}{2010-02-19}{Wieso? Hinweis auf AED wäre doch sinnvoll.} 
&
\\\addlinespace\bottomrule
\end{tabularx}
}{Bilderserie: Herz-Lungen-Wiederbelebung
(Reanimation)\label{tab:bilderserie-hlw}\index{Bilderserie!Reanimation}\index{Reanimation!Bilderserie}}{}

% {
% {Die Durchführung der Reanimation hält sich an folgenden
% Algorithmus:}\footnote{\ {Die folgende Darstellung der
% Reanimation erlaubt nur ein Überblick über die Möglichkeiten und
% Maßnahmen der Reanimation. Eine ausführlichere Darstellung findet
% sich in den nachfolgenden Kapiteln des Lehrbuches. }}}
% 
% 
% 
% 
%  Einfügen des
% BLS-Algorithmus erstellt vom Josef 2006 für die
% ASB-921-Schulung


\Layout{0}{Ende der HLW}
{\index{Reanimation!Beendigung}
Beendet wird eine Herz-Lungenwiederbelebung durch den Ersthelfer nur
dann, wenn einer der folgenden drei Umstände eintritt. Ohne diese
zusätzlichen Umstände darf der Ersthelfer die Reanimation nicht
unterbrechen: 

\begin{description}
  \item[Eintreffen von \EE{kompetenter Hilfe.}] Der Ersthelfer darf aber beim
  Eintreffen der Rettungskräfte nicht sofort mit der Wiederbelebung  aufhören, sondern muss diese weiterführen
  bis er ausdrücklich vom Einsatzpersonal abgelöst  wird.
  \item[Erschöpfung.] Sollte der Ersthelfer all seine Kräfte aufgebraucht haben,
  wird gezwungener Maßen die Wiederbelebung ab-, bzw. unterbrochen.
  % aufgrund der
  % körperlichen Anstrengung im Zusammenhang mit der Reanimation all seine Kräfte aufbrauchen, und daher fortführen können, ist es  ihm erlaubt diese
  % einzustellen. Diese Möglichkeit wird aus den Grundsätzen der Zumutbarkeit
  % und der Versuchspflicht gewonnen und darf nur als absolut
  % letzte Option betrachtet werden.
  \item[Der Patient zeigt \EE{Lebenszeichen}.] Entweder
  erwacht der Patient und ist ansprechbar \footnote{eine
  leider äußerst seltene Form der Reanimationsbeendigung},
  oder der Patient gibt ein Husten oder Räuspern von sich
  bzw. atmet beim Beatmungsvorgang dem Ersthelfer entgegen. \EE{Erbrechen zählt \E{nicht} als Lebenszeichen.}
  % Da hier der Ersthelfer nicht mit verlässlicher Sicherheit vom Erfolg seiner  Wiederbelebung
  % ausgehen kann, muss er erneut die Atmung kontrollieren. Sollte sich das
  % Lebenszeichen durch vorhandene Atmung bestätigen, muss der Patient
  % in die stabile  Seitenlagerung gebracht und in Form der unten
  % beschriebenen anschließenden Behandlung  weiterversorgt werden.\item[]
\end{description}

}{
\begin{itemize}
  \item Eintreffen von kompetenter Hilfe
  \item Erschöpfung
  \item Patient zeigt Lebenszeichen
\end{itemize}
}{}{}{}


 

% \begin{enumerate}
%     \item  Lebenszeichen:
%     
%     {\itshape
%     Lebenszeichen können in zwei Unterkategorien auftreten:}
%     
%     \begin{enumerate}
%         \item  sichere Lebenszeichen:
%         
%         {\scriptsize
%         Der Patient erwacht und ist ansprechbar. Eine leider äußerst
%         seltene Form der  Reanimationsbeendigung.}
%         
%         \item  unsichere Lebenszeichen:
%         
%         {\scriptsize
%         Der Patient gibt ein Husten oder Räuspern von sich bzw. atmet beim
%         Beatmungsvorgang dem  Ersthelfer entgegen. Da hier der Ersthelfer nicht
%         mit verlässlicher Sicherheit vom Erfolg seiner  Wiederbelebung
%         ausgehen kann, muss er erneut die Atmung kontrollieren. Sollte sich das
%         Lebenszeichen durch vorhandene Atmung bestätigen, muss der Patient
%         in die stabile  Seitenlagerung gebracht und in Form der unten
%         beschriebenen anschließenden Behandlung  weiterversorgt werden.}
%     \end{enumerate}
%     \item  Eintreffen der Rettung:
%     
%     {\scriptsize
%     Der Ersthelfer darf bei Eintreffen der Rettungskräfte nicht sofort
%     mit der Wiederbelebung  aufhören, sondern muss diese weiterführen
%     bis er ausdrücklich vom Einsatzpersonal abgelöst  wird.}
%     
%     \item  Aufbrauchen der eigenen Kräfte:
%     
%     {\scriptsize
%     Sollte der Ersthelfer aufgrund der körperlichen Anstrengung im
%     Zusammenhang mit der  Reanimation alle seine Kräfte aufbrauchen,
%     sodass er sie nicht weiter fortführen kann, ist es  ihm erlaubt diese
%     einzustellen. Diese Möglichkeit wird aus den Grundsätzen der
%     Zumutbarkeit  und der Versuchspflicht gewonnen und darf nur als absolut
%     letzte Option verwendet werden.}
%     
% \end{enumerate}


% \subsubsection{8.1) Exkurs -- Defibrillation:}
% 
% {
% Auch für den Ersthelfer ist die Anwendung eines Defibrillators
% möglich. Die Durchführung orientiert sich hierbei an einem
% ähnlichen Ablauf wie beim Rettungssanitäter, wodurch an dieser
% Stelle nicht näher auf die Anwendung eingegangen wird. Zur näheren
% Auseinandersetzung siehe das Kapitel
% {\quotedblbase}{\dots}{\textquotedblleft}.}

\clearpage % TODO temp clearpage
\Layout{0}{Anschließende Behandlung}
{%
Die anschließenden Behandlungen orientieren sich an den
verschiedenen Beendigungsmöglichkeiten der
Reanimation, wodurch auch sie in ihrer Ausgestaltung verschieden sind. Zeigt
der Patient \EE{Lebenszeichen}, ist eine
Atemkontrolle durchzuführen und wenn eine normale Atmung
vorhanden ist, wird der Patient in die stabile Seitenlage umgelagert.
Anschließend weiterführende Maßnahmen gemäß Notfalldiagnose Bewusstlosigkeit. 

Bei Eintreffen von \EE{kompetenter Hilfe} muss der Ersthelfer
dem Rettungspersonal Auskunft über relevante Umstände und
die gesetzten Maßnahmen geben. 

Wenn der \EE{Ersthelfer erschöpft} ist, sollte er (falls überhaupt
konditionell möglich) erneut versuchen andere  Ersthelfer zu
finden. Scheitert auch dieser Versuch, kann der Ersthelfer keine weitere 
Maßnahme setzen als den Patienten zuzudecken und auf das
Eintreffen des  Rettungspersonals zu warten. Der Ersthelfer hat alles
in seiner Macht stehende getan, und es ist ihm kein Vorwurf
zu machen. 
}{
\begin{itemize}
  \item Lebenszeichen kontrollieren
  \item Atemkontrolle alle 30 Sek
  \item Stabile Seitenlage.
  \item Wärmeerhalt (zudecken)
  \item Frischluftzufuhr (Fensteröffnen, \ldots)
  \item Psychische Betreuung
  \item Kompetente Hilfe kommt ${\rightarrow}$ Auskunft über gesetzte
  Maßnahmen geben
  \item Erschöpfung: ${\rightarrow}$ versuchen weitere Ersthelfer zu
  finden
\end{itemize}

}{}{}{}




% \subsubsection{8.2) anschließende Behandlung:}
% 
% {
% Die anschließenden Behandlungen orientieren sich an den
% verschiedenen Beendigungsmöglichkeiten einer ersthelferischen
% Reanimation, wodurch auch sie in ihrer Ausgestaltung verschieden sind.}
% 
% {
% \foreignlanguage{english}{{
% a}}\foreignlanguage{english}{{d i) Lebenszeichen:}}}

% {
%  ad i) a) sichere Lebenszeichen:}
% 
% {\scriptsize
%  Der sich bei Bewusstsein befindende Patient ist in eine für ihn
% angenehme Position zu  bringen. Er ist psychisch zu betreuen und
% laufend zu beobachten. Um einen Wärmeverlust zu  vermeiden muss der
% Patient zugedeckt werden.}

% {
%  ad i) b) unsichere Lebenszeichen:}
% 
% {\scriptsize
%  Nach Umlagerung in die stabile Seitenlagerung sind die
% {\quotedblbase}ersthelferischen Allzweckwaffen{\textquotedblleft} 
% (Wärmeerhalt sichern, psychische Betreuung, laufende Kontrolle)
% durchzuführen.}

% {
%  ad ii) Eintreffen der Rettung:}
% 
% {\scriptsize
%  Der Ersthelfer muss dem Rettungspersonal Auskunft über relevante
% Umstände und gesetzte  ersthelferische Maßnahmen geben.}
% 
% {
%  ad iii) Aufbrauchen der eigenen Kräfte:}
% 
% {\scriptsize
%  Sollte der Ersthelfer sich nicht mehr körperlich in der Lage fühlen
% eine Reanimation  fortzuführen, sollte er (falls überhaupt
% konditionell möglich) erneut versuchen weitere  Ersthelfer zu finden.
% Scheiterte auch dieser Versuch, kann der Ersthelfer keine weiteren 
% Maßnahmen setzen als den Patienten zuzudecken und auf das
% Eintreffen des  Rettungspersonals zu warten. Der Ersthelfer hat alles
% in seiner Macht stehende getan und  darf sich keine Vorwürfe eines
% Versagens machen.}




% \subsection{9) Atemwegsverlegung durch Fremdkörper:}
% 
% {
% Unter dem Begriff {\quotedblbase}Bolus-Geschehen{\textquotedblleft}
% werden Fälle erfasst, bei denen das durch einen Biss (Bolus) in den
% Mundraum gelangte Objekt in die falsche Röhre, nämlich die
% Luftröhre, wandert. Bei der Behandlung muss unterschieden werden
% zwischen einer teilweisen, leichten Verlegung oder einer gänzlichen,
% schweren Verlegung.}
% 
% \subparagraph{Leichte Verlegung:}
% 
% \begin{enumerate}
%     \item  Kennzeichen:
%     \begin{enumerate}
%         \item  effektives Husten und Würgen
%     \end{enumerate}
%     \item  ersthelferische Behandlung:
%     \begin{enumerate}
%         \item  Patient beim Husten unterstützen durch kreisförmiges Reiben am
%         Rücken (andere Manipulationen sind zu unterlassen!)
%         \item Notrufabgabe
%         \item  Wärmeerhalt sichern, psychische Betreuung, laufende Kontrolle
%     \end{enumerate}
% \end{enumerate}
% 
% 
% \subparagraph{
% Schwere Verlegung:}
% Kennzeichen:
% \begin{enumerate}
%     \item  ineffektives (stummes) Husten
%     \item   Aussetzen der Atmung
%     \item   nach gewisser Zeit Bewusstlosigkeit
% 
%     \item  ersthelferische Behandlung:
% \end{enumerate}
% \begin{enumerate}
%     \item  Rückenschläge: 5x zwischen die Schulterblätter
% tangential schlagen,    danach Kontrollieren ob Objekt abgehustet
% wurde\footnote{\ {Für den Sanitäter besteht abwechselnd zu
% den Rückenschlägen, die Pflicht zur Durchführung eines
% Heimlich-Manövers. Die nähere Durchführung siehe im Kapitel
% {\quotedblbase}{\dots}{\textquotedblleft}.}}
% 
%     \item  wurde das Objekt nicht abgehustet und ist der Patient weiterhin bei
%    Bewusstsein: erneute 5 Rückenschläge mit anschließender
% Kontrolle
% 
%     \item  wurde das Objekt nicht abgehustet und ist der Patient bewusstlos
% (und   zwingend ohne Atmung und Kreislauf):    
% 
% {
%   Notrufabgabe + Wiederbelebung (auch hier ist die Beatmung zwingend   
% durchzuführen; es besteht die Chance durch die Beatmung das Objekt in
% den   rechten Bronchus zu bewegen)}
% 
%     \item  wurde das Objekt abgehustet und ist der Patient bei Bewusstsein: 
% 
% {
%   Lagerung mit erhöhten Oberkörper, Notrufabgabe, Wärmeerhalt
% sichern,   psychische Betreuung, laufende Kontrolle}
% 
%     \item wurde das Objekt abgehustet und ist der Patient bewusstlos:
% 
% {
%   bei aufrechter Atmung/Kreislauf: stabile Seitenlagerung, Notrufabgabe,
%    Wärmeerhalt sichern, psychische Betreuung, laufende Kontrolle}
% 
% {
%   bei fehlen von Atmung/Kreislauf: Notrufabgabe + Wiederbelebung}
% \end{enumerate}
% {
% Selbst bei relativ schneller Abhustung des Fremdkörpers, ist es ratsam
% eine ärztliche Kontrolle aufzusuchen, da es im Zuge eines
% Bolus-Geschehn leicht zu {\quotedblbase}hypoxischen
% Schäden{\textquotedblleft} (Schäden durch Sauerstoffmangel) kommen
% kann.}

\subsection{Starke Blutung}
\label{topic:starke-blutung}

% Bei staren, lebensbedrohlichen Blutung verliert der Körper innerhalb kurzer
% Zeit viel Blut. Diese kann entweder spritzend bei arteriellen Blutungen oder
% quellend bei venösen Blutungen sein. Als \EE{lebensrettende Sofortmassnahme}
% dient die \EE{Blutstillung}. 
% Vor Durchführung der Blutstillung muss sich der Ersthelfer selbst
% schützen indem er sich medizinische Einweghandschuhe anzieht, um der
% Gefahr einer Ansteckung mittels von Blut übertragender Krankheiten
% vorzubeugen. Bei der Durchführung der Blutstillung sind folgende
% Methoden möglich, wobei leider nicht alle immer sinnvoll sind.

\MassnahmenNG{Erste-Hilfe-Maßnahmen}{Starke Blutungen}{}{
\begin{itemize}
  \item Blutstillung
  \begin{itemize}
    \item Zudrücken,
    \item Abdrücken,
    \item Druckverband,
    \item Abbindung.
  \end{itemize}
\end{itemize}

}

\Layout{22}{}{\index{Blutstillung}
Für die Blutstillung werden folgende Materialien
benötigt: Handschuhe, sterile Wundauflagen (z.\,B.
Zetuvit\texttrademark{} Saugkompresse) in verschiedenen
Größen, Dreiecktücher und Mullbinden. Oft sind auch Verbandspäckchen vorhanden. Hier
ist eine Wundauflage bereits in die Mullbinde eingearbeitet.
}{}{./images/koch-aass/00800-gamma/MaterialDruckverbandb.JPG}{}{}

\Layout{60}{Blutstillung durch Zudrücken}
{Mit einem \EE{direkten Druck in die stark blutende Wunde} mittels Finger,
Handballen oder Faust -- je nach Größe der Verletzung -- mit einer Wundauflage 
kann diese gestoppt werden. Ein \EE{Hochhalten der verletzten Extremität} verstärkt
die Wirkung. Dies ist auch die \EE{einzige Möglichkeit Halsschlagaderblutungen zu
stoppen}. Nachteil an dieser Methode ist, dass der Ersthelfer ab nun nicht mehr
loslassen darf und somit gehandicapt ist.
}{
\begin{itemize}
  \item Direkter Druck in die Wunde
  \item Hochhalten der verletzten Extremität
  \item Einzige Möglichkeit Halsschlagaderblutungen zu stoppen
\end{itemize}
}
{./images/hauer-david-aass/blutstillung_druck-00800.jpg}
{Simpel, aber Wirkungsvoll: Draufdrücken!}
{Hauer}

\Layout{0}{Blutstillung durch Abdrücken
}{
Der Ersthelfer drückt in diesem Fall die blutzuführende Arterie herzwärts der
stark blutenden Wunde ab. Die Durchführung ist an
verschiedenen Stellen möglich. Bei Blutungen am Schädeldach
kann jeweils an der betroffenen Seite an der Schläfe abgedrückt werden. Bei
Blutungen im Bereich des Gesichts wird je nach Seite
am Kinn abgedrückt. Bei Blutungen an Schulter oder Oberarm liegt die
Abdrückstelle in der jeweiligen Schlüsselbeingrube. Unterarmblutungen werden an
der betroffenen Oberarminnenseite abgedrückt. Blutungen im Bereich des
Oberschenkels sind in der Leistenbeuge und Unterschenkelblutungen an der
Oberschenkelinnenseite abzudrücken.

Nachteil hierbei ist nicht nur dasselbe Handicap wie
bei direkten Druck in die Wunde, sondern auch die Tatsache, dass  die gesamte
betroffene Stelle ohne Blutversorgung bleibt. Außerdem ist das Auffinden der
Druckstellen für ungeübte Ersthelfer nicht immer möglich.
%%%%%%%%%%%%%%%%%%%%%%%%%%%%%%%%%%%%%%%%%%%%%%%%%%%%%%%%%%%
% TODO : Neuerstellung Abdrückpunkte
\TODO{GABS}{yyyy-mm-dd}{Abdrückpunkte: Neuerstellung von Grafik/Text 
                          notwendig}{}
%%%%%%%%%%%%%%%%%%%%%%%%%%%%%%%%%%%%%%%%%%%%%%%%%%%%%%%%%%%
}{
\begin{itemize}
  \item Abdrücken der blutzuführenden Arterie
\end{itemize}
}{}{}{}

\Layout{0}{Druckverband
}{\index{Druckverband}
Der Druckverband ist die beste Möglichkeit zur Stillung einer Blutung, da sie
in fast allen Situationen funktioniert und einfach und
schnell umsetzbar ist.

Eine Ausnahme der Anwendung ist die Blutung aus der Halsschlagader. Hierbei
besteht nur die Möglichkeit des Zudrückens. Außerdem kann der Druckverband nur
an den Gliedmaßen angelegt werden. Beim Druckverband wird nur Druck auf das
verletzte Gefäß ausgeübt. Somit ist die Durchblutung der Extremitäten unterhalb
der Wunde gewährleistet. 

Ein Druckverband besteht aus drei Teilen. Einer
\EE{sterilen Wundauflage}, einem \EE{Druckkörper} \footnote{Druckkörper: Als Druckkörper
eignen sich nur kompakte, weiche und saugfähige Objekte. Ein solches
Objekt, das sich auch in der Verbandsbox befindet, ist die Mullbinde.} und einem
\EE{Befestigungsmittel} \footnote{Befestigungsmittel: Als Befestigungsmittel eignet sich eine
Dreiecktuchkrawatte oder eine Mullbinde}. Ein \EE{Hochhalten der verletzten
Extremität} verstärkt die Wirkung.
}{
\begin{itemize}
  \item 3 Teile:
  \begin{itemize}
    \item Sterile Wundauflage,
    \item Druckkörper,
    \item Befestigungsmittel.
  \end{itemize}
  \item Hochhalten der verletzten Extremität
\end{itemize}
}{}{}{}

\Layout{57}{}
{}{}
{
\begin{tabularx}{\linewidth}{XXX}
\includegraphics[width=\linewidth]{./images/koch-aass/00800-gamma/Abdrueckenb.JPG}
&
\includegraphics[width=\linewidth]{./images/koch-aass/00800-gamma/Druckverband1b.JPG}
&
\includegraphics[width=\linewidth]{./images/koch-aass/00800-gamma/Druckverband2b.JPG}
\\
Schritt 1 &
Schritt 2 &
Schritt 3 
\\\addlinespace\bottomrule
\end{tabularx}
}{Bilderserie: Zudrücken und Druckverband}{}

\Layout{0}{Abbindung
}{\index{Abbindung}
Die Abbindung ist weitestgehend aus der Ersten Hilfe eliminiert worden. Grund
dafür ist, dass sich das Verhältnis zwischen Nutzen und Folgeschäden
massiv zu den Folgeschäden verschoben hat. Außerdem hat
sich die rettungsdienstliche Vollversorgung inklusive rascher Eintreffzeiten und die 
Frühversorgung in unfallchirugischen Fachabteilungen deutlich verbessert. 

\Cave{Abbindung nur bei lebensbedrohlichen Blutungen, die nicht anders
stillbar sind!\footnote{Z.B. bei totaler Amputation, teilweiser Abtrennung/Zerfetzung oder Einklemmung.}

\Ca{Die Abbindung ist die letzte Wahl!}}
 
Das Abbinden \EE{funktioniert} aus anatomischen Gründen
\EE{nur am Oberarm und Oberschenkel}, denn dort befindet
sich nur ein Knochen, gegen den sämtliche Blutgefäße
gedrückt werden können. Dazu darf nur ein \EE{breites und
schonendes Material} verwendet werden (z.\,B. Dreiecktuchkrawatte,
Blutdruckmanschette, etc.). Nachteil ist die nicht mehr gegebenen Blutversorgung
der Extremität. Für die weitere Versorgung des Patienten ist es wichtig, den \EE{Zeitpunkt der Abbindung} zu notieren, am besten gleich am Oberarm bzw.
Oberschenkel des Patienten.\footnote{Manche Autoren
empfehlen, dass die maximale Abbindezeit ca. 30 Minuten
betragen sollte. Nach heutigem Stand der Wissenschaft
wird empfohlen, einmal
angebrachte Abbindungen bis zur definitiven Versorgung
\E{nicht} (auch nicht probeweise) zu öffnen.
\cite{PHTLS6,BuchfelderErsteHilfe4}}
}{ 
\begin{itemize}
  \item Nur am Oberarm und Oberschenkel
  \item Nur breites und schonendes Material
  \item Zeitpunkt der Abbindung notieren
\end{itemize}
}{}{}{}

\Layout{57}{}
{}{}
{
\begin{tabularx}{\linewidth}{XXX}
\includegraphics[width=\linewidth]{./images/koch-aass/00800-gamma/AbbindungArm1b.JPG}
&
\includegraphics[width=\linewidth]{./images/koch-aass/00800-gamma/AbbindungArm2b.JPG}
&
\\
% Schritt 1 
Abbinden am Oberarm
&
% Schritt 2 
Zeit notieren
&
\\\addlinespace\bottomrule
\end{tabularx}
}{Bilderserie: Abbinden am Oberarm}{}

\Layout{57}{}
{}{}
{
\begin{tabularx}{\linewidth}{XXX}
\includegraphics[width=\linewidth]{./images/koch-aass/00800-gamma/AbbindungBein1b.JPG}
&
\includegraphics[width=\linewidth]{./images/koch-aass/00800-gamma/AbbindungBein2b.JPG}
&
\includegraphics[width=\linewidth]{./images/koch-aass/00800-gamma/AbbindungBein3b.JPG}
\\
Schritt 1 &
Schritt 2 &
Schritt 3 
\\\addlinespace
\includegraphics[width=\linewidth]{./images/koch-aass/00800-gamma/AbbindungBein4b.JPG}
&
&
\\
Schritt 4 &
&
\\\addlinespace\bottomrule
\end{tabularx}
}{Bilderserie: Abbinden am Oberschenkel}{}

\Layout{0}{Versorgung abgetrennter Gliedmaßen}
{
Die abgetrennten Gliedmaßen sind \EE{immer in das Spital mitzunehmen}. Nach
Versorgung des Patienten ist das \EE{Amputat steril} zu \EE{verbinden}. Für
den Transport gilt es \EE{kühl und trocken} zu \EE{lagern}. Dazu eignen
sich idealerweise Kühlakkus für Kühltaschen falls vorhanden. Wichtig ist aber,
dass das Amputat  nicht direkt mit Eis oder dem Kühlakku in Berührung
kommt, da es sonst zum Gefrieren und Absterben von Gewebe kommt.
}{
\begin{itemize}
  \item immer ins Spital mitnehmen
  \item Amputat ist steril zu verbinden
  \item für den Transport kühl und trocken lagern
\end{itemize}
}{}{}{}





\subsection{Schock}
\index{Schock!Erste-Hilfe}
\label{topic:schock-ehi}

\Definition{{\quotedblbase}Der Schock ist eine lebensbedrohliche
Kreislaufstörung, bei der es zu einer Unterversorgung der
Organe mit Sauerstoff kommt.{\textquotedblleft}}

\emph{Vgl. \See{topic:schock}.}

Das Krankheitsbild, welches umgangssprachlich als Schock bezeichnet wird, ist
eine Schreckreaktion mit Muskelzittern, weichen Knien, Schwächegefühl und
Blässe. Es hat aber mit dem echten Schock nichts zu tun, sondern ist nur eine
Reaktion des Nervensystems und normalerweise nicht bedrohlich. Beruhigen und
hinlegen lassen führt bei den Patienten zu einer schnellen Erholung.

\Layout{0}{\TxBeschreibung{} und Verlauf}
{
Der Schock ist eine \EE{lebensbedrohliche Kreislaufstörung}, bei der es zu einer
\EE{Unterversorgung der lebenswichtigen Organe}, wie Hirn,
Herz, Lunge, Leber und Niere mit Sauerstoff kommt. 

\subparagraph{Zentralisation} Der Körper hilft sich
selbst indem er das Blut zu den lebenswichtigen Organen bringt und äußere Körpergebiete wie Arme,
Beine und Haut weniger bis gar nicht durchblutet. 

\subparagraph{Dezentralisation} Hält
der Zustand länger an, bzw. ist nach wie vor zu wenig Blut vorhanden, beginnt
der Körper wieder zu dezentralisieren und lässt das Blut in weniger
wichtige Körperregionen fließen. Der dadurch entstehende Sauerstoffmangel
bei den lebenswichtigen Organen führt zu einer \EE{Schädigung der Organe}, die
in weiterer Folge zum \EE{Versagen der Organe} führt. Der Schock
ist daher unbehandelt in seiner finalen Wirkung tödlich!
}{
\begin{itemize}
  \item Lebensbedrohliche Kreislaufstörung
  \item Unterversorgung lebenswichtiger Organe
  \begin{itemize}
    \item Hirn, Herz, Lunge, Leber, Nieren
  \end{itemize}
  \item Zentralisation
  \item Dezentralisation
  \item Schädigung der Organe, Versagen der Organe
  \item Unbehandelt führt der Schock zum Tod!
\end{itemize}
}{}{}{}


% {Der Schock ist eine lebensbedrohliche Kreislaufstörung, bei der es zu
% einer Unterversorgung der lebenswichtigen Organe wie Herz, Hirn, Lunge,
% Leber, Niere mit Sauerstoff kommt. Der Sauerstoffmangel führt zu
% einer Schädigung des Organgewebes, der im ersten Schritt zum Versagen
% dieser Organe und im zweiten Schritt zum Versagen des gesamten
% Organismus führt. Der Schock ist daher in seiner finalen Wirkung
% tödlich!}


\Layout{0}{Ursachen}
{
Es gibt drei mögliche Ursachen, welche zum Schock führen
können:

\EE{Blut- und Flüssigkeitsverlust}. Der Blutverlust kommt von Verletzungen mit
starken Blutungen. Flüssigkeitsverlust kann bei großflächigen Verbrennungen und
bei extremen Brechdurchfall vorkommen. 

Bei \EE{Herzversagen} ist zwar
genügend Blut vorhanden, aber die Pumpe -- das Herz -- ist
zu schwach um es durch den Körper zu pumpen. 

Bei \EE{Vergiftungen und allergischen Reaktionen} bleibt
zwar auch die Blutmenge konstant, aber die Gefäße erweitern sich. Daher
entsteht ein unzureichender Druck in den Gefäßen, und das Blut kann nicht weiter
transportiert werden. 
}{
\begin{itemize}
  \item Blut- und Flüssigkeitsverlust
  \item Herzversagen
  \item Vergiftungen und allergische Reaktionen
\end{itemize}
}{}{}{}

% \subsubsection{10.1) Schockursachen:}
% 
% {
% Die Ursachen des Schocks sind:}
% 
% {
%  i) Blut- und Flüssigkeitsverlust:}
% 
% {
% {\textit{ Im Idealzustand sind alle Gefäße im Körper bis
% zum Rand mit Blut gefüllt. Durch diese 
% {\quotedblbase}Vollfüllung{\textquotedblleft} wird ein Druck in den
% Gefäßen aufgebaut, der gewährleistet, dass das Blut in alle 
% Körperregionen transportiert werden kann. Verliert der Körper
% massiv Blut und kann er  dadurch den Druck in den Gefäßen nicht
% mehr aufrechterhalten, geht ebenso die Fähigkeit  verloren den
% gesamten Körper ausreichend mit dem verbleibenden Blutvolumen
% zu}}{\textit{  }}{\textit{versorgen. Die Organe werden
% unterversorgt, dadurch geschädigt und schließlich absterben.}}}
% 
% {
%  ii) Vergiftungen und allergische Reaktionen:}
% 
% {\itshape
%  Bleibt zwar das Blutvolumen konstant, die Gefäßwände erweitern
% sich aber in Folge einer  Vergiftung oder allergischen Reaktion,
% entsteht derselbe unzureichende Druck in den  Gefäßen. Die Organe
% werden unterversorgt, dadurch geschädigt und schließlich
% absterben.}
% 
% {
%  iii) Herzversagen:}
% 
% {\itshape
%  Selbst bei besten Druckverhältnissen in den Gefäßen kann es zu
% einer Unterversorgung  kommen, wenn der Motor des Blutkreislaufes
% ausfällt. Pumpt folglich das Herz zu wenig Blut  in die Arterien,
% kommt es zu einer Organschädigung, die zum Absterben des gesamten 
% Organismus führen kann.}

\Layout{0}{\TxSymptome{} und Kennzeichen
}{
Das Problem bei den Schockkennzeichen ist, dass nicht immer alle sofort erkannt
werden. Die häufigsten bzw. typischsten sind hier angeführt.
Der Patient kann eine \EE{blasse Haut} haben. Weiters besteht die Möglichkeit, dass er
\EE{zittert und friert}, aber paradoxerweise dazu einen
\EE{kalten, klebrigen Schweiß} auf der Stirn hat. Der Patient kann außerdem \EE{verwirrt, unruhig
und ängstlich} sein. Er besitzt eine \EE{beschleunigte, flache Atmung} und
einen \EE{beschleunigten, kaum tastbaren Puls}.
\bigskip
\bigskip
}{
\begin{itemize}
  \item Blasse, kühle Haut
  \item Zittern und frieren
  \item Kalter, klebriger Schweiß
  \item Verwirrt, unruhig, ängstlich
  \item Beschleunigte, flache Atmung
  \item Beschleunigter, kaum tastbarer Puls
\end{itemize}
}{}{}{}


% \subsubsection{10.2) Schockkennzeichen:}
% 
% {
% {Die }{\textbf{\textit{allgemeinen Kennzeichen}}}{
% des Schocks:}}
% 
% {
%  i) blasse Haut}
% 
% {
%  ii) Schwitzen bei gleichzeitigem Zittern (Paradoxon)}
% 
% {
%  iii) schnelle, flache Atmung}
% 
% {
%  iv) große Angst und Unruhe (Patienten verwenden des Öfteren den
% Begriff \  {\quotedblbase}Todesangst{\textquotedblleft})}
% 
% {
%  v) zeitliche und örtliche Verwirrtheit}
% 
% \subparagraph{
% Ursachenbezogene Kennzeichen:}
% 
% {
%  i) Blut- und Flüssigkeitsverlust:}
% 
% {
%   a) stark blutende Wunde oder}
% 
% {
%   b) Verbrennungen/Verbrühungen}
% 
% {
%   c) chronischer Flüssigkeitsmangel durch beispielsweise dauernde
% körperliche   Überbelastung oder starker Hitzeeinwirkung}
% 
% {
%  ii) Vergiftung und allergische Reaktion:}
% 
% {
%   a) vergiftungstypische Symptome (näheres siehe unten)}
% 
% {
%   b) allergietypische Symptome (näheres siehe unten)}
% 
% {
%  iii) Herzversagen:}
% 
% {
%   a) stechender oder brennender Schmerz hinter dem Brustbein}
% 
% {
%   b) meist ein Ausstrahlen des Schmerzes in den linken Arm }
% 
% {
%   c) Atembeschwerden in Form des Gefühls eines schweren Gegenstandes
% auf   dem Brustkorb der einem daran hindert tief Luft zu holen}
% 
% {
%   c) je nach näheren Fall sind die Schmerzen entweder
% belastungsabhängig   oder nicht}

% \Layout{0}{Schockbekämpfung
% }{
% Da der Schock aus drei verschiedenen Ursachen entstehen kann, muss auch
% in der Behandlung ursachengerecht und damit differenziert vorgegangen
% werden. Z.B. die \EE{Ursache beseitgen} kann der Ersthelfer nur bei
% Blutverlust mittels Blutstillung. \EE{Beengende Kleidungsstücke öffnen} kann
% dagegen immer helfen. Das wichtigste ist jedoch die \EE{situationsgerechte
% Lagerung} des Patienten. Bei Bewusstlosigkeit stabile Seitenlage, bei Blut- und
% Flüssigkeitsverlust Oberkörper flach Beine hoch, bei Herzversagen mit leicht
% erhöhtem Oberkörper und bei Atemnot mit stark erhöhtem Oberkörper. Der
% \EE{Schutz vor Wärmeverlust}, ein \EE{beruhigender Zuspruch} und die
% \EE{Beobachtung} des Patienten sind Massnahmen, die der Ersthelfer bei jedem
% Patienten durchzuführen hat.
% }{
% \begin{itemize}
%     \item Ursache beseitigen
%     \item Beengende Kleidungsstücke öffnen
%     \item Situationsgerechte Lagerung
%     \item Schutz vor Wärmeverlust
%     \item Beruhigender Zuspruch
%     \item Beobachtung
% \end{itemize}
% }

\Layout{0}{Lagerung
}{
Da der Schock aus drei verschiedenen Ursachen entstehen kann, muss auch
in der Behandlung ursachengerecht und damit differenziert vorgegangen
werden.  Das wichtigste ist jedoch die \EE{situationsgerechte
Lagerung} des Patienten. Bei Bewusstlosigkeit stabile Seitenlage, bei Blut- und
Flüssigkeitsverlust wird der Oberkörper flach und die Beine hoch gelagert. Bei
Herzversagen wird Oberkörper leicht erhöht und bei Atemnot stark
erhöht gelagert. 

\bigskip

\bigskip

\bigskip
}{
\begin{itemize}
  \item Situationsgerecht, z.\,B.:
  \begin{itemize}
    \item Oberkörper hoch
    \item Beine hoch
    \item Stabile Seitenlage
    \item \ldots
  \end{itemize}
\end{itemize}
}{}{}{}

\MassnahmenNG{Erste-Hilfe-Maßnahmen}{Schockbekämpfung}{}{

\begin{itemize}
  \item Ursache beseitigen (wenn möglich), ggfs. Blutstillung
  \item Beengende Kleidungsstücke öffnen
  \item Situationsgerechte Lagerung
  \item Wärmeerhalt
  \item Beruhigender Zuspruch
  \item Beobachtung
\end{itemize}
}


% \subsubsection{10.3) Schockbekämpfung:}
% 
% {
% Da der Schock aus drei verschiedenen Ursachen entstehen kann, muss auch
% in der Behandlung ursachengerecht und damit differenziert vorgegangen
% werden. Der grds. Beginn der Behandlung ist bei allen gleich: man muss
% den Patienten unverzüglich in die für den ursächlichen Fall
% richtige Lagerung bringen.}
% 
% {
%  i) Blut- und Flüssigkeitsverlust sowie ii) Vergiftung und allergische
% Reaktion:}
% 
% {
%  Der Patient ist flach am Boden zu legen, seine Beine sind erhöht zu
% lagern.}
% 
% {\scriptsize
%  Durch die erhöhte Lagerung der Beine wird das in den Beinen
% vorhandene Blutvolumen durch  die Schwerkraft in den Korpus samt Kopf
% zurückgezogen und kann somit den vorhanden  Blutmangel ausgleichen.}
% 
% {
%  iii) Herzversagen:}
% 
% {
%  Der Patient ist mit aufrechtem Oberkörper am Boden zu legen.}
% 
% {\scriptsize
%  Durch die erhöhte Oberkörperlagerung wird das Blut von
% geschwächten Herzen mittels  Schwerkraft in den Großteil des
% Körpers gezogen, wodurch eine Entlastung des Herzens  eintritt. Bei
% besonders starken Atembeschwerden soll sich der Patient mit den Armen
% am  Boden aufstützen, um die Atemhilfsmuskulatur im
% Schultergürtelbereich einsetzen zu können}
% 
% 
% 
% 
% {
% Nachdem der Patient in die richtige Lagerung gebracht wurde, muss
% anschließend die Ursache des Schocks (falls überhaupt möglich)
% schnellstmöglich beseitigt werden.}
% 
% {
%  i) Blut- und Flüssigkeitsverlust:}
% 
% {
%  Blutstillung bzw. Flüssigkeitszufuhr:}
% 
% {\scriptsize
%  Nähere Informationen zur Blutstillung siehe unten. }
% 
% {
%  ii) Vergiftung und allergische Reaktion:}
% 
% {
% { Entfernen des Giftes }{vom}{ Körper des
% Patienten, nicht aber }{aus}{ dem Körper des 
% Patienten. }}
% 
% {\scriptsize
%  Sollte der Patient daher Gift über den Mund aufgenommen haben, darf
% er nicht zum  Erbrechen gebracht werden. Grund für dieses Verbot
% liegt in der Tatsache, dass das Gift  momentan im Magen eingeschlossen
% ist. Der Magen selbst ist mit einer äußerst resistenten 
% Magenschleimhaut beschichtet, die selbst der salzsäurehältigen
% Magensäure entgegenwirkt.  Wird der Patient nun zum Erbrechen
% gebracht, wandert das Gift samt der Magensäure die  Speiseröhre
% hinauf und führt zu einer erneuten Beschädigung des Mund-Rachen-
% Speiseröhrenraums. Dem Patienten ist stattdessen Wasser schluckweise
% zum Trinken zu  geben. Cave: Sollte der Patient Schaumbildner getrunken
% haben, darf auch kein Wasser  verabreicht werden, da es sonst zu einer
% Aufschäumung im Magen kommt, die wiederum ein  Erbrechen hervorruft.}
% 
% {\scriptsize
%  Ebenso darf das durch einen Schlangebiss oder eine Spritze in den
% Körper injizierte Gift nicht  aus dem Körper herausgesaugt werden.
% Grund hierfür ist die Tatsache, dass der Ersthelfer  beim Ansaugen
% das Gift über die Schleimhaut im eigenen Mundraum aufnehmen würde
% und  sich somit selbst gefährdet.}
% 
% {
%  iii) Herzversagen:}
% 
% {
%  Es gibt für den Ersthelfer keine direkten Behandlungsmöglichkeiten!
% Er muss sich  vorrangig auf die nachstehenden Maßnahmen
% beschränken.}
% 
% 
% 
% 
% {
% Weitere Maßnahmen sind:}
% 
% {
%  i) Öffnen von beengenden Kleidungsstücken}
% 
% {
%  ii) positive Beeinflussung der Atmung durch ruhiges Zusprechen und
% Voratmen}
% 
% {
%  iii) ersthelferischen Allzweckwaffen:}
% 
% {
%   a) Wärmeerhalt sichern}
% 
% {
%   b) laufende Kontrolle}
% 
% {
%   c) psychische Betreuung}




% \subsubsection{10.4) Konkretisierung - Blutstillung:}
% 
% {
% Die Blutstillung wird durchgeführt mittels Verbandsmaterial, welches
% der Ersthelfer ua. aus einem Verbandskoffer bzw. --tasche entnehmen
% kann. Bei diesen Verbandsboxen gibt es verschiedene Ausführungen, die
% sich entscheidend in ihrem Umfang unterscheiden. Der Umfang variiert je
% nachdem, für welchen Zweck und nach welcher (uU. gesetzlichen)
% Vorschrift die Box konzipiert wurde. Am Markt vorhanden sind:}
% 
% {
%  {\S} 102 KFG-Boxen: }
% 
% {
% { Nach dem Gesetz ist der Lenker eines Fahrzeuges (egal ob LKW,
% PKW, Motorrad,  etc.) verpflichtet zur Mitfuhr von
% }{\textit{{\quotedblbase}Verbandzeug, das zur Wundversorgung
% geeignet  und in einem widerstandsfähigen Behälter staubdicht
% verpackt und gegen  Verschmutzung geschützt
% ist{\textquotedblleft}}}{.}\footnote{\ {{\S} 102 Abs. 10
% erster Satz Kraftfahrzeuggesetz in BGBl. 267/1967 idF. BGBl.
% 94/2009}}{ In der Praxis sind Boxen, die nur dem {\S} 102 KFG 
% genüge tun, äußerst gering ausgestattet, weil das Gesetz einen
% viel zu großen  Spielraum frei lässt. Mitunter reichen {\S} 102
% KFG-Boxen nicht einmal zur Versorgung  eines einzelnen Patienten aus.}}
% 
% {
% Hersteller von Verbandsboxen können sich freiwillig den
% ÖNORM-Richtlinien unterwerfen, wodurch sie eine genau bestimmte
% Anzahl und Qualität von Verbandsmaterial in ihren Koffern und Taschen
% zu verkaufen haben. Von diesen ÖNORM-Verbandsboxen gibt es folgende
% Typen:}
% 
% {
%  ÖNORM V 5100: für Krafträder}
% 
% {
%  ÖNORM V 5101: für mehrspurige Kraftfahrzeuge}
% 
% {
%  ÖNORM 7 1020: für Betriebe und Schutzräume}
% 
% 
% 
% 
% {
% Materialen, die zur Blutstillung und Wundversorgung geeignet sind und
% sich unbedingt in einer Verbandsbox finden sollten, sind wie folgt:}
% 
% {
% i) medizinische Einweghandschuhe}
% 
% {
% ii) Dreiecktuch }
% 
% Graphische
% Darstellung des Dreiecktuchs und Bezeichnung seiner Bestandteile
% 
% 
% 
% 
% {
% Schlägt man die Basis mehrmals zu einem breiten Band zusammen, spricht
% man von einer Dreiecktuchkrawatte.}
% 
% {
% iii) Wundauflage iv) Mullbinde  v) Leukoplast (Klebestreifen)}
% 
% {
% vi) Schere  vii) Beatmungstuch iix) Pflaster}
% 
% {
% ix) Rettungsdecke}
% 
% 
% 
% 
% {
% Vor Durchführung der Blutstillung muss sich der Ersthelfer selbst
% schützen indem er sich medizinische Einweghandschuhe anzieht, um der
% Gefahr einer Ansteckung mittels von Blut übertragender Krankheiten
% vorzubeugen. Bei der Durchführung der Blutstillung sind folgende
% Methoden möglich, wobei leider nicht alle immer sinnvoll sind:}
% 
% {
%  i) Zudrücken:}
% 
% {
%  Der Ersthelfer drückt mit einer Wundauflage so fest auf die stark
% blutende Wunde, bis  diese aufhört zu bluten. }
% 
% {\itshape
%  Nachteil an dieser Methode ist, dass der Ersthelfer ab nun nicht mehr
% loslassen darf und somit  gehandicapt ist. }
% 
% {
%  ii) Abdrücken:}
% 
% {
%  Der Ersthelfer drückt bei Blutungen am Unterarm oder Unterschenkel
% das im  Oberarm oder Oberschenkel verlaufende blutzuführende
% Gefäß (Arterie) so fest zu,  bis die Wunde aufhört zu bluten. }
% 
% {\itshape
%  Nachteil hierbei ist nicht nur dasselbe Handicap wie oben, sondern auch
% die Tatsache, dass  der gesamte Unterarm oder Unterschenkel ohne
% Blutversorgung bleibt.}
% 
% {
%  iii) Abbinden:}
% 
% {
%  Der Ersthelfer versucht bei Blutungen am Unterarm oder Unterschenkel
% durch  Verwendung einer Dreiecktuchkrawatte den Oberarm oder
% Oberschenkel soweit  abzuschnüren, dass die Wunde aufhört zu
% bluten.}
% 
% {\itshape
%  Der Nachteil eines Handicaps ist zwar durch Verwendung der
% Dreiecktuchkrawatte nicht mehr  gegeben, es bleibt aber das Problem der
% nicht mehr gegebenen Blutversorgung. Um diesem  Nachteil zumindest ein
% wenig entgegenzutreten, muss der Ersthelfer die Abbindezeit am 
% Patienten notieren und spätestens 30 Min. nachher die Abbindung
% kurzzeitig wieder öffnen,  um den in 2-3 Stunden eintretenden
% Gewebstod hinauszuschieben. Aufgrund dieser massiven  Nachteile, darf
% das Abdrücken sowie die Abbindung nur bei Amputation oder 
% Gliedmaßeneinklemmung durchgeführt werden. Des Weiteren darf zur
% Abbindung nur ein  breites, weiches Tuch (wie die Dreiecktuchkrawatte)
% verwendet werden.}
% 
% {
%  iv) Druckverband:}
% 
% {
% { Der Ersthelfer deckt mit einer Wundauflage die blutende Wunde
% ab und positioniert  längs verlaufend der Wunde einen
% Druckkörper.}\footnote{\ {Druckkörper: Als Druckkörper
% eignen sich nur kompakte, weiche und saugfähige Objekte. Ein solches
% Objekt, das sich auch in der Verbandsbox findet, ist die
% Mullbinde.}}{ Über Wundauflage und Druckkörper  wird eine
% Dreiecktuchkrawatte dermaßen streng gewickelt, dass auf die Wunde
% ein  ausreichend starker Druck zur Stillung der Blutung ausgeübt
% wird. Dabei muss der  abschließende Knoten direkt auf dem
% Druckkörper gezogen werden.}}
% 
% {\scriptsize
%  Der Druckverband ist die beste Möglichkeit zur Stillung einer
% Blutung, da sie in fast allen  Situationen funktioniert und einfach
% sowie schnell umsetzbar ist. Eine Ausnahme der  Anwendung ist die
% Blutung aus der Halsschlagader. Hierbei besteht nur die Möglichkeit
% des  Zudrückens.}

\section{Notruf}
\index{Notruf}

% {
% Der Ersthelfer muss sich schon während der Durchführung seiner
% Erste-Hilfe-Maßnahmen daran denken weitere, professionelle Hilfe zu
% verständigen. }
% 
% {
% Für die zeitliche Koordinierung der Notrufabgabe muss der Ersthelfer
% Rücksicht auf die jeweils vorliegenden Situationen nehmen. Bei einem
% sich bei Bewusstsein befindlichen Patienten muss der Notruf
% spätestens nach setzen der letzten lebensrettenden
% Sofortmaßnahmen gesetzt werden. Bei einem bewusstlosen Patienten
% mit aufrechter Atmung und vorhandenem Kreislauf nach Umlagerung des
% Patienten in die stabile Seitenlagerung. Ist der Patient weder bei
% Bewusstsein noch hat er Atmung und Kreislauf, muss vor der
% Wiederbelebung eine Notruf abgegeben werden. }
% 
% {
% Dabei muss sich der Ersthelfer zwei Fragen stellen: Unter welcher Nummer
% kann ich Hilfe anfordern? Was muss ich beim Notruf angeben?}

\Layout{0}{Absetzen eines Notrufes}{
Der Ersthelfer muss darauf vorbereitet sein beim Notruf
folgende Fragen beantworten zu können.

\EE{Wo} ist der Notfallort? Ist es nicht möglich eine
Adresse anzugeben, muss der Notfallort so eindeutig wie
möglich beschrieben werden.

\EE{Was} ist passiert? Die Art des Notfalls bekannt geben. Es
reicht die Wahrnehmungen weiterzugeben, es müssen keine
Diagnosen gestellt werden. 

\EE{Wie viele} Leute sind betroffen? Es ist
notwendig, die Anzahl der Patienten bekannt zu geben, um
eine entsprechende Anzahl an Rettungsmitteln einzuteilen.

\EE{Wer} ruft an? Es ist sinnvoll
eine Rückrufnummer bekannt zu geben. Sollte der Notfallort z.\,B. nicht gefunden
werden, besteht die Möglichkeit zurückzurufen und nachzufragen.

Wichtig für die Aufnahme vollständiger Daten ist es, dass der
Ersthelfer nicht als Erster auflegt. Er muss warten bis der
Rettungsdienstmitarbeiter das Gespräch beendet und damit zu
verstehen gegeben hat, alle notwendigen Informationen erhalten zu haben.
}{
\begin{itemize}
  \item Wo?
  \item Was?
  \item Wie viele?
  \item Wer?
  \item Niemals als Erster auflegen!
\end{itemize}
}{}{}{}

\Layout{0}{Notrufnummern
}{
\index{Notruf!-nummern}
\index{112}
\index{122}
\index{133}
\index{144}
Die drei in Österreich geltenden Notrufnummern
lauten:
\EE{\Code{122} Feuerwehr}, \EE{\Code{133} Polizei} und \EE{\Code{144}
Rettung}. Vor einigen Jahren kam durch die EU eine weitere Notrufnummer dazu:
\EE{\Code{112} Euro-Notruf}. Die Besonderheit dieser Nummer liegt darin, dass
sie nicht nur in Österreich, sondern in allen Mitgliedsstaaten der EU anwendbar ist.
Bei Wählen dieser Nummer wird man zu einer Stelle weitergeleitet, bei
der das Personal mehrere Sprachen spricht, die
Heimatsprache und Englisch sind verpflichtend. 
% Für spezielle Notfälle gibt es auch spezielle 
% Notrufnummern. Die beiden relevantesten sind: \EE{\Code{01/406 43 43}
% Vergiftungsinformationszentrale} und \EE{\Code{141} Ärztefunkdienst}. 
}{
\begin{itemize}
  \item \Code{122} -- Feuerwehr
  \item \Code{133} -- Polizei
  \item \Code{144} -- Rettung
  \item \Code{112} -- Euro-Notruf
\end{itemize}
}{}{}{}

\Layout{0}{Besondere Dienste}
{\index{Aerztefunkdienst@Ärztefunkdienst}
\index{AeFD@ÄFD|see{Ärztefunkdienst}}
\index{Vergiftungsinformationszentrale}
Die
\EE{Vergiftungsinformationszentrale} (VIZ) dient nur dazu
Auskunft über die Risiken und Behandlungen von sicher festgestellten Vergiftungen zu
geben (\E{Expertenrat}). Die
Vergiftungsinformationszentrale dient grundsätzlich dem
medizinischen Fachpersonal um Expertenrat einholen zu
können. Der Ersthelfer muss im Vergiftungsnotfall in
erster Linie die Rettung alarmieren.

Der \EE{Ärztefunkdienst} (\Abk{ÄFD})
ist in Wien nur von Montag bis Freitag in der Nacht und Samstag, Sonn- und
Feiertag rund um die Uhr besetzt. Er
ist weder personal- noch materialtechnisch dafür ausgerüstet
Notfälle zu bewältigen. Sein Zweck liegt in der
Behandlung von allgemeinmedizinischen
"`Ordinationsfällen"'. Außerdem kann es zu Wartezeiten von
bis zu drei Stunden kommen.\footnote{Ein Umstand,
welcher leider viele Patienten dazu verleitet, lieber bei
der Rettung anzurufen, welche durchschnittlich in 11 Minuten
vor Ort ist.}. 
}{
\begin{description}
  \item[Vergiftungsinformationszentrale] \Code{01\,/\,406\,43\,43} --
  Expertenrat
  \item[Ärztefunkdienst] \Code{141} -- Allgemeinmedizinische Versorgung in
  der Nacht, an Feiertagen und Wochenenden
\end{description}
}{}{}{}




% \subsection{6.1) Notrufnummern:}
% 
% {
% Die drei traditionell in Österreich geltenden Notrufnummern sind wie
% folgt:}
% 
% \begin{table}[H]
% \Captionof{table}{Notrufnummern.}
% \begin{tabulary}{\linewidth}{LL}
% 122
% &
% Feuerwehr
% \\\addlinespace
% 133
% &
% Polizei
% \\\addlinespace
% 144
% &
% Rettung
% \\\addlinespace
% 112
% &
% Euro-Notruf
% \\\addlinespace
% 
% &
% 
% \\\addlinespace
% 140
% &
% Bergrettung
% \\\addlinespace
% 01/406 43 43
% &
% Vergiftungszentrale
% \\\addlinespace
% 141
% &
% Ärztefunkdienst
% \\\addlinespace
% 
% \end{tabulary}
% \end{table}
% 
% 
% 
% {
% Vor einigen Jahren kam durch die EU eine weitere Notrufnummer dazu:}
% 
% 
% 
% {
% Die Besonderheit dieser Nummer liegt darin, dass sie nicht nur in
% Österreich, sondern in allen Mitgliedsstaaten der EU anwendbar ist.
% Bei Wählen dieser Nummer wird man zu einer Stelle weitergeleitet, bei
% der das Personal mehrere Sprachen sprechen soll, insbesondere sind die
% Heimatsprache und Englisch verpflichtend.}
% 
% {
% Für spezielle Notfälle gibt es auch spezielle Notrufnummern. Die
% beiden relevantesten sind:}
% 
% 
% 
% {
% Selbst wenn ein unter den beiden speziellen Notrufnummern fallender
% Notfall vorliegt, kann der Ersthelfer die allgemeinen Nummern der
% Rettung (144 oder 112) wählen. Besonders zu behandeln ist die
% Notrufeinrichtung der Vergiftungszentrale. Diese Einrichtung ist nicht
% im eigentlichen Sinn eine Notrufnummer, da man über sie keine
% weiteren Hilfskräfte anfordern kann. Zweck der Vergiftungszentrale
% ist es lediglich Auskunft über die Risiken und Behandlungen von
% sicher festgestellten Vergiftungen zu geben. Daher sollte man in jenen
% Fällen, in denen weitere Hilfskräfte notwendig sind oder man sich
% über die Einschätzung der Lage unsicher ist, die allgemeinen
% Rettungsnotrufnummern (144 oder 112) wählen.}
% 
% {\scriptsize
% Alle Notrufnummern sind kostenlos. Daher kann man bei Telefonzellen ohne
% Einwurf von Münzen oder Karten bzw. bei Mobiltelefonen ohne
% aufgeladener Wertkarte einen Notruf abgeben. Bei Festnetzapparaten
% funktionieren die Nummern sogar bei Stromausfall.}
% 
% {
% Als eine weitere dreistellige Nummer ausgeprägt, wird eine weitere
% Telefonnummer oft in der Bevölkerung falsch als Notrufnummer
% gewertet. Gemeint ist:}
% 
% {
%  141 Ärztefunkdienst}
% 
% {
% Bei dieser Nummer handelt es sich um keine Notrufnummer!!! Es ist eine
% Einrichtung bei der Hausärzte zu jenen Zeiten mobil unterwegs sind,
% wenn keine Hausarztordination offen hat. Also unter der Woche in der
% Nacht und am Wochenende sowie Feiertags ganztags. Der Ärztefunkdienst
% ist weder personal- noch materialtechnisch dafür ausgerüstet
% Notfälle behandeln zu können. Sein Zweck liegt in der Behandlung
% von allgemein-ärztlichen Ordinationsfällen. Ein weiteres
% entscheidendes Kriterium gegen eine Notrufabgabe unter 141 liegt in den
% Ausfahrtszeiten von durchschnittlich 2-3 Stunden. Die Rettung unter 144
% bzw. 112 benötigt hingegen statistisch 11 Minuten bis sie am Notfall-
% bzw. Unfallort eingetroffen ist.}
% 
% 
% 
% 
% \subsubsection{6.2) Angaben beim Notruf:}
% 
% {
% Der Ersthelfer sollte sich beim Notruf die folgenden 5 W-Fragen selbst
% beantworten können, um dem Rettungsjournal dienliche Informationen
% geben zu können.}
% 
% \begin{enumerate}
%     \item Wer ruft an?
%     \item  Was ist geschehen?
%     \item Wo ist es geschehen?
%     \item Wie viele Leute sind betroffen?
%     \item  Wann trat das Geschehnis ein?
% \end{enumerate}
% 
% {
% Wichtig für die Aufnahme vollständiger Daten ist es, dass der
% Ersthelfer nicht als Erster auflegt. Er muss warten bis der
% Rettungsjournaler das Gespräch beendet und damit zu verstehen gegeben
% hat, alle notwendigen Informationen erhalten zu haben.}
% 
% {
% Ebenso wichtig ist es die eigene Rückrufnummer anzugeben, damit sich
% der Rettungsjournaler oder das Einsatzpersonal bei Notwendigkeit
% weiterer Informationen an den Ersthelfer richten kann. Fälle in denen
% dies notwendig wird, sind Situationen eines abgelegenen, schwer
% auffindbaren Notfall- bzw. Unfallortes.}
% 
% {
% Sollte der Ersthelfer zu jenem Zeitpunkt an dem der Notruf abzugeben
% notwendig wird, nicht (mehr) alleine sein, muss es zu einer
% Arbeitsteilung zwischen den Ersthelfern kommen. Während der eine
% Helfer sich weiter mit der Behandlung des Patienten auseinandersetzt,
% soll der andere die Notrufabgabe durchführen.}


\section{Weitere Erste Hilfe}

Die weitere Erste Hilfe umfasst Maßnahmen zur Versorgung von
Verletzungen, Erkrankungen und Vergiftungen. Manche Autoren erfassen
unter diesem Begriff Umstände, die
angeblich weniger lebensbedrohlich sein sollen, als die zuvor behandelten. Tatsächlich
können sich auch die folgenden Beeinträchtigungen der
Patientenintegrität zu lebensbedrohlichen Situationen entwickeln. 
% Das
% Vorgehen für den Ersthelfer bleibt aber -- wie oben beschrieben --
% gleich. Nachdem er sich einen Überblick verschaffen konnte und die
% Gefahren beseitigt hat, wird er bei der Bewusstseinskontrolle einen
% ansprechbaren, aber in seiner Integrität beeinträchtigten Patienten
% vorfinden. Dieser Patient ist wie folgt zu behandeln.}



\subsection{Wunden und andere spezielle Situationen}
\index{Wunden!Erste-Hilfe}
\label{topic:wunden-ehi}


Jede Wunde ist für unseren Körper eine Gefahr, auch wenn sie noch so klein
ist. Grund dafür ist das mögliche Eindringen von Krankheitserregern wie z.\,B.
Tetanus. 

\Layout{0}{Bagatellverletzungen}
{%
Bagatellverletzungen sind \EE{kleine, oberflächliche
Wunden}, wie z.\,B. Schürfwunden und kleine Schnitte, welche
nicht tiefer als \nicefrac{1}{2}\CM* und nicht länger als
2\CM* sind, gemeint. Bei idealer Behandlung durch den
Ersthelfer müssen diese nicht von einem Arzt versorgt
werden.
}{
\begin{itemize}
  \item kleine, oberflächliche Wunden
\end{itemize}
}{}{}{}

\MassnahmenNG{Erste-Hilfe-Maßnahmen}{Bagatellverletzungen}{}{

\begin{itemize}
  \item Wunde reinigen
  \item Evtl. steriler Wundverband
  \item Auf den Tetanusschutz achten
\end{itemize}
}

\Layout{0}{Wunden, die unbedingt ärztlich versorgt werden müssen.}
{
Alle anderen Wunden, welche tiefer, größer etc. sind, müssen unbedingt von
einem Arzt versorgt werden. Für diese Fälle gibt es \E{zwei Gebote} für den
Ersthelfer, welche er unbedingt einhalten muss: 

\begin{enumerate}
    \item \EE{Es kommt nichts in die Wunde
hinein}, wie z.\,B. Wundpuder oder Wundsalben. Dies muss sonst im
Spital wieder mühevoll aus der Wunde entfernt werden, um sie
gründlich zu reinigen.
\item \EE{Es kommt nichts aus der Wunde heraus.}
Fremdkörper, welche ein Blutgefäß durchtrennen, können dieses eventuell abdichten. Wenn nun der Fremdkörper entfernt wird, kommt es zu einer stark blutenden Wunde. 

\end{enumerate}
 Die anschließende
Tabelle gibt einen groben Überblick über Wunden, welche unbedingt von einem
Arzt versorgt werden müssen.
}{
\begin{itemize}
  \item Zwei Gebote:
  \begin{itemize}
    \item Es kommt nichts hinein
    \item Es kommt nichts heraus
  \end{itemize}
\end{itemize}
}{}{}{}

\Layout{57}{}
{}{}
{
\begin{tabulary}{\linewidth}{LL}
\addlinespace
\noindent\EE{Kein ausreichender Tetanus-Impfschutz}
&
Impfung innerhalb von 72 Stunden notwendig
\\\addlinespace
\noindent\EE{Klaffende Wunden}
&
\ldots\ müssen oft genäht werden. 
\\\addlinespace
\noindent\EE{Verbrennungen, Verbrühungen}
&
Einschätzung der Bedrohung anhand von Verbrennungsgrad und Ausdehnung
(\prettyref{topic:verbrennung-ehi}) 
\\\addlinespace
\noindent\EE{Wunden mit Fremdkörper}
&
sofern keine Bagatellverletzung (z.\,B. Schiefer, \ldots )
\\\addlinespace
\noindent\EE{Augenverletzungen aller Art}
&
Bei Augenverletzungen immer beide Augen verbinden.
\\\addlinespace
\noindent\EE{Stichwunden}
&
Vom Ersthelfer kann nicht festgestellt werden wie tief diese sind.
\\\addlinespace
\noindent\EE{Wunden an Gelenken, Hand- und Fußrücken}
&
\ldots\ mit gleichzeitiger Bewegungseinschränkung
\\\addlinespace
\noindent\EE{Bisswunden}
&
Infektionsgefahr
\\\addlinespace
\noindent\EE{Insektenstiche}
&
Bei allergischen Reaktionen oder bekannter entsprechender Allergie
\\\addlinespace
\noindent\EE{Schusswunden}
&

\\\addlinespace
\noindent\EE{Wunden im Genitalbereich}
&

\\\addlinespace\bottomrule
\end{tabulary}
}
{Wunden, die unbedingt vom Arzt versorgt werden müssen!}{}


\subsubsection{Verbrennungen, Verbrühungen}
\index{Verbrennung!Erste-Hilfe}
\label{topic:verbrennung-ehi}

\Layout{0}{Einteilung
}{%
\index{Verbrennungspanzer}%
Um festzustellen, wie bedrohlich die Verbrennung für den Patienten ist, muss der
Schweregrad der Schädigung der Haut eingeteilt werden. Dies geschieht mittels drei Grade. Beim \EE{1.~Grad} kommt es zur Rötung der Haut, einer
Schwellung des Gewebes und zu Schmerzen. Der \EE{2.~Grad} ist charakterisiert
durch Blasenbildung. Es bestehen weiterhin
Schmerzen. Beim \EE{3.~Grad} kommt es zur \T{Schorfbildung}
(\E{Verbrennungspanzer}). Die Schmerzen sind abhängig von
der Tiefe der Schädigung der Haut. }{
\begin{itemize}
  \item 1\textdegree : Rötung, Schwellung, Schmerzen
  \item 2\textdegree : Blasenbildung, Schmerzen
  \item 3\textdegree : Schorfbildung, Verbrennungspanzer
\end{itemize}
}{}{}{}



\Layout{22}{Ausdehnung: Handregel und Neunerregel}
{
Genauso wichtig wie der Schweregrad ist die Ausdehnung der
Verbrennung. Damit ist gemeint, wieviel \% der Körperoberfläche betroffen sind.
Dazu gibt es zwei Möglichkeiten. Einerseits die \T{Handregel}, welche
besagt, dass die \EE{Handfläche des Patienten  1\PC*}
seiner Körperoberfläche entspricht. Da dies in der Praxis
bei großflächigen Verbrennungen meist umständlich ist, gibt
es auch noch die \T{Neunerregel}. Hier wird der Körper in
9\PC*-Teile eingeteilt.

\begin{center}
\begin{tabulary}{\linewidth}[H]{LCC}
\toprule
 & \% &
\\\midrule
Kopf:
 & 9  &
\\
Arm, je:  &  9  &
\\
Rumpf vorne: & 18  &
\\
Rumpf hinten: & 18 &
\\
Bein, je: & 18 &
\\
Genitalien:  & 1 &
\\\bottomrule
\end{tabulary}
\end{center}
% \Captionof{table}{Neunerregel} \label{tab:Neunerregel}
}{
\begin{itemize}
  \item Handregel
  \item Neunerregel
\end{itemize}
}
{./images/noauth/AASdev20090228-img94.png}
{Neunerregel.}{}

\Cave{\Ca{Schockgefahr!} 

Bei Kindern ist ab \EE{5\PC*}, bei Erwachsenen ab \EE{10\PC*}
mind. 2\textdegree{} verbrannter Körperoberfläche mit einem Schock zu
rechnen!}

\MassnahmenNG{Erste-Hilfe-Maßnahmen}{Verbrennungen}{}{
\begin{itemize}
  \item Hitzeeinwirkung sofort stoppen.
  \item Wenn innerhalb von 2 Minuten möglich, für max. 10\Min* mit handwarmen
  Wasser kühlen (Richtwert 20\,\textcelsius); andernfalls
  keine Kühlung durchführen \cite{ASBOEAkademieUpdate2010Verbrennung}.
  
  \Ca{Achtung: Gefahr der Unterkühlung!}
  \item Wärmeerhalt ist besonders wichtig!
  \item Keimfreier Wundverband
  \item Schockbekämpfung
  \item \Ca{Keine Salben oder Puder etc. verwenden}
\end{itemize}
}

\clearpage % TODO temporäres clearpage
\subsubsection{Erfrierungen}

\Definition{Durchblutungsstörungen mit Gewebstod an den Körperenden}

\MassnahmenNG{Erste-Hilfe-Maßnahmen}{Erfrierungen}{}{

\begin{itemize}
  \item Keimfreier Wundverband
  \item Restlichen Körper aufwärmen
  \item Warme, gezuckerte Getränke verabreichen
  \item Keinen Alkohol zu trinken geben!
  \item Nicht frottieren!
\end{itemize}
}

\subsubsection{Unterkühlung}

\Definition{Absinken der Körperkerntemperatur unter 36\,\textcelsius{}}

\Layout{0}{\TxSymptome{} und Kennzeichen}
{
Am Beginn hat der Patient \EE{Schmerzen}, Atem- und Herzfrequenz sind erhöht.
Je weiter die Temperatur sinkt, desto \EE{teilnahmsloser}
und \EE{müder} wird der Patient. Zu Letzt kommt es zum \EE{Absinken
von Atem- und Herzfrequenz}. Bei einer Temperatur unter 27\,\textcelsius{} tritt der Tod ein.
}{
\begin{itemize}
  \item Schmerzen
  \item Teilnahmslosigkeit/Müdigkeit
  \item Absinken von Atem- und Herzfrequenz
\end{itemize}
}{}{}{}

\MassnahmenNG{Erste-Hilfe-Maßnahmen}{Unterkühlung}{}{

\begin{itemize}
  \item Aufwecken und wach halten
  \item Vorsichtig in Decken oder Kleidungsstücke einwickeln
  \item Rettungsdienst verständigen
  \item Warme Getränke wenn Patient bewusstseinsklar ist
  \item Keinesfalls bewegen oder frottieren: sog.
  \Ca{Bergungstod!}\footnote{Unter dem Begriff \emph{Bergungstod} versteht
  man das Versterben des Patienten während der Rettung in Folge eines Blutrückflusses
  aus der kalten Körperschale in den wärmeren Körperkern, welcher durch die
  Bewegung des Patienten während der Rettung ausgelöst wird.}
\end{itemize}
}

\subsubsection{Verstauchung}
\label{topic:verstauchung-eh}

\Definition{Überdehnung der Bänder bzw. des Kapselapparates eines Gelenkes.}

\Layout{0}{\TxSymptome{} und Kennzeichen}
{%
Der Patient hat \EE{Schmerzen} im Bereich des betroffenen Gelenkes, es
kann auch zu einer \EE{Schwellung} rund um die geschädigte Stelle kommen. Mit
unter nicht gleich vorhanden, sondern erst am nächsten Tag ersichtlich,
kann sich auch ein \EE{Bluterguss} bilden. Auch eine \EE{Einschränkung der
Bewegung} ist -- schmerzbedingt -- möglich. \bigskip
}{
\begin{itemize}
  \item Schmerzen
  \item Schwellung
  \item Bluterguss
  \item Bewegungseinschränkung
\end{itemize}
}{}{}{}

\MassnahmenNG{Erste-Hilfe-Maßnahmen}{Verstauchung}{}{

\begin{itemize}
  \item Ruhigstellung
  \item Kalte Umschläge
  \item Unfallabteilung aufsuchen
  \item Hochlagern % ZWING 2010-03-xx
\end{itemize}
}


\subsubsection{Verrenkung}
\label{topic:verrenkung-eh}

\Definition{Gelenkverletzung, bei der zwei das Gelenk bildende Knochenenden
aus ihrer normalen Stellung zueinander verschoben werden. Dabei kommt es oft zu
Begleitverletzungen an der Gelenkkapsel bzw. den Gelenkbändern. \cite{Gorgass7}}



\Layout{0}{\TxSymptome{} und Kennzeichen}
{%
Der Patient hat starke \EE{Schmerzen} im Bereich des verletzten Gelenkes,
die betroffene Extremität kann nicht mehr bewegt werden
(\EE{Bewegungsunfähigkeit}). Außerdem kommt es zu einer \EE{abnormen
Stellung}, welche \EE{federnd fixiert} gelagert ist. 
\bigskip
}{
\begin{itemize}
    \item Schmerzen
    \item Bewegungsunfähigkeit
    \item Abnorme Stellung, "`federnd fixiert"'
\end{itemize}
}{}{}{}

\MassnahmenNG{Erste-Hilfe-Maßnahmen}{Verrenkung}{}{

\begin{itemize}
    \item Vorgefundene Stellung beibehalten
    \item Ruhigstellung
    \item Rettungsdienst verständigen
    \item Keine Einrenkversuche!
\end{itemize}
}


\subsubsection{Knochenbrüche}
\label{topic:knochenbruch-eh}

\Layout{0}{Arten}
{
Grundsätzlich wird zwischen zwei Arten von Knochenbrüchen
unterschieden: Beim \EE{geschlossenen Knochenbruch} ist
die Haut über der Bruchstelle noch intakt, während beim
\EE{offenen Knochenbruch} aufgrund der Gewalteinwirkung von
außen die Haut über der Bruchstelle verletzt ist. Es muss dabei nicht der Knochen herausragen.
}{
\begin{itemize}
    \item Geschlossener Knochenbruch
    \item Offener Knochenbruch
\end{itemize}
}{}{}{}

\Layout{0}{\TxSymptome{} und Kennzeichen}
{
Es gibt \E{sichere} und \E{unsichere} Kennzeichen für einen Knochenbruch. 

Zu
den \T{sicheren Kennzeichen} zählt die \EE{Fehlstellung} (abnorme Stellung,
Verkürzung, Treppenbildung), die \EE{abnorme Beweglichkeit} und die
\EE{Krepitation} (das hörbare und fühlbare Aneinanderreiben
der Knochen). 

\T{Unsichere
Knochenbruchzeichen}\index{Knochenbruchzeichen!Unsichere}
sind \EE{Schmerzen}, \EE{Schwellungen},
\EE{Bewegungseinschränkungen} und \EE{Blutergüsse}.

\bigskip
\bigskip
\bigskip
\bigskip
\bigskip
}{
\begin{itemize}
    \item Sichere Kennzeichen
    \begin{itemize}
        \item Fehlstellung
        \item Abnorme Beweglichkeit
        \item Krepitation
    \end{itemize}
    \item Unsichere Kennzeichen
    \begin{itemize}
        \item Schmerzen
        \item Schwellung
        \item Bewegungseinschränkung
        \item Bluterguss
     \end{itemize}
\end{itemize}
}{}{}{}

\Layout{0}{\TxGefahren{}}
{
Ein Knochen ist lebendes
Gewebe, daher ist er auch gut durchblutet. Eine der Gefahren des
Knochenbruchs ist daher der \EE{Blutverlust}. Je größer der Blutverlust ist
umso größer ist wiederum die Gefahr, dass der Patient in den \EE{Schock}
gerät. Bei einem offenen Knochenbruch besteht außerdem die Gefahr einer \EE{Wundinfektion}. Bei Rippenbrüchen kann
es zu \EE{Begleitverletzungen innerer Organe} kommen.
\bigskip
}{
\begin{itemize}
    \item Blutverlust
    \item Schock
    \item Wundinfektion
    \item Begleitverletzungen innerer Organe
\end{itemize}
}{}{}{}

\MassnahmenNG{Erste-Hilfe-Maßnahmen}{Knochenbrüche}{}{

\begin{itemize}
    \item Ruhig liegen lassen
    \item Wundversorgung bei offenen Knochenbrüchen
    \item Beengende Kleidungsstücke lockern
    \item Schuhe öffnen, nicht ausziehen
    \item Ruhigstellung durch unterstützte Lagerung
    \item Schockbekämpfung
\end{itemize}
}

\subsubsection{Vergiftungen}
\label{topic:vergiftungen-ehi}

\Layout{2}{Allgemeines}
{
Es gibt mehrere Möglichkeiten der Giftaufnahme. Man kann
es beispielsweise verschlucken, einatmen oder über die Haut aufnehmen. Das Gift
selbst kann aus dem Haushalt, der Umwelt (Pflanzen, Pilze, \ldots), der
Industrie, dem Gewerbe oder der Landwirtschaft kommen.
}{
}{}{}{}

\Layout{0}{\TxSymptome{} und Kennzeichen}
{
Je nach Art des Giftes und der Menge die der Körper abbekommen hat, gibt es
unterschiedliche Kennzeichen. Es kann mit \EE{Rauschzuständen} und
\EE{Erbrechen bzw. Durchfall} beginnen, über \EE{Bewusstseinsstörung} zu
\EE{Atem- und Kreislaufstörung} führen.
\bigskip
\bigskip
}{
\begin{itemize}
    \item Rauschzustände
    \item Erbrechen, Durchfall
    \item Bewusstseinsstörung
    \item Atem- und Kreislaufstörung
\end{itemize}
}{}{}{}

\MassnahmenNG{Erste-Hilfe-Maßnahmen}{Vergiftungen}{}{
\begin{itemize}
    \item Patient nicht ansprechbar:
    \begin{itemize}
        \item Retten aus der Gefahr (Selbstschutz!).
        \item Lebensrettende Sofortmaßnahmen
    \end{itemize}
    \item Patient ansprechbar:
    \begin{itemize}
        \item Wenn das Gift bekannt ist, Vergiftungsinformationszentrale
        (\Code{01\,/\,406\,43\,43}) anrufen und \E{Expertenrat} einholen.
        \item Wenn Gift unbekannt: Betreuen und → Notruf/Krankenhaus.
        \item Selbstschutz!
        \item Giftstoffe sichern und mitnehmen.
    \end{itemize}
\end{itemize}
}

\subsubsection{Verätzungen}

\MassnahmenNG{Erste-Hilfe-Maßnahmen}{Verätzungen der Haut}{}{

\begin{itemize}
   \item \Ca{Selbstschutz} beachten!
    \item Mit ätzender Substanz getränkte Kleidung entfernen
    \item Haut mit reinem Wasser abspülen
    \item Keimfreier Wundverband
    \item Evtl. Schockbekämpfung
\end{itemize}
}

\MassnahmenNG{Erste-Hilfe-Maßnahmen}{Verätzungen der Augen}{}{

\begin{itemize}
    \item Auge mit reinem Wasser von innen nach außen ausspülen. Abwasser
    \E{nicht über das gesunde Auge} abrinnen lassen
    \item Auge mit keimfreiem Wundverband bedecken; zum Ruhigstellen \E{beide}
    Augen bedecken
    \item Selbstschutz beachten!
\end{itemize}
}

\MassnahmenNG{Erste-Hilfe-Maßnahmen}{Verätzungen des
Verdauungstraktes}{}
{ 
\begin{itemize}
    \item Mund mit Wasser ausspülen
    \item Substanz bekannt → Vergiftungsinformationszentrale anrufen
    \item ev. Schockbekämpfung
    \item Selbstschutz beachten!
\end{itemize}
}

\TODO{????}{0.8.0}{Fehlt: EHI Tierbisse, Tollwut, Schlangenbiss, ZEckenbiss,
Insektenstiche Mund/Rachenraum}{}
\subsection{Wundverbände}
\label{topic:wundverbaende}

\Layout{0}{\TxBeschreibung}{
Ein Wundverband besteht aus einer keimfreien Wundauflage und einer Befestigung.
}{
\begin{itemize}
  \item keimfreie Wundauflage
  \item Befestigung
\end{itemize}
}{}{}{} 

\subsubsection{Pflasterwundverband}
\Layout{20}{}{
Pflasterverbände können nahezu überall eingesetzt werden,
sie müssen dazu unter Umständen passend zugeschnitten
werden. Wenn ein Pflasterverband über ein bewegliches
Gelenk angelegt werden muss, wird an der beweglichen Stelle
der Klebestreifen von außen keilförmig eingeschnitten.
}{}{./images/koch-aass/00800-gamma/Pflasterb.JPG}{}{}

\subsubsection{Dreiecktuchverbände}

Kopfverbände:

\Layout{57}{}
{}{}
{
\begin{tabularx}{\linewidth}{XXX}
\includegraphics[width=\linewidth]{./images/koch-aass/00800-gamma/Kopfverband1b.JPG}
&
\includegraphics[width=\linewidth]{./images/koch-aass/00800-gamma/Kopfverband2b.JPG}
&
\\
Verband 1 &
Verband 2 &
\\\addlinespace\bottomrule
\end{tabularx}
}{Bilderserie: Kopfverbände}{}

Schulterverband:

\Layout{57}{}
{}{}
{
\begin{tabularx}{\linewidth}{XXX}
\includegraphics[width=\linewidth]{./images/koch-aass/00800-gamma/Schulterverband1b.JPG}
&
\includegraphics[width=\linewidth]{./images/koch-aass/00800-gamma/Schulterverband2b.JPG}
&
\includegraphics[width=\linewidth]{./images/koch-aass/00800-gamma/Schulterverband3b.JPG}
\\
Schritt 1 &
Schritt 2 &
Schritt 3 
\\\addlinespace\bottomrule
\end{tabularx}
}{Bilderserie: Schulterverband}{}

Handverbände:

\Layout{57}{}
{}{}
{
\begin{tabularx}{\linewidth}{XXX}
\includegraphics[width=\linewidth]{./images/koch-aass/00800-gamma/Handverband1b.JPG}
&
\includegraphics[width=\linewidth]{./images/koch-aass/00800-gamma/Handverband2b.JPG}
&
\includegraphics[width=\linewidth]{./images/koch-aass/00800-gamma/Handverband3b.JPG}
\\
Schritt 1 &
Schritt 2 &
Schritt 3 \\\addlinespace
\includegraphics[width=\linewidth]{./images/koch-aass/00800-gamma/Handkrawatte1b.JPG}
&
\includegraphics[width=\linewidth]{./images/koch-aass/00800-gamma/Handkrawatte2b.JPG}
&
\includegraphics[width=\linewidth]{./images/koch-aass/00800-gamma/Handkrawatte3b.JPG}
\\
Schritt 1 &
Schritt 2 &
Schritt 3
\\\addlinespace\bottomrule
\end{tabularx}
}{Bilderserie: Zwei verschiedene Handverbände}{}


Knieverband:

\Layout{57}{}
{}{}
{
\begin{tabularx}{\linewidth}{XXX}
\includegraphics[width=\linewidth]{./images/koch-aass/00800-gamma/Knieverband1b.JPG}
&
\includegraphics[width=\linewidth]{./images/koch-aass/00800-gamma/Knieverband2b.JPG}
&
\includegraphics[width=\linewidth]{./images/koch-aass/00800-gamma/Knieverband3b.JPG}
\\
Schritt 1 &
Schritt 2 &
Schritt 3 
\\\addlinespace\bottomrule
\end{tabularx}
}{Bilderserie: Knieverband}{}


\section*{Weiterführende Literatur}

\cite{BuchfelderErsteHilfe4,Erc2005Sec9,Erc2005Sec2,Erc2005Sec8,Erc2005Sec6,WirthEH,HolaubekASB16h,KoehnleinEH,KeggenhoffEH,KarutzKursbuchEH}

% {
% Verletzungen sind Beeinträchtigungen der körperlichen (und uU.
% geistigen) Integrität des Patienten durch äußere
% Gewalteinwirkung. Sie können genauso wie ihre Ursachen vielfältig
% auftreten. Die nachfolgende Abhandlung ist keine vollständige
% Auflistung von denkbaren Verletzungen oder Behandlungsarten! Sie
% versucht nur einen Überblick über die meist für einen Ersthelfer
% relevanten Verletzungen zu bieten!}
% 
% 
% 
% \subsubsection{11.1) Arten von Verletzungen und ihre
% ersthelferische Behandlung:}
% 
% {
% An sich gibt es keine {\quotedblbase}harmlosen{\textquotedblleft}
% Verletzungen, da jede Verletzung Eintrittsstelle für gefährliche
% Keime sein kann. Daher ist auch die grds. Behandlung einer Wunde immer
% gleich, ein keimfreier Wundverband. Jene Verletzungen die aber
% schwerwiegend genug sind, um eine ärztliche Untersuchung aufzusuchen,
% sind im Überblick wie folgt:}
% 
% {
%  i) jede Wunde ohne Tetanus-Impfschutz}
% 
% {
%  ii) klaffende Wunde}
% 
% {
%  iii) Knochenbruch}
% 
% {
%  iv) Verrenkung/Verstauchung}
% 
% {
%  v) Verbrennung/Verbrühung}
% 
% {
%  v) Erfrierung}
% 
% {
%  vi) Wunden mit Fremdkörpern}
% 
% {
%  vii) Stichwunden}
% 
% {
%  iix) Schusswunden}
% 
% {
%  ix) Augenverletzungen}
% 
% {
%  x) Bisswunden}
% 
% {
%  xi) Insektenstiche}
% 
% {
%  xii) Wunden im Genitalbereich}
% 
% {
%  xiii) sonstige körperregionsspezifische Verletzungen}
% 
% 
% 
% 
% \subsubsection{11.1.1) jede Wunde ohne ausreichenden Tetanus-Impfschutz: }
% 
% {
% Die auch als Wundstarrkrampf bezeichnete Infektion mit Tetanusbakterien
% führt nach einer Inkubationszeit von ca. 14 Tagen zu Krämpfen, die
% im Gesicht mit dem risus sardonicus (sartanisches Lächeln) beginnen
% und sich anschließend auf den gesamten Körper ausbreiten. Als
% Folge dieser Krämpfe droht der Tod, wodurch eine Schutzimpfung für
% jeden Menschen überlebenswichtig ist. Eine Impfung muss alle 8 Jahre
% ergehen.}
% 
% 
% 
% 
% \subsubsection{11.1.2) klaffende Wunden:}
% 
% {
% Sind Verletzungen bei denen eine Gewalteinwirkung größere
% Gewebseinrisse an der Körperoberfläche hervorgerufen hat.
% Arztpflichtig werden Verletzungen dieser Art ab einer Länge von
% über 2\cm* und/oder einer Tiefe von über 0,5 cm.}
% 
% {
% {Die ersthelferische Versorgung solcher Wunden entspricht den
% }{\textbf{\textit{allgemeinen Verbandsregeln}}}{. Bei den
% Verbandsregeln kann zwischen einer Mullbinden- und einer
% Dreiecktuchversorgung unterschieden werden. Beiden gleich ist, dass die
% Behandlung mit der Auflage einer keimfreien Wundauflage beginnt. Die
% Wundauflage darf nur an einem Rand angefasst aus der Verpackung
% genommen werden, die anhand der an einem Ende zu findenden Flügeln
% vorsichtig zu öffnen ist.}}
% 
% 
% 
% 
% \subparagraph{
% Mullbindenverband:}
% 
% {
% In früheren Zeiten mussten komplizierte Verbandswickelarten gelernt
% werden, da das Mullbindenmaterial aus schlechtwertigen, insbesondere
% nicht selbsthaftenden Gewebsstoff bestand. In den heutigen Zeiten kann
% hingegen auf diese Methoden verzichtet werden. Entscheidend ist für
% den Ersthelfer nur das Wissen über den richtigen Anfang eines
% Mullbindenverbandes. Hierfür empfiehlt es sich den ersten Umschlag
% über die Wundauflage schräg verlaufend anzusetzen, um im darauf
% folgenden Umschlag gerade zu verlaufen. Durch den ersten schrägen
% Umschlag steht eine Spitze des Mullbindenanfangs seitlich heraus, die
% beim dritten Umschlag mit eingewickelt werden kann. Durch diese Methode
% kann der Anfangspunkt nicht verrutschen und der restliche Verband
% bleibt stabil. Zu beenden ist der Verband bei selbsthaftenden
% Mullbinden durch einfaches Abschneiden des Bandes von der
% Mullbindenspule. Bei nicht selbsthaftenden Mullbinden ist das Ende
% mittels Klebestreifen anzukleben oder mittels Knoten zu verzurren.
% Hilfsmittel für Dauerverbände (sogen. Schwiegermutter: ein
% elastisches Band mit je zwei Metallhaken an jeder Seite) sind bei den
% ersthelferischen Momentverbänden nicht anzuwenden.}
% 
% 
% 
% 
% \subparagraph{
% Dreiecktuchverband:}
% 
% {
% {Je nach Position der Wunde kann das
% Dreiecktuch}\footnote{\ {Dreiecktuch vers.
% Dreieck}{s}{tuch: Die im Bundesdeutschenraum verwendete
% Bezeichnung Dreieck}{s}{tuch mit der Verwendung des
% typisch {\quotedblbase}deutschen Verbindungs-S{\textquotedblleft}
% sollte aus Gründen eines - wenn auch geringen --
% Patriotismusgefühls nicht verwendet werden. Die österreichische
% Rechtschreibung kommt auch ohne diesen aus.}}{ verschieden
% angewendet werden. Bei allen Varianten bietet das Dreiecktuch den
% Vorteil schnell, einfach und effektiv zu funktionieren, wodurch es
% selbst in dem höher technischen Sanitätsdienst immer das Mittel der
% ersten Wahl sein sollte. Anwendungsmöglichkeiten sind wie folgt:}}
% 
% {
%  i) Kopfverletzung:}
% 
% {
%  Am Beginn wird die Basis eingeschlagen, um eine stabile Kante zu
% bilden. Das Tuch  wird an der Hutkrempenkante entlang angelegt und die
% beiden Enden werden über  der Spitze einfach verknotet. Danach wird
% mittels Ziehen an der Spitze der Verband  straff gezogen, die Spitze in
% den einfachen Knoten eingewickelt und mittels  Weberknoten
% festgezogen.}
% 
% {
%  ii) Handverletzung:}
% 
% {
%  Am Beginn wird die Basis eingeschlagen, um eine stabile Kante zu
% bilden. Die Hand  wird in das auf einem festen Gegenstand ruhende
% Dreiecktuch gelegt. Die zuerst  vom Körper wegschauende Spitze wird
% nun zum Patienten hingewendet und sollte  über das Handgelenk hinaus
% am Unterarm liegen bleiben. Danach wird über der  Spitze mit den
% beiden Enden ein einfacher Knoten gemacht, in dem die Spitze 
% eingewickelt wird. Als Abschluss ist ein Weberknoten zu ziehen.}
% 
% {
%  iii) Knie- und Ellenbogenverletzung:}
% 
% {
%  Am Beginn wird die Basis eingeschlagen um eine stabile Kante zu bilden.
% Das Tuch  wird am oberen Unterschenkel angelegt, die Spitze liegt
% über dem Knie und schaut  zum Körper des Patienten. In der
% Kniekehle werden die enden überkreuzt und auf der  Vorderseite des
% Oberschenkels zu einem einfachen Knoten verbunden (Achtergang).  Danach
% wird die Spitze in den einfachen Knoten eingewickelt und mittels 
% Weberknoten fixiert.}
% 
% 
% 
% 
% \subsubsection{11.1.3) Knochenbruch:}
% 
% {
% Der Knochenbruch (Fraktur) kann in zwei Formen auftreten:}
% 
% {
%  i) offene Fraktur:}
% 
% {
%  Über der Bruchstelle kam es zu einer Verletzung der Haut. Der Knochen
% muss aber  noch nicht sichtbar sein.}
% 
% {
%  ii) geschlossene Fraktur:}
% 
% {
%  Über der Bruchstelle kam es zu keiner Verletzung der Haut.}
% 
% \subparagraph{
% Kennzeichen:}
% 
% {
%  i) sichere Kennzeichen:}
% 
% {
%   a) abnorme Fehlstellungen}
% 
% {
%  ii) unsichere Kennzeichen:}
% 
% {
%   a) Blutergüsse (Hämatome)}
% 
% {
%   b) Schwellungen}
% 
% {
%   c) Schmerzen}
% 
% {
%   d) Bewegungseinschränkungen}
% 
% \MassnahmenNG{Erste-Hilfe-Maßnahmen}{}{
% \paragraph{Erste-Hilfe-Maßnahmen}
% 
% }
% 
% {
%  i) ruhiges Positionieren der gebrochenen Körperstelle durch
% unterstützende Lagerung}
% 
% {
%  ii) Öffnen von beengenden Kleidungsstücken an der Bruchstelle}
% 
% {
%  iii) Öffnen von Schuhen bei gebrochenen Beinen, aber kein Ausziehen
% der Schuhe}
% 
% {
%  iv) Versorgung von offenen Wunden (siehe oben)}
% 
% {
%  v) Armtragetuch bei Brüchen der oberen Extremität
% einschließlich des Schlüsselbeins:}
% 
% {
%  Am Beginn wird ein Knoten in die Spitze des Dreiecktuchs gelegt, der es
% ermöglicht  aus dem Dreiecktuch eine auf eine Tasche zu bilden, in
% der der verletzte Arm  hineingelegt wird. Dabei wird die Spitze beim
% Ellenbogen positioniert, das äußere  Ende geht zur vom Knoten
% weiter entfernten Schulter, das innere Ende wird zur näher  gelegenen
% Schulter gezogen. Im Genick werden beiden Enden mittels Weberknoten 
% verbunden. }
% 
% {
%  v) uU. Schockbekämpfung}
% 
% {
%  vi) Notrufabgabe}
% 
% {
%  vii) Wärmeerhalt sichern}
% 
% {
%  iix) psychische Betreuung}
% 
% {
%  ix) laufende Kontrolle}
% 
% 
% 
% 
% \subsubsection{11.1.4) Verrenkung/Verstauchung :}
% 
% {
% {\textbf{\textit{Verstauchungen }}}{sind Verletzungen, bei
% denen ein Gelenkskopf aus seiner Pfanne herausgedreht wird und
% anschließend wieder seine ursprüngliche Position einnimmt.}}
% 
% {
% Kennzeichen:}
% 
% {
%  i) Schmerzen  ii) Schwellung  iii) Blutergüsse (Hämatome)}
% 
% {
%  iv) uU. Bewegungseinschränkungen}
% 
% 
% 
% 
% \MassnahmenNG{Erste-Hilfe-Maßnahmen}{}{
% \paragraph{Erste-Hilfe-Maßnahmen}~
% 
% }
% 
% {
%  i) Ruhigstellung ii) Kühlung  iii) Arztaufsuchung}
% 
% {
%  iv) Wärmerhalt sichern   v) psychische Betreuung}
% 
% {
%  vi) laufende Kontrolle}
% 
% 
% {
% {\textbf{\textit{Verrenkungen}}}{ sind Verletzungen, bei
% denen ein Gelenkskopf aus seiner Pfanne herausgedreht wird und
% anschließend in dieser abnormen Position verbleibt.}}
% 
% {
% Kennzeichen:}
% 
% {
%  i) Schmerzen  ii) Schwellung  iii) Blutergüsse (Hämatome)}
% 
% {
%  iv) Bewegungseinschränkungen  iv) federnde Fixation}
% 
% \MassnahmenNG{Erste-Hilfe-Maßnahmen}{}{
% \paragraph{Erste-Hilfe-Maßnahmen}~
% 
% 
% 
% {
%  i) Ruhigstellung ii) Kühlung  iii) Arztaufsuchung}
% 
% {
%  iv) Wärmerhalt sichern   v) psychische Betreuung}
% 
% {
%  vi) laufende Kontrolle}
% }
% 
% 
% \subsubsection{11.1.4) Verbrennung/Verbrühung :}
% 
% {
% Verletzungen dieser Art entstehen durch massive Einwirkung von
% großer Hitze, die von unterschiedlichen Materialen ausgehen kann.
% Dabei erfasst die Verbrühungen Verletzungen die durch Flüssigkeit
% oder Dampf hervorgerufen wurden, während die Verbrennung Verletzungen
% durch Feuer, Strahlung, Reibung, Strom, hieße feste Gegenstände
% oder Sonstigem erfasst.}
% 
% \subparagraph{
% Kennzeichen:}
% 
% {
%  1. Grad: {}-) Schmerzen}
% 
% {
%    {}-) Schwellung}
% 
% {
%    {}-) Rötungen}
% 
% {
%  2. Grad: {}-) Schmerzen}
% 
% {
%    {}-) Blasenbildung}
% 
% {
%  3. Grad: {}-) uU. keine Schmerzen}
% 
% {
%    {}-) Schorfbildung}
% 
% {
%  Schockgefahr ab: 5\% verbrannter Körperoberfläche beim Kind}
% 
% {
%     10\% verbrannter Körperoberfläche beim Erwachsenen}
% 
% {
% Zur Ermittlung der 5\% oder 10\% nimmt man als Maßstab die Hand des
% Patienten ohne seinen Daumen. Dieser Bereich stellt 1\% seiner
% Körperoberfläche dar. Kann der Patient die Brandwunde 10mal mit der
% eigenen Hand abtasten, sind 10\% seiner Körperoberfläche
% verbrannt.}
% 
% \MassnahmenNG{Erste-Hilfe-Maßnahmen}{}{
% \paragraph{Erste-Hilfe-Maßnahmen}
% 
% }
% 
% {
%  i) mind. 10 Min. Kühlen mit mäßig kaltem Wasser (kein Eis oder
% eiskaltes Wasser)}
% 
% {
%  ii) danach keimfrei Verbinden}
% 
% {
%  iii) uU. Schockbekämpfung}
% 
% {
%  iv) Notrufabgabe}
% 
% {
%  v) Wärmeerhalt}
% 
% {
%  vi) psychische Betreuung}
% 
% {
%  vii) laufende Kontrolle}
% 
% 
% 
% 
% \subparagraph{
% Verbote:}
% 
% {
%  i) keine Puder, Salben, Mehl oder Sonstiges in die Wunde geben}
% 
% {
%  ii) die Brandblasen nicht aufstechen, sie sind eine kühlende Abwehr
% des Körpers}
% 
% {
%  iii) klebende Kleidung nicht entfernen}
% 
% {
%  iv) keinen aluminiumbeschichteten Verband verwenden, dieser reflektiert
% die  Resthitze und verschlimmert die Verletzung}
% 
% 
% 
% \subsubsection{11.1.5) Erfrierung:}
% 
% {
% Erfrierungen sind Gewebsschädigungen hervorgerufen durch unzureichende
% Durchblutung wegen massiver Einwirkung von Kälte.}
% 
% \subparagraph{
% Kennzeichen:}
% 
% {
%  i) blasse, gefühlslose Haut}
% 
% {
%  ii) blaugraue, kleinfleckige (marmorierte) Verfärbungen der Haut}
% 
% {
%  iii) Bewegungseinschränkung bis hin zur Bewegungsunfähigkeit}
% 
% {
%  iv) Schmerzen}
% 
% \MassnahmenNG{Erste-Hilfe-Maßnahmen}{}{
% \paragraph{Erste-Hilfe-Maßnahmen}
% 
% 
% 
% {
%  i) Erwärmen des gesamten Körpers}
% 
% {
%  ii) warme, gezuckerte Getränke verabreichen}
% 
% {
%  iii) keimfreier, gepolsteter Verband}
% 
% {
%  iv) Öffnen von beengenden Kleidungen an der Erfrierungsstelle}
% 
% {
%  v) Notrufabgabe}
% 
% {
%  vi) psychische Betreuung, laufende Kontrolle}
% 
% \subparagraph{
% Verbote:}
% 
% {
%  i) Alkohol }
% 
% {
%  ii) Rauchen  }
% 
% {
%  iii) Reiben der Erfrierungsstelle (auch nicht in Schnee)}
% 
% {
%  iv) aktives Bewegen der erfrorenen Körperstelle (passives Bewegen
% durch den  Ersthelfer kann uU. notwendig sein, wenn es anders nicht
% möglich ist den Patienten  zu transportieren)}
% 
% }
% 
% 
% \subsubsection{11.1.6) Wunden mit Fremdkörpern:}
% 
% {
% Die Fremdkörper dürfen nicht entfernt werden, solange es sich nicht
% um kleinste Teile handelt, wie Schotter oder Kiesel. Der Fremdkörper
% muss in seiner Position stabilisiert werden. Eine Möglichkeit dabei
% ist zwei aufgerollte Mullbinden seitlich des Fremdkörpers durch
% ankleben mittels Leukoplast anzubringen. Vor dieser Stabilisierung ist
% die Wunde mit einer Wundauflage, in deren Mitte der Ersthelfer ein Loch
% ausgeschnitten hat, abzudecken. Anschließend durchzuführen ist: }
% 
% {
%  i) Notrufabgabe}
% 
% {
%  ii) Wärmeerhalt}
% 
% {
%  iii) psychische Betreuung}
% 
% {
%  iv) laufende Kontrolle}
% 
% 
% 
% \subsubsection{11.1.7) Stichwunden:}
% 
% {
% Hierbei ist zu unterscheiden ob das Messer oder sonstige Stichwerkzeug
% noch im Körper steckt oder ob es schon entfernt wurde. Bei steckendem
% Messer gilt das zur Versorgung einer Wunde mit Fremdkörper gesagte.
% Befindet sich das Messer hingegen nicht mehr im Körper, ist die Wunde
% entsprechend den allgemeinen Verbandsregeln zu versorgen. }
% 
% {
% Cave: Bei Stichwunden reicht die Wunde meist tiefer als man es von
% außen erwarten würde!}
% 
% {
% Anschließend durchzuführen ist: }
% 
% {
%  i) Notrufabgabe}
% 
% {
%  ii) Wärmeerhalt}
% 
% {
%  iii) psychische Betreuung}
% 
% {
%  iv) laufende Kontrolle}
% 
% 
% 
% \subsubsection{11.1.8) Schusswunden:}
% 
% {
% Auch die Schusswunde besteht meist aus einer größeren
% unsichtbaren, inneren Verletzungen als man es von außen
% zuverlässig sagen kann. Bei der Versorgung ist nicht nur auf die
% Eintritts-, sondern auch auf eine mögliche Austrittswunde zu achten.
% Ansonsten verläuft die Versorgung nach den allgemeinen
% Verbandsregeln.}
% 
% {
% Anschließend durchzuführen ist: }
% 
% {
%  i) Notrufabgabe}
% 
% {
%  ii) Wärmeerhalt}
% 
% {
%  iii) psychische Betreuung}
% 
% {
%  iv) laufende Kontrolle}
% 
% 
% 
% 
% 
% 
% 
% 
% 
% 
% \subsubsection{11.1.9) Augenverletzungen:}
% 
% {
% Sollten ätzende Flüssigkeiten in das Auge geraten sein, ist das
% betroffene Auge unter fließendem Wasser zu reinigen. Dabei muss das
% verletzte Auge nach unten gehalten werden, damit das unverletzte Auge
% nicht im Zuge der Reinigung mit beeinträchtigt wird. Größere
% Fremdkörper sind im Auge zu belassen und dürfen nicht versucht
% werden hinausgespült zu werden.}
% 
% {
% Anschließend muss nicht nur das verletzte Auge verbunden werden,
% sondern auch das unverletzte Auge, damit eine höhere Chance besteht,
% dass der Patient beide Augen ruhig hält und gefährliche
% Restsubstanzen nicht weiterverbreitet. Der Verband besteht aus einer
% Wundauflage, die auf das verletzte Auge gelegt wird und einer
% Dreiecktuchkrawatte, die über beide Augen sanft gebunden wird. Der
% Patient ist ab nun äußerst hilfsbedürftig. Bevor man einen
% weiteren Schritt setzt muss dies dem Patienten ausdrücklich
% mitgeteilt werden.}
% 
% {
% Abschließend durchzuführen ist: }
% 
% {
%  i) Notrufabgabe}
% 
% {
%  ii) Wärmeerhalt}
% 
% {
%  iii) psychische Betreuung}
% 
% {
%  iv) laufende Kontrolle}
% 
% 
% 
% 
% \subsubsection{11.1.10) Bisswunden:}
% 
% {
% Bisswunden tragen zwei Gefahrenquellen in sich. Die eine Gefahr entsteht
% durch den Biss selber. Hierbei gilt, umso größer und gewaltiger
% die Beißwerkzeuge, desto größer die Wunde. Als plakatives
% Beispiel dient der Biss eines Hamsters im Verhältnis zum Biss eines
% weißen Hais. \ Die andere Gefahr liegt in der Wundentzündung, in
% Folge der durch den Speichel übertragenen Keime. So ist in einem
% Beispiel der Biss eines Menschen gefährlicher als der eines
% Kampfundes, weil der Speichel des Mensch über mehr Keime verfügt,
% als der des Hundes.}
% 
% {
% Die Versorgung besteht aus einem keimfreien Wundverband entsprechend der
% allgemeinen Verbandsregeln.}
% 
% {
% Anschließend durchzuführen ist: }
% 
% {
%  i) Notrufabgabe}
% 
% {
%  ii) Wärmeerhalt}
% 
% {
%  iii) psychische Betreuung}
% 
% {
%  iv) laufende Kontrolle}
% 
% 
% 
% 
% \subsubsection{
% 11.1.11) Insektenstiche:}
% 
% {
% Bei Insektenstichen entstehen an sich nur kleinste Verletzungen, die
% grds. durch Kühlen und Abdecken mittels keimfreien Pflaster zu
% versorgen sind. }
% 
% {
% Anschließend durchzuführen ist: }
% 
% {
%  i) Notrufabgabe}
% 
% {
%  ii) Wärmeerhalt}
% 
% {
%  iii) psychische Betreuung}
% 
% {
%  iv) laufende Kontrolle}
% 
% {
% Die eigentliche Gefahr liegt in dem injizierten Gift. Da auch dies ein
% Fall der Vergiftung ist, wird für nähere Ausführungen auf
% gleichnamiges, unten stehendes Kapitel verwiesen.}
% 
% 
% 
% 
% \subsubsection{
% 11.1.12) Wunden im Genitalbereich:}
% 
% {
% Beim Mann wie bei der Frau sind die Genitalbereiche außerordentlich
% gut versorgt mit Nerven und Gefäßen. Daher können auch
% geringfügige Gewalteinwirkungen große Schäden und insbesondere
% Schmerzen hervorrufen. Nach schmerzhaften Gewalteinwirkungen, die uU.
% zu einer unnatürlichen Anschwellung des Genitalbereiches führen,
% ist eine mäßige Kühlung zu empfehlen bei raschem Aufsuchen
% ärztlicher Kontrolle. }
% 
% {
% Anschließend durchzuführen ist: }
% 
% {
%  i) Notrufabgabe}
% 
% {
%  ii) Wärmeerhalt}
% 
% {
%  iii) psychische Betreuung}
% 
% {
%  iv) laufende Kontrolle}
% 
% 
% 
% 
% \subsubsection{
% 11.1.13) sonstige körperregionsspezifische Verletzungen}
% 
% \subparagraph{
% Brustkorbverletzungen:}
% 
% {
% {Bei Verletzungen des Brustkorbs besteht immer die Gefahr, dass
% inneren lebenswichtigen Organe, wie Herz oder Lunge, beeinträchtigt
% wurden. Worin genau die Gefahren liegen, wird in den späteren
% Kapiteln behandelt (Stichwort: Pneumothorax). Für den Ersthelfer
% relevant zu wissen ist, dass die Wunde mittels Verband
% }{nicht}{ luftdicht verschlossen werden darf und der
% Patienten mit aufrechtem Oberkörper auf die verletzte Körperseite
% zu lagern ist.}}
% 
% {
% Anschließend durchzuführen ist: }
% 
% {
%  i) Notrufabgabe}
% 
% {
%  ii) Wärmeerhalt}
% 
% {
%  iii) psychische Betreuung}
% 
% {
%  iv) laufende Kontrolle}
% 
% 
% 
% 
% \subparagraph{
% Bauchverletzungen:}
% 
% {
% Verletzungen des Bauchraums können aufgrund der vielen gut
% durchblutenden Organe schnell zu einem Schock führen. Hierbei sammelt
% sich das Blut meistens in Körperhöhlen, die sich wegen des
% konstanten Blutstroms rapide ausdehnen und dadurch verhindern, dass
% auffällig viel Blut nach außen sichtbar wird. Die tatsächliche
% Gefährlichkeit einer Bauchverletzung bleibt also in den ersten
% Minuten nach Gewalteinwirkung meist unbekannt. Nach ca. einer Stunde
% spannt sich die Bauchdecke auffällig stark an und der Patient begibt
% sich von alleine in eine Schonhaltung, bestehend aus dem Anziehen der
% Beine zum Bauch und der Krümmung des Brustkorbs zum Bauch. In dieser
% Schonhaltung muss der Patient auch vom Ersthelfer unterstützt werden,
% in dem er eine Deckenrolle unter die Knie legt und den Oberkörper
% leicht erhöht lagert. Eine offene Wunde wird versorgt, indem sie
% mittels einer Wundauflage abgedeckt wird und an den Rändern durch
% Leukoplast fixiert wird.}
% 
% {
% Anschließend durchzuführen ist: }
% 
% {
%  i) Notrufabgabe}
% 
% {
%  ii) Wärmeerhalt}
% 
% {
%  iii) psychische Betreuung}
% 
% {
%  iv) laufende Kontrolle}
% 
% 
% 
% \subsection{12) Erkrankungen:}
% 
% {
% {Erkrankungen sind Beeinträchtigungen der physischen und/oder
% psychischen Integrität ohne äußere Gewalteinwirkung. Sie
% können genauso wie ihre Ursachen vielfältig auftreten. Die
% nachfolgende Abhandlung ist keine vollständige Auflistung von
% denkbaren Erkrankungen }{oder Behandlungsarten! Sie versucht nur
% einen Überblick über die meist für einen Ersthelfer relevanten
% Erkrankungen zu bieten! Solche relevanten Beeinträchtigungen können
% im Sinn der San-AV folgende }{\textit{{\quotedblbase}plötzlich
% auftretende Erkrankungen{\textquotedblleft}
% }}{sein.}\footnote{\ {Anlage 1, Modul 1
% Sanitäter-Ausbildungsverordnung-San-AV BGBl. 420/2003}}}
% 
% 
% 
% \subsubsection{
% 12.1) Arten von Erkrankungen und ihre ersthelferische Behandlung:}
% 
% \paragraph{
% 12.1.1) Herzinfarkt und Angina Pectoris:}
% 
% {
% Beiden Erkrankungen ist gleich, dass sie in den Herzkranzgefäßen
% (also jenen Gefäßen die das Herz mit Blut versorgen) zu einer
% Verengung oder Verstopfung führen, wodurch die nachfolgende
% Gewebsregion nicht mehr ausreichend mit Blut versorgt werden kann.}
% 
% 
% 
% 
% \subparagraph{
% Herzinfarkt:}
% 
% {
% Die Herzkranzgefäße werden derart stark verstopft, dass die Region
% nach der Verstopfung gar nicht mehr mit Blut versorgt wird und folglich
% abstirbt (sogen. Infarkt). Mitunter kann es sogar zu einem Platzen des
% Gefäßes vor der Verstopfungsstelle kommen. Da der Gewebstod zu
% einer Beeinträchtigung der Herztätigkeit führt, ist der Schock
% oft gesehener Begleiter des Herzversagen. Der Herzinfarkt ist daher
% lebensbedrohlich und fordert zur Behandlung lebenserhaltende
% Sofortmaßnahmen.}
% 
% {
%  i) Kennzeichen:}
% 
% {
%   a) stechender oder brennender Schmerz hinter dem Brustbein}
% 
% {
%   b) meist ein Ausstrahlen des Schmerzes in den linken Arm }
% 
% {
%   c) Atembeschwerden in Form des Gefühls eines schweren Gegenstandes
% auf   dem Brustkorb der einem daran hindert tief Luft zu holen}
% 
% {
%   c) die Schmerzen sind belastungsunabhängig (d.\,h. sie sind sogar in
% Ruhelage   zu spüren)}
% 
% {
%  ii) ersthelferische Behandlung:}
% 
% {
%   a) Lagerung mit erhöhtem Oberkörper}
% 
% {
%   b) bei besonders starken Atembeschwerden soll sich der Patient mit den
%   Armen am Boden aufstützen, um die Atemhilfsmuskulatur im
% Schultergürtel-  bereich einsetzen zu können}
% 
% {
%   c) Vermeidung von unnötigen Bewegungen des Patienten }
% 
% {
%   d) Öffnen von beengenden Kleidungsstücken}
% 
% {
%   e) uU. Öffnen eines Fensters}
% 
% {
%   f) Notrufabgabe}
% 
% {
%   g) abschließende Behandlung durch Wärmerhalt sichern, laufende
% Kontrolle,   psychische Betreuung}
% 
% 
% 
% 
% \subparagraph{
% Angina Pectoris:}
% 
% {
% Hierbei werden die Herzkranzgefäße nicht komplett verstopft,
% sondern nur teilweise verlegt. Die teilweise Verlegung führt zu einer
% unzureichenden Versorgung der nachfolgenden Gewebsregion und kann sich
% schnell zu einer Komplettverlegung verschlechtern. Daher ist auch diese
% Situation lebensbedrohlich und fordert zur Behandlung lebenserhaltende
% Sofortmaßnahmen.}
% 
% {
%  i) Kennzeichen:}
% 
% {
%   a) stechender oder brennender Schmerz hinter dem Brustbein}
% 
% {
%   b) meist ein Ausstrahlen des Schmerzes in den linken Arm }
% 
% {
%   c) Atembeschwerden in Form des Gefühls eines schweren Gegenstandes
% auf   dem Brustkorb der einem daran hindert tief Luft zu holen}
% 
% {
%   c) die Schmerzen sind belastungsabhängig (d.\,h. sie bestehen grds.
% nur bei   Bewegung)}
% 
% {
%  ii) ersthelferische Behandlung:}
% 
% {
%   a) Lagerung mit erhöhtem Oberkörper}
% 
% {
%   b) bei besonders starken Atembeschwerden soll sich der Patient mit den
%   Armen am Boden aufstützen, um die Atemhilfsmuskulatur im
% Schultergürtel-  bereich einsetzen zu können}
% 
% {
%   c) Vermeidung von unnötigen Bewegungen des Patienten}
% 
% {
%   d) Öffnen von beengenden Kleidungsstücken}
% 
% {
%   e) uU. Öffnen eines Fensters}
% 
% {
%   f) Notrufabgabe}
% 
% {
%   g) abschließende Behandlung durch Wärmerhalt sichern, laufende
% Kontrolle,   psychische Betreuung}
% 
% 
% 
% 
% \paragraph{
% 12.1.2) Schlaganfall:}
% 
% {
% Genauso wie das Herz muss auch das Gehirn äußerst stark
% durchblutet sein. Kommt es in einem Hirngefäß zum Verlegen oder
% Verstopfen, kann dies genauso zu einem Absterben der nachfolgenden
% Hirnregion führen. Die unzureichende Blutversorgung führt zu
% (meist) auffällig neurologischen Symptomen, die sich wie folgt
% kennzeichnen lassen:}
% 
% \subparagraph{
% Kennzeichen:}
% 
% {
%  i) plötzlich auftretende Lähmungserscheinungen (insbes. einer
% Körperhälfte)}
% 
% {
%  ii) uU. plötzliches auf den Boden Schlagen des Patienten (daher auch
% die  Bezeichnung {\quotedblbase}Schlag{\textquotedblleft}-Anfall)}
% 
% {
%  iii) hängende Mundwinkel, Augenlider oder Wagen- und
% \ Stirnmuskulatur}
% 
% {
%  iv) Sehstörungen}
% 
% {
%  v) Verwirrtheit bis hin zur Bewusstlosigkeit}
% 
% {
%  vi) Sprachstörungen}
% 
% {
%  vii) Kopfschmerzen}
% 
% \MassnahmenNG{Erste-Hilfe-Maßnahmen}{}{
% \paragraph{Erste-Hilfe-Maßnahmen}
% 
% 
% 
% {
%  i) Patient mit Bewusstlosigkeit: }
% 
% {
%   a) stabile Seitenlagerung}
% 
% {
%  ii) Patient ohne Bewusstlosigkeit:}
% 
% {
%   a) Lagerung mit erhöhtem Oberkörper}
% 
% {
%   b) Vermeidung von unnötigen Bewegungen des Patienten}
% 
% {
%   c) Öffnen von beengenden Kleidungsstücken}
% 
% {
%   d) uU. Öffnen eines Fensters}
% 
% {
%   e) Notrufabgabe}
% 
% {
%   f) abschließende Behandlung durch Wärmerhalt sichern, laufende
% Kontrolle,   psychische Betreuung}
% 
% }
% 
% 
% \paragraph{
% 12.1.3) Krampfanfall (Epilepsie):}
% 
% {
% Ein Krampfanfall ist der Ausdruck von fehlerhaften, unkontrollierten
% elektrischen Entladungen im Gehirn. Die Epilepsie kann sich dabei
% unterschiedlich äußern, so kann es von kurzfristig geistigen
% Abwesenheitszuständen bis hin zu länger dauernden Krämpfen mit
% anschließender Bewusstlosigkeit kommen. Eindeutige Kennzeichen sind
% wie folgt:}
% 
% \subparagraph{
% Kennzeichen:}
% 
% \begin{itemize}
%     \item  i) plötzliches Hinfallen, oft verbunden mit einem Aufschrei
%     \item  ii) generalisierte Krämpfe
%     \item  iii) Schaumaustritt aus dem Mund (uU. mit Blutbeimengung)
%     \item  iv) Urin- und Harnabgang
%     \item  v) während des Anfalls keine Atmung erkennbar
%     \item  vi) nach Beendigung des Krampfes Eintritt von Verwirrtheit,
% Tiefschlafphase oder  Bewusstlosigkeit
% \end{itemize}
% 
% \MassnahmenNG{Erste-Hilfe-Maßnahmen}{}{
% \paragraph{Erste-Hilfe-Maßnahmen}
% 
% 
% 
% {
%  i) während des Krampfes:}
% 
% {
%   a) der Patient darf nicht festgehalten werden, es müssen 
% Gegenstände an   denen sich der Patient verletzen kann entfernt
% werden}
% 
% {
%  ii) nach dem Krampf:}
% 
% {
%   a) stabile Seitenlagerung (für den bewusstlosen sowie bloß
% schlafenden    Patient); Lagerung mit erhöhtem Oberkörper für den
% bloß verwirrten Patient  b) Mundausräumen (ein Tuch um zwei
% Finger wickeln und sorgsam den    Mundraum vom Schaum befreien)}
% 
% {
%   c) Notrufabgabe}
% 
% {
%   d) abschließende Behandlung mittels Wärmerhalt sichern, laufende
% Kontrolle,   psychische Betreuung}
% 
% }
% 
% 
% \paragraph{12.1.4) Hitzeohnmacht:}
% 
% 
% {
% Durch Fehlverteilung des vorhandenen Blutvolumens entsteht eine
% vorübergehende Minderdurchblutung des Gehirns.}
% 
% \subparagraph{Kennzeichen:}
% 
% {
%  i) anfänglich rotes Gesicht  }
% 
% {
%  ii) später blass durch Blutdruckabfall}
% 
% {
%  iii) Schwindel, Schwächegefühl }
% 
% {
%  iv) Übelkeit, Erbrechen}
% 
% {
%  v) kühle, schweißige Haut  }
% 
% {
%  vi) Temperatur normal}
% 
% \MassnahmenNG{Erste-Hilfe-Maßnahmen}{}{
% \paragraph{Erste-Hilfe-Maßnahmen}
% 
% {
%  i) Flachlagerung mit erhöhten Beinen }
% 
% {
%  ii) beengende Kleidung öffnen}
% 
% {
%  iii) Kühlung des Körpers mit feuchten Tüchern}
% 
% {
%  iv) Notrufabgabe}
% 
% {
%  v) psychische Betreuung}
% 
% {
%  vi) bei Schockkennzeichen: Schocklagerung, laufende Kontrolle}
% 
% {
%  vii) bei Bewusstlosigkeit:  stabile Seitenlagerung, laufende Kontrolle}
% 
% }
% 
% 
% \paragraph{
% 12.1.5) Hitzeerschöpfung:}
% 
% {
% Infolge hitzebedingtem Schweissverlust kommt es zum Versagen der
% Kreislaufregulation.}
% 
% \subparagraph{
% Kennzeichen:}
% 
% {
%  i) anfänglich rote Haut, dann blass, kaltschweißig \& feucht}
% 
% {
%  ii) Kopfschmerzen, Schwindel, Übelkeit}
% 
% {
%  iii) Seh- \& Hörstörungen  }
% 
% {
%  iv) Durstgefühl}
% 
% {
%  v) schnelle Atmung    }
% 
% {
%  vi) ev. Krämpfe, Bewusstseinsstörung}
% 
% \MassnahmenNG{Erste-Hilfe-Maßnahmen}{}{
% \paragraph{Erste-Hilfe-Maßnahmen}
% 
% }
% 
% {
%  i) Flachlagerung mit erhöhten Beinen }
% 
% {
%  ii) beengende Kleidung öffnen}
% 
% {
%  iii) Kühlung des Körpers mit feuchten Tüchern}
% 
% {
%  iv) Notrufabgabe}
% 
% {
%  v) psychische Betreuung}
% 
% {
%  vi) bei Schockkennzeichen: Schocklagerung, laufende Kontrolle}
% 
% {
%  vii) bei Bewusstlosigkeit:  stabile Seitenlagerung, laufende Kontrolle}
% 
% 
% 
% 
% \paragraph{
% 12.1.6) Hitzeschlag:}
% 
% {
% Hitzeschlag ist das Versagen der körpereigenen Temperaturregulation
% als Folge der versagenden Kreislaufregulation einer Hitzerschöpfung.}
% 
% \subparagraph{
% Kennzeichen:}
% 
% {
% { }{\textbf{\textit{bis 40{\textdegree} C:}}}{ {}-)
% Haut warm, gerötet, Schweiß möglich}}
% 
% {
%    {}-) Kopfschmerzen, Übelkeit, Schwindel}
% 
% {
%    {}-) schnelle Atmung}
% 
% 
% 
% 
% {
% { }{\textbf{\textit{bis 41{\textdegree} C:}}}{ {}-)
% kein Schweiß}}
% 
% {
%    {}-) Haut grau, heiß und trocken}
% 
% {
%    {}-) Bewusstseinstrübung}
% 
% 
% 
% 
% {
% { }{\textbf{\textit{ab 41{\textdegree} C:}}}{ {}-)
% Haut grau und trocken}}
% 
% {
%    {}-) Krampfanfälle}
% 
% {
%    {}-) Atemstillstand}
% 
% {
%    {}-) Herzrhythmusstörung}
% 
% {
%    {}-) Bewusstlosigkeit}
% 
% \MassnahmenNG{Erste-Hilfe-Maßnahmen}{}{
% \paragraph{Erste-Hilfe-Maßnahmen}
% 
% }
% 
% {
%  i) Flachlagerung mit erhöhten Beinen }
% 
% {
%  ii) beengende Kleidung öffnen}
% 
% {
%  iii) Kühlung des Körpers mit feuchten Tüchern (insbes. an Kopf-
% und Arterienposition)}
% 
% {
%  iv) Notrufabgabe}
% 
% {
%  v) psychische Betreuung}
% 
% {
%  vi) bei Schockkennzeichen: Schocklagerung, laufende Kontrolle}
% 
% {
%  vii) bei Bewusstlosigkeit:  stabile Seitenlagerung, laufende Kontrolle}
% 
% 
% 
% \paragraph{
% 12.1.7) Sonnenstich:}
% 
% 
% {
% Durch direkte intensive Sonnenbestrahlung des Kopfes wird eine Reizung
% der Gehirnhaut hervorgerufen.}
% 
% \subparagraph{
% Kennzeichen:}
% 
% {
%  i) heißer Kopf   }
% 
% {
%  ii) Kopfschmerzen, Schwindel }
% 
% {
%  iii) Übelkeit, Erbrechen}
% 
% {
%  iv) Körpertemperatur anfänglich normal (erst später erhöht)  }
% 
% {
%  v) Nackensteifigkeit  }
% 
% {
%  vi) Bewusstseinstrübung     }
% 
% {
%  vii) beschleunigte Atmung}
% 
% {
%  iix) eventuell Krampfanfälle und Bewusstlosigkeit}
% 
% \MassnahmenNG{Erste-Hilfe-Maßnahmen}{}{
% \paragraph{Erste-Hilfe-Maßnahmen}
% 
% \begin{enumerate}
% 
%     \item  Lagerung mit erhöhtem Oberkörper
%     \item  beengende Kleidung öffnen
%     \item  Kühlung des Kopfes mit kalten Tüchern
%     \item  Notrufabgabe
%     \item  psychische Betreuung
%     \item  bei Schockkennzeichen: Schocklagerung, laufende Kontrolle
%     \item  bei Bewusstlosigkeit:  stabile Seitenlagerung, laufende Kontrolle
% 
% \end{enumerate}
% 
% }
% 
% 
% \paragraph{
% 12.1.8) Fieber:}
% 
% 
% {
% Von Fieber spricht man erst ab einer Erhöhung der
% Körperkerntemperatur auf über 38{\textdegree}C. Eine
% Körperkerntemperatur von 37{\textdegree}C bis 38{\textdegree}C stellt
% nur eine erhöhte Temperatur dar.}
% 
% \subparagraph{Kennzeichen:}
% 
% \begin{table}[H]
% \Captionof{table}{Fieber.}
% 
% \begin{tabulary}{\linewidth}{LL}
% erhöhte Temperatur: (37{\textdegree}-38{\textcelsius})
% &
%   a) Kopfschmerzen
% 
%   b) uU. leichtes Unwohlsein
% \\\addlinespace
%  ii) leichtes bis hohes Fieber: (38{\textdegree}-40{\textcelsius})
% &
%   a) Kopfschmerzen  
% 
%   b) trockene, gerötete, heiße Haut
% 
%   c) Übelkeit, Erbrechen 
% 
%   d) glasige Augen
% 
%   e) Schwindel   
% 
%   f) starkes Unwohlsein
% 
%   g) Schüttelfrost 
% 
%   h) beschleunigte Atmung
% \\\addlinespace
%  iii) sehr hohes Fieber: ({\textgreater} 40{\textcelsius})
% &
%   a) verstärkte Kennzeichen des leichten bis hohen Fiebers
% 
%   b) Fieberdelir (= Störungen von Wahrnehmung und Bewusstsein)
% \\\addlinespace
%  Fiebertod: (42,6{\textcelsius})
% &
% Eiweiß beginnt zu gerinnen; Stoffwechsel bricht zusammen; Tod 
% \\\addlinespace\bottomrule
% \end{tabulary}
% \end{table}
% 
% \MassnahmenNG{Erste-Hilfe-Maßnahmen}{}{\paragraph{Erste-Hilfe-Maßnahmen}
% 
% {
%  i) laufende Fiebermessung und allgemeine laufende Kontrolle}
% 
% {
%  ii) reichlich Flüssigkeitszufuhr}
% 
% {
%  iii) Ruhe (insbes. Vermeidung von optischen und akustischen Reizen),
% dazwischen  psychische Betreuung}
% 
% {
%  iv) uU. Notrufabgabe}
% }
% 
% 
% 
% \paragraph{
% 12.1.9) Fieberkrampf:}
% 
% 
% {
% Umfasst Krampfanfälle infolge raschen Fieberanstiegs bei Säuglingen
% und Kleinkindern bis zum 4. Lebensjahr. }
% 
% \paragraph{Kennzeichen:}
% 
% {
%  i) Schüttelkrämpfe am gesamten Körper (Dauer Sek.)}
% 
% {
%  ii) Biss auf eigene Zunge}
% 
% {
%  iii) ev. blaue Verfärbung der Lippen (Zyanose) }
% 
% {
%  iv) während des Anfalls keine Atmung erkennbar }
% 
% {
%  v) nach Anfall Verwirrtheit, Tiefschlafphase oder Bewusstlosigkeit}
% 
% \MassnahmenNG{Erste-Hilfe-Maßnahmen}{}{
% \paragraph{Erste-Hilfe-Maßnahmen}
% 
% 
% 
% {
%  i) während des Krampfes:}
% 
% {
%   a) der Patient darf nicht festgehalten werden, es müssen 
% Gegenstände an   denen sich der Patient verletzen kann entfernt
% werden}
% 
% {
%  ii) nach dem Krampf  }
% 
% {
%   a) stabile Seitenlagerung (für den bewusstlosen sowie bloß
% schlafenden    Patient); Lagerung mit erhöhtem Oberkörper für den
% bloß verwirrten Patient}
% 
% {
%   b) Mundausräumen (ein Tuch um zwei Finger wickeln und sorgsam den   
% Mundraum vom Schaum befreien)}
% 
% {
%   c) Notrufabgabe}
% 
% {
%   d) abschließende Behandlung durch Wärmerhaltsicherung, laufender
%    Kontrolle, psychische Betreuung}
% }
% 
% \paragraph{12.1.10) Unterkühlung:}
% 
% Von Unterkühlung spricht man ab dem Absinken der
% Körperkerntemperatur unter 36{\textcelsius}.
% 
% \paragraph{Kennzeichen:}
% 
% {
%  i) Abwehrstadium: (34{\textdegree}-36{\textcelsius})}
% 
% {
%   a) Atmung schnell \& tief}
% 
% {
%   b) Haut blass -- zyanotisch, kühl}
% 
% {
%   c) Muskelzittern}
% 
% {
% ii) Erschöpfungsstadium: (31{\textdegree}-33{\textcelsius} )}
% 
% {
%   a) Atmung deutlich flacher}
% 
% {
%   b) Bewusstseinstrübung  }
% 
% {
% iii) Lähmungsstadium: (27{\textdegree}-30{\textcelsius})}
% 
% {
%   a) Minimalatmung}
% 
% {
%   b) tiefe Bewusstlosigkeit  }
% 
% {
%   c) Muskellähmung}
% 
% {
% iv) Scheintod: (unter 27{\textcelsius})}
% 
% {
%   a) Atemstillstand}
% 
% {
%   b) tiefe Bewusstlosigkeit}
% 
% {
%   c) Muskellähmung   }
% 
% {
%   d) Pupillenstarre}
% 
% \MassnahmenNG{Erste-Hilfe-Maßnahmen}{}{
% \paragraph{Erste-Hilfe-Maßnahmen}~
% 
% 
% {
%  i) Patient langsam erwärmen}
% 
% {
%  ii) nur schonende aktive/passive Bewegungen des Patienten}
% 
% {
%  iii) keine alkoholischen Getränke oder Rauchen (stattdessen warme
% gezuckerte  Getränke)}
% 
% {
%  iv) Notrufabgabe}
% 
% {
%  v) psychische Betreuung}
% 
% {
%  vi) laufende Kontrolle}
% 
% }
% 
% 
% \subsection{13) Vergiftungen:}
% 
% {
% Vergiftung stellt die Beeinträchtigung des Körpers durch
% schädliche Substanzen dar. An sich kann jede Substanz ab einer
% gewissen Dosis giftig sein (vgl. Paracelsus: Die Dosis macht das
% Gift!). Die Vergiftung kann in ihrer finalen Folge zum Schock führen,
% der - wie oben näher beschrieben - unbehandelt zum Tod führt. Daher
% umfasst uU. auch die Behandlung einer Vergiftung lebenserhaltende
% Sofortmaßnahmen.}
% 
% {
% Die folgende Ausarbeitung versucht nur die relevantesten
% Vergiftungsfälle zu behandeln und stellt keinen vollständigen
% Katalog dar.}
% 
% 
% 
% 
% \subsubsection{13.1) Arten von Vergiftungen und ihre ersthelferische Behandlung:}
% 
% {
% Bei der Bekämpfung einer Vergiftung gilt der Grundsatz:}
% 
% {
% {Das Gift ist vom }{vom}{ Körper des Patienten zu
% entfernen, nicht aber }{aus}{ dem Körper des Patienten.
% }}
% 
% 
% 
% 
% \subparagraph{
% Aufnahme durch den Mund:}
% 
% {
% Am häufigsten werden Gifte über den Mund aufgenommen in Form von
% Trinken oder Schlucken. Die näheren Gründe einer solchen Vergiftung
% liegen zumeist in einem Suizidversuch oder der Unwissenheit über die
% Gefährlichkeit der zugenommenen Substanz. Als Beispiele zu nennen
% sind:}
% 
% {
%  i) Medikamente}
% 
% {
%  ii) Drogen}
% 
% {
%  iii) giftige Pilze}
% 
% {
% Sollte der Patient Gift über den Mund aufgenommen haben, darf er nicht
% zum Erbrechen gebracht werden. Grund für dieses Verbot liegt in der
% Tatsache, dass das Gift momentan im Magen eingeschlossen ist. Der Magen
% selbst ist mit einer äußerst resistenten Magenschleimhaut
% beschichtet, die selbst der salzsäurehältigen Magensäure
% entgegenwirkt. Wird der Patient nun zum Erbrechen gebracht, wandert das
% Gift samt der Magensäure die Speiseröhre hinauf und führt zu
% einer erneuten Beschädigung des Mund-Rachen-Speiseröhrenraums. Dem
% Patienten ist stattdessen Wasser schluckweise zum Trinken zu geben. }
% 
% {
% Cave: Sollte der Patient Schaumbildner getrunken haben, darf auch kein
% Wasser verabreicht werden, da es sonst zu einer Aufschäumung im Magen
% kommt, die wiederum ein Erbrechen hervorruft.}
% 
% 
% 
% 
% \subparagraph{
% Aufnahme über die Atemwege:}
% 
% {
% Ebenso häufig finden sich Vergiftungen, hervorgerufen durch das
% Einatmen von giftigen Dämpfen oder Gasen. }
% 
% {
% Eine in Wien jährlich mehrmals auftretende Gefahr liegt in dem Gas
% Kohlenmonoxid (CO). Dieses geruch- und farblose Gas entsteht bei
% Verbrennung von organischen Substanzen unter unzureichender
% Sauerstoffzufuhr. In den Wiener Altbauwohnungen ist dies bei defekten
% Gasdurchlauferhitzern der Fall. }
% 
% {
% Kennzeichen der CO-Vergiftung:}
% 
% {
%  i) Kopfschmerzen}
% 
% {
%  ii) Übelkeit}
% 
% {
%  iii) Erbrechen}
% 
% {
%  iv) Verwirrtheit}
% 
% {
% { v) aber }{keine}{ rosa-rote Verfärbung der
% Haut}\footnote{Oft wird als allgemeines Kennzeichen der CO-Vergiftung
% eine rosa-rote Verfärbung der Haut angeführt. Dies ist so nicht
% richtig! Denn der lebende Mensch erfährt äußerst selten eine
% Verfärbung seiner Haut. Nur bei Leichen ist zu erkennen, dass sie zu
% einem späteren Zeitpunkt grau verfärben und dadurch länger
% Hautfarben rosa wirken, weil das im Blut befindliche CO eine spätere
% Zersetzung der roten Farbträger bewirkt.
% \cite{ForMEd2}}}
% % (vgl.
% % Hochmeister/Grassberger/Stimpf, forensische Medizin, S. 148)
% 
% \MassnahmenNG{Erste-Hilfe-Maßnahmen}{}{
% \paragraph{Erste-Hilfe-Maßnahmen}
% 
% }
% 
% {
%  i) bei Gefahr einer Selbstvergiftung geht Selbstschutz vor Fremdschutz}
% 
% {
%  ii) Rettung aus dem Gefahrenbereich}
% 
% {
%  iii) situationsabhängige Behandlung:}
% 
% {
%   a) Patient ist ansprechbar: }
% 
% {
%   Lagerung mit erhöhtem Oberkörper}
% 
% {
%   b) Patient ist nur bewusstlos: }
% 
% {
%   stabile Seitenlagerung}
% 
% {
%   c) Patient hat keine Atmung und Kreislauf:   }
% 
% {
%   Wiederbelebung}
% 
% {
%  iv) anschließende Behandlung durch Wärmeerhaltsicherung, laufende
% Kontrolle und  psychische Betreuung}
% 
% 
% 
% 
% \subparagraph{
% Aufnahme über die Haut:}
% 
% {
% Giftige Substanzen, die über die Haut aufgenommen werden (wie manche
% Pestizide), sind durch Abspülen der betroffenen Körperstelle zu
% beseitigen. Dabei darf der Ersthelfer nicht in direkten Kontakt mit dem
% Gift kommen. Die weiterführenden Behandlungen sind
% situationsabhängig und entsprechen dem oben gesagten.}
% 
% 
% 
% 
% \subparagraph{
% Insektenstich:}
% 
% {
% Für den Durchschnittsmenschen ist der einzelne Stich eines heimischen
% Insekts nicht gefährlich. Für Allergiker genügt uU. schon ein
% einzelner Stich, um Schocksymptomatiken hervorzurufen. Ab welcher
% Anzahl Insektenstiche gefährlich werden, kann man generalisiert nicht
% sagen, da jeder Mensch unterschiedlich auf das Gift reagiert. Um zu
% erkennen welches Insekt zu gestochen hat, gibt es eine einfache
% Unterscheidung zwischen Biene und Wespe. Blieb nach der Attacke der
% Stachel samt Eingeweidesack an der Einstichstelle stecken, handelte es
% sich um eine Biene. Im gegensätzlichen Fall war es eine Wespe.}
% 
% {
% Der am Körper des Patienten hängenden Eingeweidesack der Biene darf
% nicht bei dessen Entfernung zusammengedrückt werden, da man ansonsten
% das gesamte Restvolumen an Gift in den Körper injiziert. Der Stachel
% ist mittels einer Pinzette anzufassen und vorsichtig herauszuziehen.}
% 
% {
% Das Gift darf in keinem Fall vom Ersthelfer aus der Einstichstelle
% herausgesaugt werden. Grund hierfür ist die Tatsache, dass der
% Ersthelfer beim Ansaugen das Gift über die Schleimhaut im eigenen
% Mundraum aufnehmen würde und sich somit selbst gefährdet. Die
% Einstichstelle ist zu kühlen und keimfrei zu verbinden. Der
% Wärmeerhalt ist (falls notwendig) zu sichern, der Patient ist
% psychisch zu betreuen und laufend zu kontrollieren. Spätestens bei
% Eintritt von Schocksymptomatiken ist ein Notruf abzugeben.}
% 
% 
% 
% 
% \paragraph{Schlangenbiss:}
% 
% {
% Das durch einen Schlangebiss in den Körper injizierte Gift darf nicht
% aus dem Körper herausgesaugt werden. Die Einstichstelle ist zu
% kühlen und keimfrei zu verbinden. Der Wärmeerhalt ist (falls
% notwendig) zu sichern, der Patient ist psychisch zu betreuen und
% laufend zu kontrollieren. In jedem Fall ist ein Notruf abzugeben.}
% 
% 
% 
% 
% \subsection{14) abschließende Behandlung:}
% 
% {
% Zum Abschluss des Themenblock Weitere Erste Hilfe, aber eigentlich auch
% lebensrettende Sofortmaßnahmen, ist nochmals zu wiederholen, dass
% es in der Ersten Hilfe drei {\quotedblbase}ersthelferische
% Allzweckwaffen{\textquotedblleft} gibt, die in so gut wie jedem Fall
% verpflichtend anzuwenden sind. Diese sind:}
% 
% \begin{enumerate}
% 
% 
%     \item  i) Wärmeerhalt sichern:
% 
% {\scriptsize
%  Bei 1{\textdegree}C gesunkener Körperkerntemperatur sinkt die
% Überlebenschance des Patienten um  ganze 30\%. Eine Erhaltung der
% Körperwärme ist daher absolut notwendig. Der Ersthelfer kann  sich
% hierfür unterschiedlicher Mittel bedienen, eines ist die in einer
% Autoapotheke gemäß  ÖNORM befindliche Rettungsdecke.}
% 
%     \item ii) psychische Betreuung:
% 
% {\scriptsize
%  Die  Wichtigkeit psychologischer Zuwendungen wurde ohnedies von
% unzähligen  medizinischen und psychologischen Studien bewiesen.
% Aufbauen kann sich die  psychologische Betreuung aus  einem guten
% Zuspruch und einer Hand die auf die Schulter  des Patienten gelegt
% wird.}
% 
%     \item  iii) laufende Kontrolle:
% 
% {\scriptsize
%  Der Ersthelfer soll sich alle 30 Sek. durch Überprüfung der
% Atemtätigkeit und allgemeinen  Konstitution des Patienten davon
% überzeugen, dass der Patientenzustand unverändert  geblieben ist.}
% 
% \end{enumerate}
% 
% 
% {
% Sollte in einem Themenbereich nicht ausdrücklich auf diese
% Maßnahmen hingewiesen worden sein, bedeutet es nicht, dass sie
% nicht anzuwenden sind. Sie sind nur in Ausnahmefällen ungeeignet!
% Bsp.: Der Ersthelfer möchte einen an Hitzeschock leidenden Patienten
% übermäßig wärmen.}




% \section{IV. Katastrophenmedizin in der Ersten Hilfe}
% 
% {
% Die zuvor behandelten Themenbereiche sind in ihren Regeln immer von
% einem Prinzip ausgegangen, dem der Individualmedizin. Hierbei gibt es
% genau so viele oder sogar mehr Ersthelfer als Patienten. In der
% Realität kommen hingegen Fälle einer Überzahl von Patienten im
% Verhältnis zu Ersthelfern sehr häufig vor, wodurch wir uns solcher
% Problematiken nicht verschließen sollten. Die oben erwähnten
% Grundsätze gelten auch für den Katastrophenbereich, mit der
% Besonderheit, dass die Interessenabwägung einen größeren
% Stellenwert erhält. Sollten also mehrere Patienten gleich stark
% beeinträchtigt sein, muss der Ersthelfer alle gleich gut bzw. gleich
% schlecht (in Form von weniger umfangreich) behandeln. Er kann nicht
% jedem Patienten seine volle Aufmerksamkeit und Fürsorge zu kommen
% lassen, sondern muss eine ausgewogene Gesamtbehandlung finden. Steht er
% sogar vor der Entscheidung eine Maßnahme nur an einen Patienten
% durchführen zu können, obwohl sie mehrere Personen benötigen,
% muss der Ersthelfer jene Person behandeln, die seiner Meinung nach die
% höhere Überlebenschance hat.}
% 
% {
% Das diese Entscheidungen und Maßnahmen äußerst schwierig zu
% treffen bzw. umzusetzen sind, stellt hier keiner in Frage. Der
% Ersthelfer muss sich aber bewusst sein, dass es physikalische Grenzen
% gibt an die jeder Mensch gebunden und die weder er noch andere Personen
% an seiner Stelle brechen hätten können.}



% \ChapterBibliography
% 
% 
% {
% {\textsc{Böhmer, }}{Roman }{\textsc{/ Schneider,
% }}{Thomas }{\textsc{/ Wolcke }}{Benno: Reanimation
% kompakt}{\textsuperscript{2}}{. Nach den internationalen
% ERC-Richtlinien 2005 basierend auf den CoSTR 2005 (ILCOR/AHA) (Mainz:
% Naseweis Verlag, 2006)}}
% 
% 
% 
% 
% {
% {\textsc{Hackstock, }}{Horst }{\textsc{/
% Kühberger, }}{Rudolf: Arbeiter-Samariter-Bund Österreichs
% Bundes-schulungsreferat: Leitfaden für Vortragende. Lebensrettende
% Sofortmaßnahmen am Ort des Verkehrsunfalls (Wien:
% Arbeiter-Samariter-Bund Österreichs intern, Mai 2006)}}
% 
% 
% 
% 
% {
% {\textsc{Hackstock, }}{Horst }{\textsc{/
% Kühberger, }}{Rudolf: Arbeiter-Samariter-Bund Österreichs
% Bundes-schulungsreferat: Leitfaden für Vortragende.
% Breitenschulungskurs (Wien: Arbeiter-Samariter-Bund Österreichs
% intern, Sept. 2008)}}
% 
% 
% 
% 
% {
% {\textsc{Hochmeister}}{, Manfred /
% }{\textsc{Grassberger}}{, Martin /
% }{\textsc{Stimpfl}}{, Thomas: forensische Medizin. für
% Studium und Praxis (Wien: Maudrich Verlag, 2007)}}
% 
% 
% 
% 
% {
% {\textsc{Holaubeck}}{,
% Karl/}{\textsc{Nistler}}{,
% Kurt/}{\textsc{Scherer}}{, Nicolas: Erste Hilfe. 16
% Stunden für das Leben (Wien: Arbeiter-Samariter-Bund Österreichs,
% 2005)}}
% 
% 
% 
% 
% {
% {\textsc{Karutz,}}{ Harald / }{\textsc{von
% Buttlar}}{, Manfred: Kursbuch Erste
% Hilfe}{\textsuperscript{2}}{ (München: Deutscher
% Taschenbuch Verlag, 2008)}}
% 
% 
% 
% 
% {
% {\textsc{Keggenhoff, }}{Franz: Erste Hilfe. Das offizielle
% Handbuch. Sofortmaßnahmen bei Babys, Kindern und Erwachsenen
% (München: Südwest Verlag, 2007)}}
% 
% 
% 
% 
% {
% {\textsc{Köhnlein, }}{Edzard }{\textsc{/ Weller
% }}{Siegfried [Hrsg.]: Erste
% Hilfe}{\textsuperscript{10}}{ (Stuttgart: Georg Thieme
% Verlag, 2004)}}
% 
% 
% 
% 
% {
% {\textsc{Köhnlein, }}{Edzard }{\textsc{/ Weller
% }}{Siegfried [Hrsg.]: Erste Hilfe. Aktuelle Empfehlungen zur
% kardiopulmonalen Reanimation (Stuttgart: Georg Thieme Verlag, 2006)}}
% 
% 
% 
% 
% {
% {\textsc{Pschyrembel}}{, Willibald: Klinisches
% Wörterbuch (Berlin: Walter de Gruyter Verlag, 2004)}}
% 
% 
% 
% 
% {
% {\textsc{Wirth}}{, Armin: Erste Hilfe
% unterwegs}{\textsuperscript{3}}{. Effektiv und praxisnah
% (Bielefeld: Reise-Know-How Verlag Peter Rump, 2007)}}

